\documentclass[12pt, a4paper, openany]{book}

\usepackage{fontspec}
\setmainfont{TeX Gyre Pagella}
\usepackage{geometry}
\geometry{margin=1.25in, top=1.2in, bottom=1.3in}
\usepackage{setspace}
\onehalfspacing
\usepackage{amsmath}
\usepackage{amssymb}
\usepackage{amsthm}
\usepackage{mathtools}
\usepackage{stmaryrd}
\usepackage{titlesec}
\usepackage{titletoc}
\usepackage{fancyhdr}
\setlength{\headheight}{14pt}
\usepackage{booktabs}
\usepackage{array}
\usepackage{longtable}
\usepackage{listings}
\usepackage{xcolor}
\usepackage{csquotes}
\usepackage{enumitem}
\usepackage{hyperref}
\usepackage{parskip}
\usepackage{microtype}

\hypersetup{colorlinks=true, linkcolor=black, citecolor=black, urlcolor=black}

\setlength{\parindent}{1.5em}
\setlength{\parskip}{0pt}

% Theorem environments
\theoremstyle{plain}
\newtheorem{theorem}{Theorem}[chapter]
\newtheorem{proposition}[theorem]{Proposition}
\newtheorem{lemma}[theorem]{Lemma}
\newtheorem{corollary}[theorem]{Corollary}
\theoremstyle{definition}
\newtheorem{definition}[theorem]{Definition}
\newtheorem{example}[theorem]{Example}
\newtheorem{construction}[theorem]{Construction}
\theoremstyle{remark}
\newtheorem{remark}[theorem]{Remark}

% Chapter and section formatting
\titleformat{\chapter}[display]
  {\normalfont\large\bfseries}
  {\chaptertitlename\ \thechapter}{12pt}{\Large}
\titleformat{\section}{\normalfont\large\bfseries}{}{0em}{}
\titleformat{\subsection}{\normalfont\normalsize\bfseries}{}{0em}{}
\titleformat{\subsubsection}{\normalfont\normalsize\itshape}{}{0em}{}

% Code listings
\definecolor{codegreen}{rgb}{0,0.5,0}
\definecolor{codegray}{rgb}{0.5,0.5,0.5}
\definecolor{codeblue}{rgb}{0,0,0.6}
\lstdefinestyle{rust}{
  language=C,
  basicstyle=\small\ttfamily,
  keywordstyle=\color{codeblue}\bfseries,
  commentstyle=\color{codegray}\itshape,
  stringstyle=\color{codegreen},
  breaklines=true,
  frame=single,
  framesep=4pt,
  xleftmargin=12pt,
  xrightmargin=4pt,
}

% Operators
\newcommand{\Pop}{\mathsf{Pop}}
\newcommand{\Refuse}{\mathsf{Refuse}}
\newcommand{\Bind}{\mathsf{Bind}}
\newcommand{\Collapse}{\mathsf{Collapse}}
\newcommand{\Hist}{\mathbf{Hist}}
\newcommand{\Obs}{\mathbf{Obs}}
\newcommand{\calH}{\mathcal{H}}
\newcommand{\calC}{\mathcal{C}}
\newcommand{\calA}{\mathcal{A}}
\newcommand{\calN}{\mathcal{N}}
\newcommand{\calR}{\mathcal{R}}
\newcommand{\Adm}{\mathsf{Admissible}}
\newcommand{\Rfsd}{\mathsf{Refused}}
\newcommand{\Col}{\mathsf{Collapsed}}
\newcommand{\Proc}{\mathsf{Process}}

\pagestyle{fancy}
\fancyhf{}
\fancyhead[LE]{\small\textit{Spherepop: Toward a Complete Theory}}
\fancyhead[RO]{\small\textit{\leftmark}}
\fancyfoot[C]{\thepage}

\title{
  \vspace{2cm}
  {\Huge\bfseries Spherepop}\\[1em]
  {\large\itshape Histories, Continuations, and the Algebra\\
  of Reachable Computation}\\[3em]
  {\normalsize Toward a Complete Theory: History-Primary Computation,\\
  Type Systems, Error Correction, and Garbage Collection}
}
\author{Flyxion\\[0.5em]\small Independent Researcher}
\date{}


\begin{document}

\maketitle
\thispagestyle{empty}

\frontmatter

\chapter*{Abstract}
\addcontentsline{toc}{chapter}{Abstract}

Most programming languages begin with values, variables, and state transitions.
Spherepop begins elsewhere: with the claim that histories are ontologically prior
to states, and that computation is the progressive restriction of possibility rather
than the manipulation of values.

This monograph presents the complete theory of the Spherepop programming language
and its philosophical foundations. We proceed from two converging arguments. The
first is a theory of notes: the observation that notes are not passive records of
the past but anti-collapse artifacts that preserve access to continuations. The
second is the Spherepop inversion: the observation that state is a collapsed
history and that taking this seriously demands a new computational ontology.

These arguments converge because they are the same argument. A note is an
externalized continuation. A Spherepop history is an internalized note. Civilization
builds note infrastructure; Spherepop builds note infrastructure for computation.

The monograph develops this convergence through nine parts. Part~I establishes
the dual failure of the passive theory of notes and the state-centric theory of
computation. Part~II derives the four Spherepop operators---Pop, Refuse, Bind, and
Collapse---from first principles as the minimal complete set for managing possibility
space. Part~III develops the full operational semantics, the History category, and
collapse as quotient. Part~IV presents the complete admissibility type system with
Progress, Preservation, and Collapse Soundness theorems. Part~V derives the
seven note classes as physical realizations of the four operators, establishing
that notes are Spherepop programs and Spherepop programs are notes. Part~VI covers
compilation, the Event IR, history equivalence as compiler correctness, garbage
collection as admissibility-guided collapse, and error correction as certified
refusal. Part~VII develops the deep mathematical structure: sheaf semantics, CLIO
projections, the Santoro--Waghmare--Panaretos unification, and active geodesic
inference. Part~VIII situates the framework within civilization-scale note
infrastructure, MEM\textbar8, and Repair Theory. Part~IX collects open problems.

Four principal theorems anchor the work: the Conservation of Possibility, the
Ephemerality Inversion, the Civilizational Collapse Theorem, and the
History--Observation Adjunction. Four implementation discoveries retroactively
sharpened the theory: refusal must carry reasons; collapse must specify quotient
rules; type checking depends on the assumed boundary of the world; and compiler
correctness requires history equivalence, not mere output equivalence.


\begin{quote}\small\textit{This manuscript is not presented as a finished formalization. It is a staged research program. Some results are proved fully, some are proof sketches, and some are conjectural bridges to be developed through implementation and subsequent work. The goal is to make the program precise enough to be falsifiable, not to claim it complete.}\end{quote}

\newpage

\tableofcontents

\mainmatter

%% ============================================================
%% PART I: THE PROBLEM
%% ============================================================
\part{The Problem}

\chapter{The Passive Theory of Notes and Its Failure}

There is a theory of notes so pervasive that it rarely receives explicit statement.
On this theory, a note is a device for compensating the limitations of biological
memory. Something happens---a meeting is scheduled, a formula is derived, a
perception is formed---and because the mind cannot reliably retain it, the note
is made. The note is therefore a prosthetic memory: it externalizes information
onto a more durable substrate, where it may be retrieved when biological recall
fails.

This is the passive theory of notes. It places the note in a retrospective
relation to time. The note exists because something happened in the past that is
worth preserving. Its value is proportional to the fidelity with which it
reproduces a prior state: a good note accurately records what occurred; a poor note
distorts it; an ephemeral note is low-value because it preserves little and does
not endure.

The passive theory is not entirely wrong. Many notes are made for the reason it
describes. But as a general account of what notes are and why they matter, it
fails on its own terms and in a direction that reveals something important about
cognition, language, and civilization.

The failure becomes apparent the moment one examines not celebrated notes but
common ones. A grocery list does not record anything that happened. The items on
the list do not exist yet in the form the list anticipates: they have not been
purchased, assembled, or carried home. The list records a future act. A blueprint
describes a building not yet constructed. Source code records computations not yet
run. A theorem records consequences not yet derived. A calendar records commitments
not yet honored. A playlist records a future act of listening. A bookmark records
a future act of reading.

None of these are memories. All of them are notes.

The passive theory cannot accommodate this class of cases without becoming
incoherent. It might be extended to say that a note records an \emph{intention}
rather than an event, but this extension merely postpones the difficulty. The
deeper problem is that the passive theory places every note in a backward-looking
relation to time. Even when the recorded content is a future intention, the note
is still conceived as a trace of a prior mental event. The orientation is always
retrospective.

This gets the phenomenology of note-taking almost exactly wrong. When one makes a
to-do list, one is not primarily recording a past intention; one is constructing
a future. The note is not oriented toward what was; it is oriented toward what
has not yet occurred and what might not occur without the note's assistance.

\section{The Continuation Account}

A more adequate theory begins from a distinction that theoretical computer science
has found indispensable: the difference between a \emph{state} and a
\emph{continuation}. A state is a snapshot of a system at a moment in time---a
summary, a compression, a projection. A continuation is what comes next: the
future execution enabled by a present configuration, the program of further
development that a state makes possible.

Most theories of information focus on states. A database stores states. A
photograph is described as capturing a state. A memory is treated as the retention
of a past state. But what makes any stored state useful is not the state itself.
It is the continuation that state makes accessible.

Consider a Git commit. The commit is not valuable because it stores bytes. It is
valuable because it preserves a path through development space---a specific
trajectory through the possibility space of a codebase that would otherwise require
expensive reconstruction. The stored state is instrumentally useful only because
it is a node in a navigable graph of continuations.

\begin{definition}[Note]
Let $\calC$ denote a space of continuations and let $A$ be an agent. A
\emph{note} is any artifact $N$ such that there exists a continuation
$c \in \calC$ for which
\[
P_A(c, t \mid N) > P_A(c, t),
\]
where $P_A(c, t)$ denotes the probability that agent $A$ can successfully
reconstruct or traverse continuation $c$ at future time $t$.
\end{definition}

This definition is deliberately broad. It immediately encompasses bookmarks, maps,
photographs, source code, language, and libraries without privileging any medium.
What unifies them is the functional condition: a note increases the reachability
of at least one continuation.

\begin{definition}[Continuation Volume]
The \emph{continuation volume} of a note $N$ is
\[
V_C(N) = \bigl|\{c \in \calC : P_A(c, t \mid N) > P_A(c, t)\}\bigr|.
\]
\end{definition}

\begin{definition}[Reachability Leverage]
The \emph{reachability leverage} of a note $N$ is
\[
\Lambda(N) = \frac{\sum_{c} \bigl[C_r(c) - C_r(c \mid N)\bigr]}{C_{\text{create}}(N)},
\]
where $C_r(c)$ is the cost of reconstructing continuation $c$ from first principles
and $C_r(c \mid N)$ is the cost given access to $N$.
\end{definition}

Reachability leverage varies over many orders of magnitude. A ten-second bookmark
may save hours of future search. A one-line equation may preserve centuries of
accumulated derivation. The leverage analysis predicts which notes are most worth
making and preserving, entirely independently of their fidelity to past states.

\section{Notes as Anti-Collapse Artifacts}

The Spherepop framework, developed in subsequent chapters, holds that histories
are primary and states are their compressed residues. A state $s$ is a projection
\[
\sigma : \calH \to S
\]
that collapses many histories $h \in \calH$ into a single observable state,
producing what the framework calls the \emph{state illusion}: the appearance that
the present configuration is self-sufficient, when it is in fact the endpoint of
a history whose fibers have been discarded.

A note is, in this framework, an anti-collapse artifact: an artifact that preserves
selected fibers of $\calH$ against the projection $\sigma$.

A note does not literally preserve ``selected fibers'' of $\calH$, since a fiber
is a set of histories induced by a projection, not a portable object. What a
note preserves is evidence that reduces uncertainty \emph{within} a fiber, or
that permits reconstruction of distinctions a particular projection suppresses.
This makes anti-collapse explicitly relative to an observational rule, rather
than an absolute property of the artifact itself.

\begin{proposition}[Notes as Anti-Collapse Artifacts]
A note $N$ is anti-collapse relative to an observation $c$ if and only if $N$
reduces uncertainty within some fiber $c^{-1}(o)$:
\[
N \text{ is anti-collapse relative to } c
\quad\Longleftrightarrow\quad
N \text{ reduces uncertainty within some } c^{-1}(o).
\]
\end{proposition}

This is the precise sense in which the passive theory inverts the correct account.
The passive theory holds that a note preserves the past against forgetting. The
correct account holds that a note preserves access to futures against collapse.
The distinction is not merely philosophical: it determines what kinds of notes
matter, how notes should be evaluated, and what it means for a note to succeed
or fail.

\chapter{The State-Centric Assumption and Its Consequences}

\section{The Standard Premise}

Open any introductory programming text and you encounter a first lesson: a variable
is a named location that holds a value. A program begins in some initial state,
instructions modify that state, and the final state is the answer.

\begin{lstlisting}[style=rust, caption={The standard state-centric view}]
x = 5
y = x + 1
\end{lstlisting}

This appears natural. It mirrors how we naively think about physical processes: an
object has properties; those properties change; the current properties constitute
the object's state. But the appearance of naturalness is produced by a prior
choice---the choice to treat present configuration as the fundamental unit of
description.

\section{The State Illusion}

Consider what this choice erases. In version control, the repository is not merely
its current file tree; it is the entire sequence of commits that produced it. In
event sourcing, database records are append-only event logs; the current state is
a fold over that log. In scientific experiment, a measurement is the terminus of
a sequence of decisions, calibrations, and interpretive choices. In evolutionary
biology, a genome is a compressed record of selective pressures.

In each case, the historical structure is epistemically and causally prior to the
present configuration. State is the summary. History is the substance---though,
as the Layered World Model below makes precise, operational state is a
maintained \emph{consequence} of history, whereas observable state is a
task-relative \emph{quotient} of it, and the two should not be conflated.

\section{The Layered World Model}
\label{sec:layered-world-model}

Before stating the State Illusion precisely, it is necessary to fix a vocabulary
that the remainder of this book will use without further qualification. The
recurring confusion in early drafts of this theory---and in much writing about
event-sourced systems generally---is the collapse of several distinct notions
under the single word ``state.'' We separate four layers.

\begin{definition}[Authoritative History]
The \emph{authoritative history} at time $t$ is
\[
H_t = [e_1, \ldots, e_t] \in E^*.
\]
It contains accepted, provenance-bearing events. It is append-only and is the
sole authority from which any replayable structure may be reconstructed.
\end{definition}

\begin{definition}[Operational State]
The \emph{operational state} is
\[
M_t = F_{\mathrm{op}}(H_t; M_0).
\]
It contains the structures needed to continue execution efficiently: the current
option space, the admissibility index, the dependency graph, symbol bindings,
and related caches. It is derived from history, but it is not an
\emph{observation} of history in the sense developed in Chapter~\ref{ch:no-direct-observation}---it is the machinery execution needs, maintained as a fold.
\end{definition}

\begin{definition}[Semantic State / Read Model]
A \emph{semantic state} or \emph{read model} under scheme $q$ is
\[
K_t^q = F_{\mathrm{sem}}^q(H_t; K_0^q).
\]
It answers a particular interpretive question posed of the history. LastWrite,
Accumulate, Projection, denotational semantics, and type extraction (developed
in later chapters) all belong to this layer.
\end{definition}

\begin{definition}[Rendered Appearance]
\emph{Rendered appearance} is
\[
x_t = R(M_t, K_t^q, \omega_t),
\]
where $\omega_t$ contains disposable presentation conditions---viewport, theme,
camera, sampling seed, level of detail. Rendering must not read the unbounded
history directly; it reads only the already-derived operational and semantic
layers.
\end{definition}

We reserve the word \emph{world} for the authoritative computational
configuration developed in Section~\ref{sec:spherepop-world} below, and we
qualify every other use of ``state'' throughout this book as operational,
semantic, observable, or rendered. Collapse, defined formally in
Chapter~\ref{ch:collapse-quotient}, is not simply ``observation'' in the sense
of a pure look; it is the attested adoption of a task-relative quotient as a
constraint on future computation. This distinction---between looking and
committing to what was seen---is the single most important repair this edition
makes to the theory, and it is developed fully in
Section~\ref{sec:view-vs-collapse}.

\begin{definition}[State Degeneracy]
Let $E$ be an event alphabet, $\calH = \bigcup_{n=0}^{\infty} E^n$ the set of all
finite histories, and $S$ an observable state space. A \emph{state projection} is
a function $\sigma : \calH \to S$. The \emph{degeneracy} of a state $s \in S$ is
\[
D(s) = |\sigma^{-1}(s)|.
\]
\end{definition}

\begin{proposition}[The State Illusion]
\label{prop:state-illusion}
For any nontrivial finite state space $S$ and any non-injective projection
$\sigma : \calH \to S$, there exist distinct histories $H_1 \neq H_2$ such that
$\sigma(H_1) = \sigma(H_2)$. Consequently,
\[
\sigma(H_1) = \sigma(H_2) \not\Rightarrow H_1 = H_2.
\]
\end{proposition}

\begin{proof}
Since $|\calH|$ is countably infinite and $|S|$ is finite, $\sigma$ cannot be
injective. Therefore $D(s) > 1$ for at least one $s \in S$, which yields the
required $H_1 \neq H_2$ with $\sigma(H_1) = \sigma(H_2)$.
\end{proof}

\begin{remark}
Proposition~\ref{prop:state-illusion} is the formal statement of the state
illusion: any program that exposes only $\sigma(H)$ to its user has silently
discarded distinguishing information about the history $H$. The conventional
debugger is an instrument for reconstructing $H$ from $\sigma(H)$ after
information has already been lost.
\end{remark}

\section{Consequences for Language Design}

Programming languages that treat state as primary inherit the degeneracy of
Proposition~\ref{prop:state-illusion} as a structural feature. Race conditions
arise because $\sigma$ is non-injective on interleaved histories. Use-after-free
vulnerabilities arise because $\sigma$ conflates histories that differ only in
deallocation events. The auditing and provenance tools that practitioners add to
production systems are, in each case, partial reconstructions of $H$ from $\sigma(H)$.

The Spherepop design decision follows directly: if history is what is lost,
history must be what is preserved.

\section{The Connection to Notes}

The state illusion and the passive theory of notes are the same mistake viewed from
different angles.

The passive theory treats a note as a record of a past state. The state-centric
theory of computation treats a program's output as its definitive product. Both
identify value with a terminal state rather than with the trajectory that produced
it. Both discard the history in favor of the summary.

The Spherepop framework and the continuation theory of notes are therefore the
same correction viewed from different angles. Both insist that the trajectory is
primary. Both insist that what matters is not what a system is at a moment but
what it has done and what it can still do. Both treat state as a derived notion:
useful for some purposes, but never the fundamental unit of description.

This convergence is the organizing principle of the present monograph.

%% ============================================================
%% PART II: FIRST PRINCIPLES
%% ============================================================
\part{First Principles}

\chapter{Deriving the Operators}
\label{ch:collapse-quotient}

\section{The Minimal Requirements}

If histories are primary and states are derived, then a computational framework
must at minimum be able to:

\begin{enumerate}
\item Commit to a specific continuation from among the available possibilities.
\item Document that a possible continuation is inadmissible, with reasons.
\item Record a dependency between two elements of the history.
\item Adopt an admissible projection as a constraining judgment while recording
its rule, result, provenance, and task boundary.
\end{enumerate}

We claim that these four requirements are not only necessary but sufficient: they
constitute the minimal complete set of operations on possibility space. All other
computational phenomena are derived from these four.

\begin{definition}[Event Alphabet]
Let $E$ be a set of events, partitioned into primitive generators:
\[
E = \{\Pop\} \cup \{\Refuse\} \cup \{\Bind\} \cup \{\Collapse\} \cup E_\lambda
\]
where $E_\lambda$ contains the lambda-calculus events $\{\mathsf{LamIntro},
\mathsf{Apply}\}$ and scope markers.
\end{definition}

\begin{definition}[Spherepop World]
\label{sec:spherepop-world}
A \emph{Spherepop world} at time $t$ is
\[
W_t = (H_t, M_t), \qquad M_t = F_{\mathrm{op}}(H_t; M_0),
\]
where $H_t \in \calH$ is the authoritative history and $M_t$ is the operational
state of Section~\ref{sec:layered-world-model}, containing the option space,
admissibility index, dependency graph, and symbol bindings. Its implementation
representation is
\[
\widehat W_t = (H_t, \widehat M_t), \qquad \widehat M_t \overset{!}{=} F_{\mathrm{op}}(H_t; M_0),
\]
where $\widehat M_t$ is a cached materialization required to satisfy the replay
invariant of Chapter~\ref{ch:compiler} ($\widehat M_t$ must equal the value the
fold $F_{\mathrm{op}}$ would recompute from $H_t$).

The option space is a projected component of operational state,
\[
\Omega_t = \operatorname{options}(M_t) \subseteq \Omega_0,
\]
rather than an independently authoritative second coordinate of the world. This
is the resolution of an apparent tension in earlier drafts, which alternated
between saying that no state is stored and observing that the World evidently
stores an option set: the correct statement is that no independently
authoritative \emph{semantic} or \emph{observable} state is stored, while
operational materializations, including $\Omega_t$, may be stored, provided
they remain replay-equivalent to the fold over history.

\begin{remark}
For notational economy, the remainder of this book continues to write $W =
(H, \Omega)$ in places where only the option-space component of $M_t$ is
relevant to the argument at hand (as in the operator definitions immediately
below). This is shorthand for $(H_t, \operatorname{options}(M_t))$ and should
not be read as reinstating $\Omega$ as an independent second coordinate.
\end{remark}
\end{definition}

\section{The Four Operators}

\begin{definition}[Pop: Commitment]
The \emph{commitment operator} for symbol $x \in \Omega$ is
\[
P_x : (H, \Omega) \mapsto (H \mathbin{++} [\Pop(x)],\; \Omega \setminus \{x\}).
\]
$P_x$ is undefined when $x \notin \Omega$. Pop records a commitment and removes
the committed symbol from the option space, foreclosing all futures in which $x$
takes a different value.
\end{definition}

\begin{definition}[Refuse: Documented Inadmissibility]
The \emph{refusal operator} for symbol $x$ with reason $r$ is
\[
R_{x,r} : (H, \Omega) \mapsto (H \mathbin{++} [\Refuse(x, r)],\; \Omega).
\]
Note that $\Omega$ is unchanged: refusal records inadmissibility without
foreclosing structural availability. The reason $r$ is a first-class component
of the operation. A refusal without a reason documents nothing.
\end{definition}

\begin{definition}[Bind: Dependency Declaration]
The \emph{bind operator} is
\[
B_{a,b} : (H, \Omega) \mapsto (H \mathbin{++} [\Bind(a, b)],\; \Omega).
\]
Bind records that two trajectory elements are coupled without consuming either
from the option space.
\end{definition}

\begin{definition}[Collapse: Attested Adoption of a Quotient]
The \emph{collapse operator} with collapse rule $c$, observed value $o = c(H)$,
provenance $\rho$, and task boundary $\kappa$ is
\[
C_{x,c,o,\rho,\kappa} : (H, \Omega) \mapsto
(H \mathbin{++} [\Collapse(x, c, o, \rho, \kappa)],\; \Omega),
\]
subject to the admissibility precondition $\Gamma \vdash t : \Adm(T)$.
Collapse requires an admissibility certificate. A term whose type is not
$\Adm(\cdot)$ cannot be collapsed---this is a hard type error, not a runtime
exception. Unlike the earlier presentation of this operator as bare
``observation,'' Collapse here is the \emph{committed} act of adopting a
task-relative quotient as constraining for later computation; the pure act of
looking that produces $o$ in the first place is not itself an event and is
treated separately in Section~\ref{sec:view-vs-collapse}.
\end{definition}

\section{VIEW versus COLLAPSE}
\label{sec:view-vs-collapse}

The decisive repair this edition makes is to separate pure observation from
committed Collapse. A definition of Collapse as ``observation under a rule''
conflates two operations that later chapters---particularly the operational
semantics and the Rust implementation---already treat differently in practice:
a pure query that returns a value without touching history, and a committed
event that appends to history. Making this explicit here prevents the
equivocation from propagating into the minimality argument, the transition
semantics, and the categorical treatment.

\begin{definition}[View: Pure Projection]
The \emph{view} of history $H_t$ under observation $c$ is
\[
\operatorname{View}_c(H_t) = c(H_t) \in O_c.
\]
VIEW does not append an event, does not change operational state, and does not
consume possibility. Looking is not a world transition:
\[
\operatorname{View}_c(H_t) = o \quad \not\Rightarrow \quad H_{t+1} \neq H_t.
\]
\end{definition}

\begin{definition}[Attest]
An \emph{attestation} is a tuple
\[
\operatorname{Attest}(c, o, \rho, \kappa),
\]
where $o = c(H_t)$, $\rho$ records provenance (who or what produced the
observation and by what mechanism), and $\kappa$ states the task or constraint
for which the observation is being admitted. An attestation is a precondition
for Collapse, not itself an event on $H$.
\end{definition}

\begin{definition}[Collapse as Committed Quotient Adoption]
\[
\operatorname{Attest}(c, o, \rho, \kappa)
\quad\Rightarrow\quad
H_{t+1} = H_t \mathbin{++} [\Collapse(z, c, o, \rho, \kappa)].
\]
\end{definition}

VIEW is not a fifth primitive operator, because it is a meta-level query rather
than an event-producing transformation; it never appears on the left-hand side
of a world transition. Collapse remains a primitive because adopting a quotient
as constraining is irreducible to Pop, Refuse, or Bind: none of those three
records that a projection's result has been accepted as a boundary on later
computation. The four-operator calculus is therefore preserved, with VIEW
sitting outside it as the pure function every Collapse presupposes.

\section{Why These Four and No Others}

The minimality of the four operators follows from an analysis of what operations
on possibility space are conceptually irreducible.

Pop is irreducible because commitment---the foreclosure of alternative
futures---cannot be expressed as a combination of documentation, dependency
recording, or observation. It is the primitive act of narrowing possibility.

Refuse is irreducible because documented inadmissibility is distinct from
commitment: it records that a path was considered and found inadmissible, which
is information that commitment cannot record. A commit of the alternative does not
document why the refused path was refused. Reasons are evidence that survives
collapse; without them, refusal is indistinguishable from non-consideration.

Bind is irreducible because dependency recording---coupling two trajectory
elements---cannot be expressed as commitment or refusal. A bind creates a
relationship between elements without consuming either; no combination of Pop
and Refuse produces this effect.

Collapse is irreducible because observation under a specific rule is distinct from
all three of the above. Pop, Refuse, and Bind all operate on the history and
option space; Collapse maps from history space to observable state space under
a specified projection. It is the meta-operation that makes histories visible.

\begin{proposition}[Representation Adequacy of the Four Operators]
\label{thm:completeness}
The four operators $\{P_x, R_{x,r}, B_{a,b}, C_{x,c,o,\rho,\kappa}\}$ suffice to
encode the following constructs: variable assignment, conditional branching,
function application, exception handling, memory deallocation, and concurrent
event recording.
\end{proposition}

\begin{remark}[Why this is a proposition, not a theorem]
This statement is downgraded here from ``theorem'' to a representation claim.
Exhibiting encodings for six constructs demonstrates representation adequacy
for those constructs; it does not establish general completeness over an
arbitrary source language. A genuine completeness theorem would need to fix a
source calculus $\mathcal{L}$, a compositional translation
$\llbracket - \rrbracket : \mathcal{L} \to \mathsf{Spherepop}$, and an
operational correspondence theorem relating reduction in $\mathcal{L}$ to
reduction of the translated term in Spherepop. No such calculus or
correspondence theorem is given in this manuscript; ``Completeness of the Four
Operators'' should therefore be read as a research proposition motivating
future work, not as an established result.
\end{remark}

\begin{proof}[Proof of representation adequacy]
We establish expressibility for each construct.

\textit{Variable assignment $x := v$}: Express as $P_v$ followed by $B_{x,v}$.
The commitment records the value choice; the bind records the name-value
dependency.

\textit{Conditional branch on predicate $\phi$}: Express as $P_{\phi^+}$ if $\phi$
holds, or $R_{\phi^+,\mathsf{ConstraintViolation}(\phi)}$ followed by
$P_{\phi^-}$ if not. The refusal records why the positive branch was not taken.

\textit{Function application $f(a)$}: Expressed as $\mathsf{LamIntro}(x)$ followed
by $\mathsf{Apply}(a)$---a deferred Pop. The closure freezes the option space at
the moment of abstraction; application is the moment of commitment.

\textit{Exception throwing}: Expressed as $R_{t, r}$ (documented inadmissibility
without system destruction). The type $\Rfsd(T, r)$ carries the reason as a
first-class certificate.

\textit{Memory deallocation}: Expressed as $R_{p, \mathsf{Deallocated}}$ where
$p$ is the pointer symbol. The refusal records that the memory region is no longer
admissible to access, making use-after-free a type-level error.

\textit{Concurrent event recording}: Expressed via Meld, the monoidal composition
of two histories. Two concurrent worlds $(H_1, \Omega_1)$ and $(H_2, \Omega_2)$
are combined as $(H_1 \mathbin{++} H_2, \Omega_1 \cup \Omega_2)$ under an
appropriate ordering convention.
\end{proof}

\section{A Genuine Completeness Result: Boolean Circuits}
\label{sec:boolean-completeness}

The remark above named exactly what a real completeness theorem would need: a
fixed source calculus, a compositional translation into Spherepop, and a
correspondence theorem showing the translation is faithful. For one restricted
but well-understood source calculus---Boolean circuits---all three are
available, and the result below is a genuine theorem rather than a
representation claim. It does not close the general completeness question left
open above (an arbitrary imperative or functional language is not a Boolean
circuit), but it is real, provable progress on a well-defined fragment, and
it doubles as a clean illustration of the VIEW/Collapse distinction from
Section~\ref{sec:view-vs-collapse} in a setting small enough to check by hand.

\begin{definition}[Boolean Interpretation of the Four Operators]
\label{def:boolean-interpretation}
Fix the two-element domain $\{0, 1\}$ and define a map $\llbracket - \rrbracket_B$
from Spherepop's primitives to Boolean operations:
\[
\llbracket P \rrbracket_B(x) = x, \qquad
\llbracket R \rrbracket_B(x) = \neg x, \qquad
\llbracket B \rrbracket_B(x, y) = x \wedge y, \qquad
\llbracket C \rrbracket_B(e) = \mathrm{eval}_B(e),
\]
where $\mathrm{eval}_B$ recursively evaluates a term built from $P$, $R$, $B$
under this interpretation. Under this reading, Pop supplies a committed
Boolean input; Refuse realizes complement; Bind realizes conjunction; and
Collapse \emph{projects} the resulting judgment---a VIEW in the sense of
Section~\ref{sec:view-vs-collapse}, not yet a commitment. The judgment becomes
a committed Collapse event only once it is attested: $\mathrm{eval}_B(e)$ is a
pure read of what the circuit currently evaluates to, and
$\operatorname{Attest}(\llbracket - \rrbracket_B, \mathrm{eval}_B(e), \rho, \kappa)$
is what would make adopting that Boolean judgment as binding an actual
Collapse event, exactly as in the general case.
\end{definition}

\begin{definition}[Boolean Circuit Calculus $\mathcal{L}_{\mathrm{Bool}}$]
Fix $\mathcal{L}_{\mathrm{Bool}}$ to be the calculus of Boolean circuits over a
finite set of Boolean-valued inputs, built from $\neg$ and $\wedge$ with
arbitrary fan-out (i.e., a shared input may feed multiple gates). This is
source calculus enough for the theorem below; the translation
$\llbracket - \rrbracket : \mathcal{L}_{\mathrm{Bool}} \to \mathsf{Spherepop}$
sends a circuit built from $\neg$ and $\wedge$ to the corresponding term built
from $R$ and $B$, with each circuit input realized by a Pop and the circuit's
output read by Collapse under the interpretation of
Definition~\ref{def:boolean-interpretation}.
\end{definition}

\begin{theorem}[Boolean Completeness of Refuse and Bind]
\label{thm:boolean-completeness}
Every Boolean function $f : \{0,1\}^n \to \{0,1\}$ is expressible as
$C(t)$ for some Spherepop term $t$ built only from $P$, $R$, and $B$ (no use of
$\Collapse$ except at the outermost position).
\end{theorem}

\begin{proof}
This is the translation of a classical fact---$\{\neg, \wedge\}$ is a
functionally complete basis for Boolean logic---through
$\llbracket - \rrbracket$. Every Boolean function has a disjunctive normal
form, a disjunction of conjunctions of literals; $\vee$ is eliminable in favor
of $\neg$ and $\wedge$ via De Morgan's law, $x \vee y = \neg(\neg x \wedge \neg
y)$, so every Boolean function is expressible using only $\neg$ and $\wedge$
applied to the input variables. Each input variable is realized by a Pop; each
$\neg$ becomes an $R$; each $\wedge$ becomes a $B$; the whole expression is
read by a single outer $C$. This is exactly $\llbracket - \rrbracket$ applied
to the DNF circuit for $f$, which lies in $\mathcal{L}_{\mathrm{Bool}}$.
\end{proof}

\begin{table}[h]
\centering
\begin{tabular}{lll}
\toprule
Gate & Construction & Formula \\
\midrule
Buffer & $A \to P \to C \to Y$ & $C(P(A)) = A$ \\
Not & $A \to R \to C \to Y$ & $C(R(A)) = \neg A$ \\
And & $A, B \to B \to C \to Y$ & $C(B(A,B)) = A \wedge B$ \\
Nand & $A, B \to B \to R \to C \to Y$ & $C(R(B(A,B))) = \neg(A \wedge B)$ \\
Nor & $A, B \to R, R \to B \to C \to Y$ & $C(B(R(A),R(B))) = \neg(A \vee B)$ \\
Or & $A, B \to R, R \to B \to R \to C \to Y$ & $C(R(B(R(A),R(B)))) = A \vee B$ \\
\bottomrule
\end{tabular}
\caption{Canonical gate constructions from $P$, $R$, $B$, $C$. Nor needs no
final negation because $\neg A \wedge \neg B = \neg(A \vee B)$ already, by De
Morgan; Or is Nor negated once more.}
\end{table}

\begin{remark}[Xor and Xnor, and a correction to an earlier informal diagram]
An informal circuit diagram developed alongside this manuscript rendered Xor
and Xnor with a typeset formula carrying one stray closing parenthesis---an
artifact of how it was generated, not a conceptual error, since the diagrammed
gate-level wiring (two Nots feeding a two-stage And/Or lattice) is a standard,
correct realization. The cleanly parenthesized forms, verified against the
truth table below, are
\[
\mathrm{Xor}(A,B) = C\bigl(R(B(\,R(B(A,R(B))),\; R(B(R(A),B))\,))\bigr),
\]
\[
\mathrm{Xnor}(A,B) = C\bigl(R(B(\,R(B(A,B)),\; R(B(R(A),R(B)))\,))\bigr),
\]
obtained by realizing $A \oplus B = (A \wedge \neg B) \vee (\neg A \wedge B)$
and its negation directly as nested $\mathrm{Or}$/$\mathrm{Nand}$-style
compositions of the constructions above. As with any Boolean function, other
gate-count-minimal realizations of Xor and Xnor exist; Theorem
\ref{thm:boolean-completeness} guarantees expressibility, not a unique or
minimal circuit.
\end{remark}

\begin{center}
\begin{tabular}{cccccccc}
\toprule
$A$ & $B$ & And & Or & Xor & Nand & Nor & Xnor \\
\midrule
0 & 0 & 0 & 0 & 0 & 1 & 1 & 1 \\
0 & 1 & 0 & 1 & 1 & 1 & 0 & 0 \\
1 & 0 & 0 & 1 & 1 & 1 & 0 & 0 \\
1 & 1 & 1 & 1 & 0 & 0 & 0 & 1 \\
\bottomrule
\end{tabular}
\end{center}

\begin{remark}[What this does and does not settle]
Theorem~\ref{thm:boolean-completeness} is a genuine completeness theorem for
$\mathcal{L}_{\mathrm{Bool}}$: the source calculus, the translation, and the
correspondence (DNF expressibility, then a literal term-for-gate substitution)
are all fully specified, unlike the six-construct representation-adequacy
claim above. It does not, by itself, upgrade that broader claim to a theorem:
Boolean circuits are combinational and have no notion of sequencing,
admissibility, or refusal-with-reasons, so this result says nothing about
whether an arbitrary program in some fixed imperative or functional source
language translates faithfully into Spherepop. What it does establish is that
the concern behind Proposition~\ref{thm:completeness}---that six worked
examples might not generalize---does not apply uniformly: for at least one
natural and non-trivial fragment, the four operators (in fact, only two of
them) are provably complete, not merely adequate for the cases checked.
\end{remark}

\section{The Conservation Law}

The four operators satisfy a conservation law that reflects the fundamental
structure of possibility.

\begin{theorem}[Conservation of Possibility]
\label{thm:conservation}
For any legal execution sequence
\[
W_0 \xrightarrow{e_1} W_1 \xrightarrow{e_2} \cdots \xrightarrow{e_n} W_n
\]
where each step is one of $P_x$, $R_{x,r}$, $B_{a,b}$, or $C_{x,c}$:
\begin{enumerate}
\item $|H_t|$ is strictly monotonically increasing: $|H_{t+1}| = |H_t| + 1$.
\item $|\Omega_t|$ is monotonically non-increasing.
\item Under pure Pop dynamics, $|H_t| + |\Omega_t| = |\Omega_0|$ for all $t$.
\end{enumerate}
\end{theorem}

\begin{proof}
(i) Every operator appends exactly one event to $H$, so $|H_{t+1}| = |H_t| + 1$.

(ii) $P_x$ removes $x$ from $\Omega$; the remaining operators do not modify
$\Omega$. Hence $|\Omega_{t+1}| \leq |\Omega_t|$.

(iii) Under pure Pop dynamics, every step removes exactly one element from $\Omega$
and appends one event to $H$. Starting from $|H_0| = 0$ and $|\Omega_0|$, after
$t$ steps $|H_t| = t$ and $|\Omega_t| = |\Omega_0| - t$, giving
$|H_t| + |\Omega_t| = |\Omega_0|$.
\end{proof}

The conservation law has the structure of a physical conservation principle:
possibility is not created or destroyed---it is either consumed (by Pop) or
documented as inadmissible (by Refuse). The operators Bind and Collapse record
relationships and observations without affecting the possibility budget.

\begin{definition}[Selection-Budget Invariant]
\label{def:selection-budget}
Define the event weight function $w : E \to \mathbb{N}$ by $w(\Pop(x)) = 1$,
$w(\Refuse(x,r)) = 0$, $w(\Bind(a,b)) = 0$, $w(\Collapse(x,c,o,\rho,\kappa)) = 0$.
For a world $(H, \Omega)$, the \emph{selection-budget invariant} is
\[
\Pi(H, \Omega) = |\Omega| + \sum_{e \in H} w(e).
\]
This quantity was originally called the ``generalised possibility functional,''
but that name overstates what it measures; see the remark below.
\end{definition}

\begin{theorem}[Conservation of the Selection-Budget Invariant]
For any legal execution sequence, $\Pi(H_t, \Omega_t) = |\Omega_0|$ for all $t$.
\end{theorem}

\begin{proof}
At each step, exactly one of four operators is applied. Pop: $|\Omega|$ decreases
by 1; $w(\Pop) = 1$ is added. Net change: $(-1) + 1 = 0$. Refuse, Bind, Collapse:
$|\Omega|$ unchanged; $w = 0$ is added. Net change: $0$. In every case
$\Pi(H_{t+1}, \Omega_{t+1}) = \Pi(H_t, \Omega_t)$.
\end{proof}

\begin{remark}[$\Pi$ does not measure reachable possibility]
Because Refuse and Collapse can constrain admissibility, $\calA(H)$, without
changing $\Omega$ (Section~\ref{sec:admissibility-set} below), $\Pi$ is
conserved trivially by construction and does not track total reachable
possibility: a history full of refusals that narrow $\calA(H)$ down to a
single admissible symbol has the same $\Pi$ as one with no refusals at all,
even though the two differ enormously in what remains reachable. $\Pi$ should
therefore be read as a bookkeeping identity over $(\Pop\text{-consumption}) +
(\text{remaining options})$, not as a general reachability measure. A separate
admissible-continuation functional, defined over $\calA(H)$ rather than
$\Omega$ alone, would be needed to track reachability properly; none is
developed in this manuscript. The committed Collapse event, in particular, can
narrow task-relative admissibility even when $|\Omega|$ is unchanged, precisely
because it commits the world to treating a quotient as constraining---this is
invisible to $\Pi$ by construction, not because Collapse has no effect.
\end{remark}

\begin{corollary}[Irreversibility of Execution]
For any non-empty execution sequence, $|H_{t+1}| > |H_t|$. Therefore no legal
execution can return to a previous world state $(H_t, \Omega_t)$.
\end{corollary}

\begin{proof}
Every operator appends exactly one event to $H$, so $|H_{t+1}| = |H_t| + 1 > |H_t|$.
Since $H$ grows strictly, $(H_{t+1}, \Omega_{t+1}) \neq (H_t, \Omega_t)$.
\end{proof}

This irreversibility is not a limitation but a feature. It is the formal
correlate of the phenomenological observation that histories are primary: a world
that could return to a prior state would be a world in which the history of
the intervening steps was undone, which is not computation but its negation.

\chapter{The Algebra of Histories}

\section{Histories as the Primary Objects}

Before introducing types, evaluators, or compilers, we must characterize the
mathematical structure of the objects that are primary in the Spherepop ontology.
Histories are not merely lists of events. They form a category, and the category
structure constrains what operations on histories are meaningful.

\begin{definition}[History]
A \emph{history} over event alphabet $E$ is a finite sequence
$H = [e_1, e_2, \ldots, e_n]$ with each $e_i \in E$. The empty history is
$\varepsilon = []$.
\end{definition}

\begin{definition}[History Concatenation]
For histories $H_1 = [e_1, \ldots, e_m]$ and $H_2 = [f_1, \ldots, f_n]$,
\[
H_1 \mathbin{++} H_2 = [e_1, \ldots, e_m, f_1, \ldots, f_n].
\]
\end{definition}

\begin{proposition}[History Monoid]
The triple $(\calH, \mathbin{++}, \varepsilon)$ is the free monoid over $E$.
\end{proposition}

\begin{proof}
Associativity of $\mathbin{++}$ follows from list concatenation associativity.
$\varepsilon$ is the left and right identity. Freeness: every history $H$ is
uniquely determined by its event sequence, and no non-trivial relation among
generators holds.
\end{proof}

\section{The History Category}

\begin{definition}[History Category $\Hist$]
Define the category $\Hist$ as follows:
\begin{itemize}
\item \textit{Objects}: option spaces $\Omega \subseteq \Omega_0$.
\item \textit{Morphisms}: $\mathrm{Hom}(\Omega, \Omega')$ is the set of histories
$H$ such that starting from $\Omega$, executing $H$ yields $\Omega'$.
\item \textit{Composition}: $H_2 \circ H_1 = H_1 \mathbin{++} H_2$ (execute
$H_1$ first, then $H_2$).
\item \textit{Identity}: $\mathrm{id}_\Omega = \varepsilon$.
\end{itemize}
\end{definition}

\begin{proposition}[$\Hist$ is a Well-Defined Category]
$\Hist$ satisfies all category axioms.
\end{proposition}

\begin{proof}
Associativity of composition follows from associativity of $\mathbin{++}$.
$\varepsilon$ is the identity morphism since $H \mathbin{++} \varepsilon = H =
\varepsilon \mathbin{++} H$.
\end{proof}

\begin{proposition}[$\Hist$ is a Free Strict Monoidal Category]
$(\Hist, \mathbin{++}, \varepsilon, \otimes)$ is a free strict monoidal category
where the tensor product $H_1 \otimes H_2$ is the Meld operation (parallel
composition of histories) and the monoidal unit is the empty option space with
empty history.
\end{proposition}

\begin{remark}
The free monoid structure forces \texttt{append} (extend by one generator) and
\texttt{meld} (monoidal composition of two histories) as the only legitimate
history operations. The absence of \texttt{remove} and \texttt{undo} is the
mechanization of freeness: there are no relations among generators.
\begin{lstlisting}[style=rust, caption={History as free monoid in Rust}]
pub fn append(&mut self, event: Event) {
    self.events.push(event);  // one generator
}
pub fn meld(&mut self, other: &History) {
    self.events.extend(other.events.iter().cloned());
}
\end{lstlisting}
History equality $H_1 = H_2$ is the primary criterion of program identity.
\end{remark}

\section{Unique Factorisation}

\begin{theorem}[Unique History Factorisation]
Every non-empty history $H = [e_1, e_2, \ldots, e_n]$ admits a unique
factorisation into primitive generators:
\[
H = [e_1] \mathbin{++} [e_2] \mathbin{++} \cdots \mathbin{++} [e_n].
\]
No other factorisation into generators exists.
\end{theorem}

\begin{proof}
Since $\calH$ is the free monoid over $E$, every element has a unique normal form
as a sequence of generators. The generators are singletons $[e]$; freeness of the
monoid guarantees there are no non-trivial relations among generators.
\end{proof}

\begin{remark}
Unique factorisation is what makes history equivalence a decidable correctness
criterion for compilation. If two execution paths produce the same history, they
must have produced the same sequence of primitive events in the same order.
\end{remark}

\section{Admissibility as a Set-Valued Function}
\label{sec:admissibility-set}

\begin{definition}[Admissible Futures]
For a history $H$, the set of admissible future symbols is
\[
\calA(H) = \Omega_0 \setminus \{x \mid \exists r.\, \Refuse(x, r) \in H\}.
\]
\end{definition}

\begin{definition}[Admissibility Contraction]
Let $H' = H \mathbin{++} [\Refuse(x, r)]$. Then $\calA(H') \subseteq \calA(H)$,
and specifically $x \notin \calA(H')$ while $x$ may have been in $\calA(H)$.
\end{definition}

\begin{remark}
The asymmetry between Pop and Refuse is critical. Pop contracts the structural
option space $\Omega$. Refuse contracts the admissibility set $\calA(H)$ without
touching $\Omega$. Two execution paths may have the same $\Omega$ while having
different $\calA(H)$---they are structurally equivalent but semantically
distinguished by their refusal histories. This distinction is invisible to any
state-centric computational model.
\end{remark}

%% ============================================================
%% PART III: OPERATIONAL SEMANTICS AND COLLAPSE
%% ============================================================
\part{Operational Semantics and Collapse}

\chapter{Small-Step Semantics}

\section{The Transition System}

We write $(H, \Omega) \xrightarrow{\alpha} (H', \Omega')$ for a single execution
step labelled $\alpha$.

\[
\frac{x \in \Omega}{(H, \Omega) \xrightarrow{\mathsf{pop}(x)} (H \mathbin{++} [\Pop(x)],\; \Omega \setminus \{x\})} \quad [\mathrm{POP}]
\]

\[
\frac{}{(H, \Omega) \xrightarrow{\mathsf{refuse}(x,r)} (H \mathbin{++} [\Refuse(x, r)],\; \Omega)} \quad [\mathrm{REFUSE}]
\]

Note: $[\mathrm{REFUSE}]$ has no premise---any symbol may be refused. The effect
on the option space is nil; the effect on admissibility is non-nil.

\[
\frac{}{(H, \Omega) \xrightarrow{\mathsf{bind}(a,b)} (H \mathbin{++} [\Bind(a, b)],\; \Omega)} \quad [\mathrm{BIND}]
\]

\[
\frac{\Gamma \vdash t : \Adm(T) \quad \operatorname{Attest}(c, o, \rho, \kappa)}{(H, \Omega) \xrightarrow{\mathsf{collapse}(x,c,o,\rho,\kappa)} (H \mathbin{++} [\Collapse(x, c, o, \rho, \kappa)],\; \Omega)} \quad [\mathrm{COLLAPSE}]
\]

The premise $\Gamma \vdash t : \Adm(T)$ is the type-theoretic admissibility guard.
Collapse requires an admissibility certificate. Without it, collapse is a hard
type error.

Alongside $[\mathrm{COLLAPSE}]$ we record a separate, non-event judgment for
pure observation, matching the \texttt{view}/\texttt{collapse} split of
Section~\ref{sec:view-vs-collapse}:
\[
\frac{}{W \Downarrow_c o} \quad [\mathrm{VIEW}]
\]
where $o = c(H)$. $[\mathrm{VIEW}]$ does not appear on the left of an
arrow-labelled transition and contributes no event to $H$; it exists only to
make explicit, at the level of the transition system, that the earlier draft's
single \texttt{observe(\&self, ...)} method (Chapter~\ref{ch:compiler}) was
already pure in the Rust implementation while its formal counterpart
$[\mathrm{COLLAPSE}]$ appended an event---an inconsistency between
implementation and semantics that this separation resolves rather than
papers over. $[\mathrm{VIEW}]$ is excluded from the conservation law of
Theorem~\ref{thm:conservation} because it is not an execution event; only the
committed $[\mathrm{COLLAPSE}]$ transition is.

\section{Big-Step Semantics}

The big-step relation $\langle t, W \rangle \Downarrow \langle v, W' \rangle$
relates terms and worlds to values and resulting worlds.

\begin{definition}[Evaluation Result]
An \emph{evaluation result} is a triple $(v, H, a)$ where $v$ is a value,
$H$ is the accumulated history, and $a \in \{\top, \bot\}$ is the admissibility
flag. The admissibility flag is $\top$ if every operation in $H$ was admissible
and $\bot$ if any refusal was encountered.
\end{definition}

The key invariant is that evaluation never discards history. Every step appends
to $H$; no step modifies or removes prior events. The world grows monotonically.

\begin{lstlisting}[style=rust, caption={Evaluation result preserving all history}]
pub struct EvalResult {
    pub value: Value,
    pub world: World,    // world.history grows monotonically
    pub admissible: bool,
}
\end{lstlisting}

\section{Multi-Step Reduction}

\begin{definition}[Multi-Step Reduction]
The multi-step reduction $\to^*$ is the reflexive transitive closure of $\to$.
A term $t$ is \emph{normal} if there is no $t'$ with $t \to t'$.
\end{definition}

\begin{theorem}[History Monotonicity]
If $(H_0, \Omega_0) \to^* (H_n, \Omega_n)$ then $H_0$ is a prefix of $H_n$:
there exists $H'$ such that $H_n = H_0 \mathbin{++} H'$.
\end{theorem}

\begin{proof}
By induction on the length of the reduction sequence. Each step appends exactly
one event to the history. The result follows by transitivity.
\end{proof}

\begin{theorem}[Option Space Monotonicity]
If $(H_0, \Omega_0) \to^* (H_n, \Omega_n)$ then $\Omega_n \subseteq \Omega_0$.
\end{theorem}

\begin{proof}
Only Pop removes elements from $\Omega$; all other operators leave it unchanged.
By induction, $\Omega$ can only shrink.
\end{proof}

\chapter{Collapse as Quotient}

\section{The Insufficiency of Unparameterized Collapse}

The original Spherepop design represented collapse as simply \texttt{Collapse(Symbol)}.
This was insufficient, and the insufficiency is philosophically illuminating.

Observable state depends on how you look. The same history, examined by different
observers with different observational frameworks, yields different observable
states. A collapse without a rule is a collapse into an unspecified notion of
visibility. The rule is not an implementation detail; it is what makes the
observation meaningful.

\begin{definition}[Collapse Rule]
A \emph{collapse rule} is a function $c : \calH \to O_c$ from the history monoid
to an observational space $O_c$.
\end{definition}

\begin{definition}[Observational Equivalence]
Two histories $H_1, H_2 \in \calH$ are \emph{observationally equivalent under
rule $c$}, written $H_1 \sim_c H_2$, if and only if $c(H_1) = c(H_2)$.
\end{definition}

Three distinct objects are gathered under the informal phrase ``observable
state,'' and it is worth separating them explicitly, since the manuscript's
earlier presentation moved between them without comment:

\begin{itemize}
\item the \emph{attainable observation values}, $O_c = \operatorname{im}(c)$;
\item the \emph{equivalence class} of a history, $[H]_c \in \calH / {\sim_c}$;
\item the \emph{actual observation}, $o = c(H)$, a single element of $O_c$.
\end{itemize}

\begin{definition}[Observable State as Quotient]
Strictly, $O_c$ need not literally equal $\calH / {\sim_c}$ as a set; it is
canonically isomorphic to the quotient once $c$ is treated as surjective onto
its own image (which it always is, tautologically, once $O_c$ is taken to be
$\operatorname{im}(c)$ rather than some a priori larger codomain):
\[
O_c \;\cong\; \calH / {\sim_c}, \qquad [H]_c \;\longleftrightarrow\; c(H).
\]
Under this identification, an observable state may be treated interchangeably
as an equivalence class of histories or as the value $c(H)$ that names it.
\end{definition}

Observable state is not a state. Observable state is a quotient. Two histories
that are globally distinct may be observationally indistinguishable under a
specific collapse rule. This is a feature: the collapse rule specifies which
distinctions are relevant to a given question.

\section{The Four Canonical Rules}

\begin{definition}[Identity Rule]
$c_I : \calH \to \calH$ is the identity: $c_I(H) = H$. Under the identity rule,
every event is visible and $O_I \cong \calH$.
\end{definition}

\begin{definition}[LastWrite Rule]
$c_{LW}$ maps $H$ to the most recent pop for each symbol, with subsequent refusals
retroactively removing symbols:
\[
c_{LW}(H) = \{x \mapsto 1 \mid \Pop(x) \in H \text{ and } \nexists r.\, \Refuse(x,r) \in H \text{ after the last } \Pop(x)\}.
\]
\end{definition}

\begin{definition}[Accumulate Rule]
$c_{\mathrm{acc}}(H)$ maps each symbol to the count of its Pop events:
\[
c_{\mathrm{acc}}(H)(x) = |\{i \mid e_i = \Pop(x)\}|.
\]
\end{definition}

\begin{definition}[Projection Rule]
For a label prefix $L$, the projection rule $\pi_L$ restricts to events whose
symbol names share prefix $L$:
\[
\pi_L(H) = [e_i \in H \mid \mathrm{sym}(e_i) \text{ starts with } L].
\]
\end{definition}

\begin{proposition}[Coarsening Order]
The collapse rules form a partial order by coarsening: $c_1 \leq c_2$ iff
${\sim_{c_2}} \subseteq {\sim_{c_1}}$ (finer equivalence = more informative
observation). Under this order: $c_I$ is the finest; $\pi_L$ is coarser than
$c_I$; $c_{LW}$ and $c_{\mathrm{acc}}$ are incomparable in general.
\end{proposition}

\section{The Collapse Functor}

\begin{definition}[Observable-State Category $\Obs_c$]
For a collapse rule $c$, define the category $\Obs_c$:
\begin{itemize}
\item \textit{Objects}: observable states $o \in O_c$.
\item \textit{Morphisms}: $\mathrm{Hom}(o_1, o_2)$ is the set of observable
transitions transforming $o_1$ into $o_2$ under rule $c$.
\item \textit{Composition}: sequential observable transition.
\end{itemize}
\end{definition}

\begin{definition}[Collapse Functor]
For collapse rule $c$, define $F_c : \Hist \to \Obs_c$ by
$F_c(\Omega) = c(\varepsilon_\Omega)$ on objects and $F_c(H) = c(H)$ on morphisms.
\end{definition}

\begin{theorem}[Functoriality of Collapse]
$F_c$ is a well-defined functor when $c$ respects composition:
\[
c(H_2 \mathbin{++} H_1) = c(H_2) \circ c(H_1).
\]
\end{theorem}

\begin{proof}
Identity preservation: $F_c(\varepsilon) = c(\varepsilon) = \mathrm{id}_{O_c}$.
Composition preservation: $F_c(H_2 \circ H_1) = F_c(H_1 \mathbin{++} H_2) =
c(H_1 \mathbin{++} H_2) = c(H_2) \circ c(H_1) = F_c(H_2) \circ F_c(H_1)$ when
$c$ respects composition.
\end{proof}

\begin{remark}
The functoriality of collapse elevates it from an implementation detail to a
category-theoretic object. Collapse is a structure-preserving map between
categories. The collapse rule specifies which structures are preserved, and
different rules preserve different structures. $c_{LW}$ is not generally functorial
because the last-write of a concatenated history is not the composition of the
last-writes of its parts---a later segment can retroactively affect the observable
state of an earlier one.
\end{remark}

\begin{remark}[Scope of the functoriality claim]
It bears restating explicitly that Theorem~1.101 (Functoriality of Collapse) is
conditional: an arbitrary observation function $c : \calH \to O_c$ does not
automatically induce a functor. Functoriality holds precisely when $c$ is a
monoid homomorphism, $c(H_2 \mathbin{++} H_1) = c(H_2) \circ c(H_1)$, and fails
otherwise, as the $c_{LW}$ example immediately above shows. ``Collapse Functor
Functoriality,'' entered in the theorem status table as proved, should be read
as proved \emph{for composition-respecting collapse rules}, not as a general
statement about all observation functions. A more ambitious construction---a
quotient category built from a congruence on histories, requiring
$H_1 \sim_c H_1'$ and $H_2 \sim_c H_2'$ to imply $H_1 H_2 \sim_c H_1' H_2'$---is
not attempted here and remains open.
\end{remark}

\section{The Universal Property of Collapse}

\begin{theorem}[Universal Property of Collapse]
Let $q_c : \calH \to O_c$ be the quotient map for rule $c$. If $f : \calH \to X$
is any function satisfying
\[
H_1 \sim_c H_2 \implies f(H_1) = f(H_2),
\]
then there exists a unique $\bar{f} : O_c \to X$ such that $f = \bar{f} \circ q_c$.
\end{theorem}

\begin{proof}
By the universal property of quotient sets. Define $\bar{f}([H]_c) = f(H)$.
This is well-defined because $f$ is constant on equivalence classes. Uniqueness:
any $\bar{f}$ satisfying $f = \bar{f} \circ q_c$ must map $[H]_c$ to $f(H)$.
\end{proof}

\begin{remark}
The universal property says that observable state is not one representation among
many. It is the universal representation among all representations that identify
the same histories as $c$ does. Every coarser observation necessarily passes
through $O_c$.
\end{remark}

%% ============================================================
%% PART IV: TYPE THEORY
%% ============================================================
\part{Type Theory}

\chapter{Types as Admissibility Certificates}

\section{The Inadequacy of Value-Set Types}

The standard interpretation of types---a type is a set of admissible values---is
adequate for state-centric languages. It is insufficient for Spherepop.

Consider what a type must certify in a history-primary setting. A value $v$ of
type $T$ is not merely a member of a set. It is the result of a path through
history that reached $v$ without violating any constraint represented by $T$.
The type is not a description of $v$ at a moment; it is a certificate about the
history that produced $v$.

\section{The Type Language}

\begin{definition}[Spherepop Types]
The type language of Spherepop is defined inductively:
\begin{align*}
T ::=\; & \mathsf{Unit} \mid \mathsf{Never} \mid X \\
\mid\; & \Pi(x : T_1).T_2 \\
\mid\; & \Adm(T) \\
\mid\; & \Rfsd(T, r) \\
\mid\; & \Col(T, c) \\
\mid\; & \Proc(T_1 \rightsquigarrow T_2)
\end{align*}
The simple function type $T_1 \to T_2$ is notation for $\Pi(\_ : T_1).T_2$.
\end{definition}

The distinctive types carry certificates, not just values:
\begin{itemize}
\item $\Adm(T)$: the path that produced this term was admissible.
\item $\Rfsd(T, r)$: the term was encountered but found inadmissible for reason
$r$. The reason is part of the type---two terms refused for different reasons
have genuinely different types.
\item $\Col(T, c)$: an admissible term was observed under collapse rule $c$ and
the observation is sealed. The rule $c$ is part of the type.
\end{itemize}

\section{The Type Rules}

We write $\Gamma \vdash t : T$ for the standard typing judgement.

\[
\frac{\Gamma \vdash t : T}{\Gamma \vdash \mathsf{pop}(t) : \Adm(T)} \quad [\mathrm{T\text{-}POP}]
\]

Pop transforms any typed term into an admissible term. It is the mechanism by
which structural commitment produces a reachability certificate.

\[
\frac{\Gamma \vdash t : T \quad r \in \mathsf{RefusalReason}}{\Gamma \vdash \mathsf{refuse}(t, r) : \Rfsd(T, r)} \quad [\mathrm{T\text{-}REFUSE}]
\]

Refusal is typed: the type carries the reason. A term refused for
$\mathsf{ConstraintViolation}(c)$ has a type formally distinct from one refused
for $\mathsf{NotAvailable}$.

\[
\frac{\Gamma \vdash t : \Adm(T) \quad c \in \mathsf{CollapseRule}}{\Gamma \vdash \mathsf{collapse}(t, c) : \Col(T, c)} \quad [\mathrm{T\text{-}COLLAPSE}]
\]

Collapse requires an admissibility certificate. A term whose type is not
$\Adm(\cdot)$ cannot be collapsed---this is a hard type error, not a runtime
exception.

\[
\frac{\Gamma \vdash a : T_1 \quad \Gamma \vdash b : T_2}{\Gamma \vdash \mathsf{bind}(a, b) : \Proc(T_1 \rightsquigarrow T_2)} \quad [\mathrm{T\text{-}BIND}]
\]

\[
\frac{\Gamma, x : T_1 \vdash b : T_2}{\Gamma \vdash \lambda x : T_1.\, b : \Pi(x : T_1).T_2} \quad [\mathrm{T\text{-}LAM}]
\]

\[
\frac{\Gamma \vdash f : \Pi(x : T_1).T_2 \quad \Gamma \vdash a : T_1}{\Gamma \vdash f\, a : T_2[x := a]} \quad [\mathrm{T\text{-}APP}]
\]

\section{Open Worlds and Closed Worlds}

During implementation of the Spherepop type checker, a conflict emerged between
how the interpreter and type checker treated free variables. The apparent bug
conceals a genuine theoretical discovery: which behavior is correct depends on
which assumption about the boundary of the world is in force.

\begin{definition}[Closed-World Context]
A \emph{closed-world context} $\Gamma_c$ satisfies:
\[
x \notin \Gamma_c \implies \Gamma_c \vdash x : \bot.
\]
The absence of a declaration is the absence of the object.
\[
\frac{x \notin \Gamma \quad \Gamma \text{ is closed-world}}{\Gamma \not\vdash x : T \text{ for any } T} \quad [\mathrm{T\text{-}VAR\text{-}CLOSED}]
\]
\end{definition}

\begin{definition}[Open-World Context]
An \emph{open-world context} $\Gamma_o$ satisfies:
\[
x \notin \Gamma_c \implies \Gamma_o \vdash x : \mathsf{Unit}.
\]
The absence of a declaration is undecided possibility.
\[
\frac{x \notin \Gamma \quad \Gamma \text{ is open-world}}{\Gamma \vdash x : \mathsf{Unit}} \quad [\mathrm{T\text{-}VAR\text{-}OPEN}]
\]
\end{definition}

\begin{proposition}[Monotonicity of Provability]
$\Gamma_c \subseteq \Gamma_o \implies \mathrm{Provable}(\Gamma_c) \subseteq
\mathrm{Provable}(\Gamma_o)$.
\end{proposition}

\begin{remark}
The open/closed world distinction is not an engineering convenience. It is the
mechanization of an epistemological commitment: a closed-world type checker assumes
complete knowledge of the reachable; an open-world type checker assumes partial
knowledge. A language whose fundamental ontology is about possibility space cannot
leave this distinction implicit. The $\mathsf{TypeMode}$ is preserved through
\texttt{extend}: a closed-world context cannot silently become open-world
mid-derivation.
\end{remark}

\section{Principal Theorems}

\begin{theorem}[Collapse Soundness]
If $\Gamma \vdash \mathsf{collapse}(t, c) : T$ then $\Gamma \vdash t : \Adm(T')$
for some $T'$, and $T = \Col(T', c)$.
\end{theorem}

\begin{proof}
By inversion on $[\mathrm{T\text{-}COLLAPSE}]$: the only rule that derives a
Collapse judgement requires the premise $\Gamma \vdash t : \Adm(T)$, so the
inner term must have an admissible type.
\end{proof}

\begin{theorem}[Type Preservation]
If $\Gamma \vdash t : T$ and $t \to t'$, then $\Gamma \vdash t' : T$.
\end{theorem}

\begin{proof}
By structural induction on the reduction relation $\to$. The key cases are
$[\mathrm{T\text{-}POP}]$, $[\mathrm{T\text{-}REFUSE}]$, and
$[\mathrm{T\text{-}COLLAPSE}]$, each mapping a reduction step to an admissibility
transformation closed under the respective rule. Beta-reduction requires a
substitution lemma: $\Gamma, x : T_1 \vdash b : T_2$ and $\Gamma \vdash a : T_1$
implies $\Gamma \vdash b[x := a] : T_2[x := a]$.
\end{proof}

\begin{theorem}[Progress]
If $\vdash t : T$ (well-typed in the empty context) then either $t$ is a value,
or there exists $t'$ such that $t \to t'$.
\end{theorem}

\begin{proof}
By structural induction on $t$. Variable terms are values. Lambda terms are values
(closures). $\mathsf{pop}$, $\mathsf{refuse}$, $\mathsf{bind}$, and
$\mathsf{collapse}$ of a well-typed inner term either reduce (if the inner term
reduces) or produce a value (if the inner term is a value).
\end{proof}

\chapter{Dependent Types and Admissibility Lattices}

\section{Dependent Process Types}

The current $\Proc(T_1 \rightsquigarrow T_2)$ is first-order. A dependent process
type allows the output type to depend on the specific value consumed.

\begin{definition}[Dependent Process Type]
A \emph{dependent process type}
\[
\Pi(x : T_1).\, \Proc(x, f(x))
\]
allows the type of the output process to depend on the value $x$ of the input.
This is the natural extension of dependent type theory into the process calculus
setting.
\end{definition}

\begin{example}
A function that reads a file and returns a proof that the file was read can be
typed as:
\[
\Pi(\mathit{path} : \mathsf{FilePath}).\, \Proc(\mathit{path}, \Adm(\mathsf{FileContents}(\mathit{path}))).
\]
The $\Adm$ in the output type certifies that the reading path was admissible---no
access violation occurred, the file existed, and the permissions were valid.
\end{example}

\section{Admissibility Lattices}

The current admissibility structure is Boolean: a term is admissible or not. A
more refined theory recognizes that admissibility is graded.

\begin{definition}[Admissibility Lattice]
An \emph{admissibility lattice} $(\mathcal{L}, \leq, \top, \bot, \wedge, \vee)$
is a bounded lattice where:
\begin{itemize}
\item $\top$ is full admissibility: the trajectory satisfies all constraints.
\item $\bot$ is total inadmissibility: the trajectory violates all constraints.
\item $\ell_1 \wedge \ell_2$ is the admissibility under the conjunction of
constraints $c_1$ and $c_2$.
\item $\ell_1 \vee \ell_2$ is the admissibility under the disjunction.
\end{itemize}
\end{definition}

\begin{definition}[Graded Admissibility Type]
$\Adm_\ell(T)$ denotes admissibility at level $\ell \in \mathcal{L}$. The
standard $\Adm(T)$ is $\Adm_\top(T)$.
\end{definition}

\begin{proposition}[Admissibility Lattice Typing]
The typing rule for graded admissibility is:
\[
\frac{\Gamma \vdash t : T \quad \ell \in \mathcal{L}}{\Gamma \vdash \mathsf{pop}_\ell(t) : \Adm_\ell(T)}
\]
Collapse requires admissibility at level at least $\ell_0$ for some threshold
$\ell_0$ determined by the collapse rule $c$.
\end{proposition}

\begin{remark}
Admissibility lattices connect the Spherepop type theory to the admissibility
geometry program of the RSVP and CLIO frameworks. The xylomorphic criterion
$\lambda < 1$ from the admissibility field is a threshold condition in a
continuous admissibility lattice, where admissibility is measured by a scalar
$\lambda$ and collapse is only valid when $\lambda$ falls below the threshold.
\end{remark}

\section{Proof-Carrying Refusal}

The $\mathsf{RefusalReason}$ type is currently a vocabulary for inadmissibility.
It is not yet a proof system. A fully realized Spherepop type system would require
$\mathsf{RefusalReason}$ to be a proof term: evidence that can be checked, not
merely a label.

\begin{definition}[Proof-Carrying Refusal]
A \emph{proof-carrying refusal} is a refusal of the form
$\mathsf{refuse}(t, \pi)$ where $\pi$ is a proof term of type
$\mathsf{Inadmissible}(t, \Omega)$---a formal certificate that $t$ is inadmissible
given the current admissibility set $\calA(H)$.
\end{definition}

\begin{definition}[Inadmissibility Proof System]
The inadmissibility propositions and their proofs are:
\begin{align*}
\mathsf{AlreadyRefused}(x) &: x \in \{y \mid \exists r.\, \Refuse(y,r) \in H\} \\
\mathsf{ConstraintViolation}(c, x) &: c(x) > \mathsf{threshold}(c) \\
\mathsf{NotAvailable}(x) &: x \notin \Omega \\
\mathsf{DependencyFailed}(x, y) &: \Bind(x,y) \in H \wedge \mathsf{AlreadyRefused}(y)
\end{align*}
Each proposition has a type-theoretic proof term, and $\mathsf{refuse}(t, \pi)$
is only well-typed when $\pi$ proves one of these propositions for $t$.
\end{definition}

\begin{theorem}[Refusal Completeness]
Every inadmissible trajectory has a proof-carrying refusal at some point.
That is, if $\calA(H) \not\ni x$ for some $x$ that is later attempted, then
there exists a proof term $\pi$ of $\mathsf{Inadmissible}(x, \calA(H))$.
\end{theorem}

\begin{proof}
$x \notin \calA(H)$ implies $\exists r.\, \Refuse(x, r) \in H$. The proof term
$\pi = \mathsf{AlreadyRefused}(x)$ witnesses that $x$ is in the refused set,
which is defined as the complement of $\calA(H)$.
\end{proof}

%% ============================================================
%% PART V: NOTES AS SPHEREPOP OBJECTS
%% ============================================================
\part{Notes as Spherepop Objects}

\chapter{The Seven Note Classes as Operator Realizations}

The four Spherepop operators---Pop, Refuse, Bind, Collapse---are abstract
operations on possibility space. Notes are their physical realization in cognition,
computation, language, and civilization. This chapter develops the correspondence
systematically.

\section{The Classification Principle}

\begin{definition}[Dominant Operator]
The \emph{dominant operator} of a note class is the Spherepop operator that
accounts for the primary function of notes in that class. Most notes mix operators,
but each class has a dominant one.
\end{definition}

\begin{theorem}[Notes as Spherepop Programs]
Every note is a Spherepop program. Every Spherepop program is a note. The
correspondence is given by the following bijection between note classes and
operator combinations:
\begin{center}
\begin{tabular}{lll}
\toprule
Note Class & Dominant Operator & Secondary Operators \\
\midrule
Capture & Refuse & --- \\
Prospective & Bind & Pop (deferred) \\
Repair & Refuse $+$ Pop & --- \\
Coordination & Bind & Refuse \\
Refusal & Refuse & --- (pure form) \\
Generative & Pop & Bind \\
Collapse & Collapse & --- \\
\bottomrule
\end{tabular}
\end{center}
\end{theorem}

\chapter{Capture Notes}

\begin{definition}[Capture Note]
A \emph{capture note} is a note whose primary function is to preserve a
continuation that exists at the time of note creation but is expected to become
inaccessible. In Spherepop terms, a capture note performs a Refuse operation
with respect to the projection $\sigma : \calH \to S$: it preserves history
fibers that the projection would collapse.
\end{definition}

A photograph refuses the collapse of a perceptual trajectory. An audio recording
refuses the collapse of a sonic trajectory. A laboratory notebook refuses the
collapse of a procedural and cognitive trajectory.

\begin{proposition}[Capture Notes and Refuse]
A capture note $N$ performs a Refuse operation with respect to the projection
$\sigma : \calH \to S$ if and only if there exists a history fiber
$h \in \sigma^{-1}(S)$ such that $P_A(c_h, t \mid N) > P_A(c_h, t)$, where
$c_h$ is the continuation associated with history $h$.
\end{proposition}

The quality criterion for a capture note is not fidelity to the original state
but adequacy as a reconstruction substrate. A photograph that captures the correct
person but not the context may be a poor capture note even at high resolution, if
the lost context is what matters for future reconstruction. The adequacy of a
capture note depends on what future continuations are anticipated.

\chapter{Prospective Notes}

\begin{definition}[Prospective Note]
A \emph{prospective note} is a note whose primary function is to preserve access
to a continuation that does not yet exist at note creation time but is anticipated.
In Spherepop terms, a prospective note performs a Bind operation between present
agent state $A_t$ and future continuation $c_{t'}$.
\end{definition}

A calendar entry binds future action to present scheduling. A blueprint binds
future construction to present design. A roadmap binds future development to
present strategic commitment. Source code binds future computation to present
specification---and is unique in that the binding is mechanically executable.

\begin{proposition}[Prospective Notes and Temporal Binding]
A prospective note $N$ performs a Bind operation if and only if
\[
P_A(c_{t'}, t' \mid N, A_t) > P_A(c_{t'}, t' \mid A_t),
\]
where $t' > t$ and the conditioning on $A_t$ indicates the agent's present state
is available in both cases.
\end{proposition}

\begin{proposition}[Collective Reachability from Coordination]
Let $A_1, \ldots, A_n$ be agents with independent continuation spaces
$\calC_1, \ldots, \calC_n$. A coordination note $N$ increases collective
reachability if
\[
|\calC(A_1, \ldots, A_n \mid N)| > |\calC_1| + \cdots + |\calC_n|.
\]
\end{proposition}

The proposition captures the supraadditive character of coordination: bound agents
can traverse paths that no individual could reach alone.

\chapter{Repair, Refusal, and Generative Notes}

\section{Repair Notes}

\begin{definition}[Repair Note]
A \emph{repair note} is a note whose primary function is to reduce the cost of
reconstructing a continuation after it has been or is anticipated to be
interrupted. It combines Refuse (preserving the history of the trajectory against
collapse) and Pop (enabling re-entry at a specific point after interruption).
\end{definition}

\begin{theorem}[Notes as Pre-emptive Repair]
Every note $N$ with positive repair value $R_v(N, c) = C_r(c) - C_r(c \mid N) > 0$
constitutes a pre-emptive repair: it reduces expected reconstruction cost prior
to any failure occurring.
\end{theorem}

\begin{proof}
$R_v(N, c) > 0$ implies $C_r(c \mid N) < C_r(c)$. Since the note exists before
any failure occurs, the cost reduction is in force from the moment of note
creation. The note therefore functions as a repair mechanism that precedes the
event it repairs for.
\end{proof}

A troubleshooting guide is a repair note: a stored Pop operator for a class of
failure states. Each entry says: if you arrive at this failure state, the
continuation from here to a functional state runs through these steps.

\section{Refusal Notes}

\begin{definition}[Refusal Note]
A \emph{refusal note} is a note whose primary function is to prevent the collapse
of an existing distinction, independently of whether the preserved distinction is
expected to be used. It is the closest realization of the pure Refuse operation.
\end{definition}

Mathematical proofs, cryptographic checksums, archival copies, legal records, and
formal specifications belong to this class.

\begin{proposition}[Proofs as Refusal Notes]
A mathematical proof $P$ of theorem $T$ from premises $\Gamma$ satisfies
\[
P_A(c_{\Gamma \vdash T}, t \mid P) > P_A(c_{\Gamma \vdash T}, t)
\]
for any agent $A$ with sufficient background knowledge, at any future time $t$.
\end{proposition}

The proof refuses the collapse of an inferential trajectory into mere assertion.
Its value is not actualized use but preserved potential: it maintains the
reachability of the derivation whether or not anyone consults it.

\section{Generative Notes}

\begin{definition}[Generative Note]
A \emph{generative note} is a note that creates access to continuations that were
not reachable before the note's creation. It is the primary realization of the Pop
operation: it opens new trajectories in possibility space.
\end{definition}

\begin{proposition}[Generative Notes and Novel Reachability]
A generative note $N$ satisfies
\[
\exists\, c \in \calC : P_A(c, t \mid N) > 0 \text{ and } P_A(c, t) = 0.
\]
That is, $N$ creates access to continuations that were previously unreachable.
\end{proposition}

Scientific theories, maps, programming languages, mathematical formalisms, and
ontologies are generative notes. A programming language is a generative note that
is mechanically executable: it defines not just a space of possible programs but
a machine for traversing that space.

\chapter{Collapse Notes and the Ephemerality Inversion}

\section{The Hidden Class}

\begin{definition}[Collapse Note]
A \emph{collapse note} is an artifact that performs a deliberate projection
$\sigma_N : \calH \to S_N$, compressing a history or artifact into a more compact
representation at the cost of detail, in order to make the representation
navigable. It is the realization of the Collapse operator.
\end{definition}

File names, titles, abstracts, indexes, summaries, taxonomies, classifications,
labels, and keywords are all collapse notes. The paradox of collapse notes is that
they appear to be the opposite of notes---they destroy information---yet they are
indispensable to the function of every other note class.

\begin{theorem}[Necessity of Collapse Notes]
For any sufficiently large collection of notes $\calN$, there exists a threshold
$|\calN| > K$ above which the reachability of individual notes in $\calN$ decreases
in the absence of collapse notes over $\calN$.
\end{theorem}

\begin{proof}
Search cost through an unstructured collection grows at least linearly with
collection size. For $|\calN| > K$, exhaustive search becomes too expensive for
practical use. A collapse note---index, catalogue, taxonomy---reduces search cost
by imposing structure, making individual notes reachable at sublinear cost.
Without such structure, notes become inaccessible through navigational failure,
not destruction.
\end{proof}

\begin{remark}
A word is a collapse note: it compresses a complex of experience, association,
inference, and usage into a single retrievable token. A noun is a collapse note
about a process. The passive theory would not recognize these as notes at all.
The Spherepop theory recognizes them immediately as the most pervasive and powerful
class of collapse notes civilization has developed.
\end{remark}

\section{The Ephemerality Inversion}

\begin{definition}[Continuation Completion]
A continuation $c$ is \emph{complete} with respect to note $N$ at time $t$ if
the trajectory encoded by $N$ has been successfully traversed.
\end{definition}

\begin{theorem}[Ephemerality Inversion]
Let $N$ be a note encoding a single continuation $c$. If $c$ has been completed,
then the destruction of $N$ does not reduce the reachability of any remaining
continuation. Therefore the ephemerality of $N$ after completion of $c$ implies
neither loss nor failure.
\end{theorem}

\begin{proof}
By assumption, $c$ is the sole continuation for which $N$ increases reachability.
After completion of $c$, $V_C(N) = 0$. Since continuation volume is zero, the
destruction of $N$ reduces reachability by zero.
\end{proof}

The grocery list discarded at the supermarket exit has not failed. It has succeeded
completely. The theorem generates an inversion of the passive theory's ranking:

\begin{center}
\begin{tabular}{lll}
\toprule
Fate & Passive Theory & Continuation Theory \\
\midrule
Executed and discarded & Low value & Optimal \\
Preserved and re-entered & Medium value & Valuable \\
Preserved, never re-entered & High value & Unrealized potential \\
\bottomrule
\end{tabular}
\end{center}

The celebrated survivors of antiquity---incomplete manuscripts, undelivered
letters, abandoned plans---survive precisely because their continuations were never
completed. The completed notes, the sent letters, the finished buildings, left no
trace because they succeeded.

%% ============================================================
%% PART VI: COMPILATION, GC, AND ERROR CORRECTION
%% ============================================================
\part{Compilation, Garbage Collection, and Error Correction}

\chapter{The Event IR and Compiler Correctness}
\label{ch:compiler}

\section{The Failure of Traditional Compiler Correctness}

The standard notion of compiler correctness is: a compiler is correct if the
compiled program produces the same output as the interpreted program for every
input. This criterion is adequate for state-centric languages. It is insufficient
for Spherepop.

Two programs that produce the same final value may have constructed radically
different histories to get there, and those histories are part of what the program
is.

\begin{definition}[History Equivalence as Correctness]
Let $I : P \to \calH$ be the interpreter function mapping programs to histories,
and $V : \mathsf{IR} \to \calH$ the VM mapping event-IR blocks to histories,
and $C : P \to \mathsf{IR}$ the compiler. A compiler $C$ is correct if for all
programs $p$ and initial worlds $W$:
\[
I(p, W).\mathsf{history} = V(C(p), W).\mathsf{history}.
\]
Output equivalence (same final value) is a strictly weaker criterion that is
insufficient for Spherepop.
\end{definition}

\section{Event IR}

Conventional intermediate representations encode instructions that mutate machine
state. Spherepop IR encodes \emph{events}---records of occurrences that persist
as part of the history the VM is constructing.

\begin{definition}[Event IR]
The Spherepop intermediate representation is a sequence of $\mathsf{IrEvent}$
terms:
\[
\mathsf{IR} ::= \Pop \mid \Refuse(r) \mid \Collapse(c) \mid \Bind \mid \mathsf{Load}(x) \mid \mathsf{Apply} \mid \cdots
\]
Each event carries its reason or rule as a first-class argument, preserving the
full certificate structure of the source term.
\end{definition}

\begin{lstlisting}[style=rust, caption={Event IR in Rust}]
pub enum IrEvent {
    Pop,
    Refuse { reason: RefusalReason },
    Collapse { rule: CollapseRule },
    Bind,
    Load(Symbol),
    Apply,
    LamIntro(Symbol, Type),
}
\end{lstlisting}

\begin{proposition}[IR Faithfulness]
The compilation function $C : \mathsf{Term} \to \mathsf{IR}$ is faithful: for
every primitive term $t$, $C(t)$ contains exactly the event variants that $I(t)$
would append to the history, with identical arguments.
\end{proposition}

\begin{theorem}[Replay Equivalence]
The Spherepop compiler satisfies history equivalence for all programs expressible
in the Spherepop core calculus.
\end{theorem}

\begin{proof}
By structural induction on program terms. For each primitive operation (Refuse,
Bind, Collapse, Pop), the compiler emits the corresponding $\mathsf{IrEvent}$
variant, and the VM step performs exactly the same World mutation as the
interpreter's eval case. The history appended is therefore identical. For $\mathsf{Seq}$:
each sub-term is compiled in order, and since histories compose by concatenation,
the composed history of the compiled form equals that of the interpreted form by
induction.
\end{proof}

\begin{theorem}[Compiler Full Faithfulness]
\[
I(p_1).\mathsf{history} = I(p_2).\mathsf{history}
\iff
V(C(p_1)).\mathsf{history} = V(C(p_2)).\mathsf{history}.
\]
\end{theorem}

\begin{proof}
($\Rightarrow$) By Replay Equivalence applied twice. ($\Leftarrow$) Similarly.
The biconditional follows from the symmetry of history equivalence.
\end{proof}

\begin{corollary}[Operational-State Faithfulness, Repaired]
\label{cor:op-state-faithfulness}
Under the same hypothesis,
\[
I(p_1).\mathsf{history} = I(p_2).\mathsf{history}
\;\Longrightarrow\;
F_{\mathrm{op}}(I(p_1).\mathsf{history}; M_0) = F_{\mathrm{op}}(I(p_2).\mathsf{history}; M_0)
\]
and likewise for the compiled/VM histories, so equal histories under
Compiler Full Faithfulness entail equal operational states $M_t$
(Section~\ref{sec:layered-world-model}) on both the interpreter and VM sides.
This closes the gap the implementation chapter's replay test targets: the
assertion \texttt{world.materialized == fold\_operational(\&world.history)}
is exactly the case $p_1 = p_2$ of this corollary applied to a single program
compared against its own operational-state cache.
\end{corollary}

\begin{proof}
$F_{\mathrm{op}}(-; M_0)$ is a function of $H$ alone (Definition of Operational
State, Section~\ref{sec:layered-world-model}); applying an equal function to
equal arguments gives equal results.
\end{proof}

\begin{remark}[History equality remains strictly stronger]
The converse of Corollary~\ref{cor:op-state-faithfulness} fails in general:
equal operational states do not entail equal histories. Two different,
non-replay-equivalent event sequences can fold to the same operational-state
cache---for instance, two orderings of Refuse events for distinct symbols
produce the same admissibility index $\calA(H)$ even though the histories
differ as sequences, since $\calA$ depends only on the \emph{set} of refused
symbols (Definition of Admissible Futures, Section~\ref{sec:admissibility-set}),
not their order. Compiler correctness should therefore compare \emph{both}
criteria---history equality via Compiler Full Faithfulness, and operational-state
equality via Corollary~\ref{cor:op-state-faithfulness}---with history equality
understood as the strictly stronger of the two, exactly as the implementation
chapter's revision requires.
\end{remark}

\chapter{Garbage Collection as Admissibility-Guided Collapse}

\section{The Standard Problem}

In conventional languages, garbage collection identifies and reclaims memory that
is no longer reachable from the program's roots. The criterion is purely
structural: an object is garbage if no live reference points to it.

In Spherepop, the analogous question is: which portions of the history are no
longer needed to support any reachable continuation? The answer is richer because
it depends not just on structural reachability but on admissibility.

\section{History Segments and Liveness}

\begin{definition}[History Segment Liveness]
A history segment $H[i:j] = [e_i, e_{i+1}, \ldots, e_j]$ is \emph{live at
time $t$} if there exists an agent $A$ and a continuation $c$ such that:
\begin{enumerate}
\item $c$ is reachable at time $t$: $P_A(c, t) > 0$, and
\item $c$ depends on $H[i:j]$: $P_A(c, t \mid \neg H[i:j]) < P_A(c, t)$.
\end{enumerate}
A segment is \emph{dead} if it is not live.
\end{definition}

\begin{definition}[Admissibility-Guided GC]
\emph{Admissibility-guided garbage collection} removes dead history segments
while preserving the admissibility certificates needed for future collapses.
Formally, let $\calA_{\mathrm{live}}(H)$ be the set of future continuations
still reachable. GC produces $H' \subseteq H$ (as a subsequence) satisfying:
\[
\calA(H') = \calA(H) \cap \{x : \exists c \in \calA_{\mathrm{live}}(H),\, x \text{ needed by } c\}.
\]
\end{definition}

\section{GC as Collapse}

\begin{theorem}[GC as Collapse Composition]
Admissibility-guided GC is equivalent to applying a collapse rule $c_{\mathrm{GC}}$
that maps $H$ to the subsequence of live events:
\[
c_{\mathrm{GC}}(H) = [e \in H : e \text{ is live in } H].
\]
$c_{\mathrm{GC}}$ is a valid collapse rule when liveness is stable under
extension: if $e$ is live in $H$, it is live in $H \mathbin{++} H'$ for all $H'$.
\end{theorem}

\begin{proof}
Liveness is defined in terms of reachable continuations. Since histories grow
monotonically and Pop operations only add to the committed history, a past
commitment that enables a reachable continuation remains enabling under history
extension. Therefore $c_{\mathrm{GC}}$ respects composition: the GC of a
concatenated history equals the concatenation of the GC of the parts.
\end{proof}

\begin{corollary}[GC Preserves Admissibility]
If $c_{\mathrm{GC}}(H) = H'$, then $\calA(H') \supseteq \calA_{\mathrm{live}}(H)$.
\end{corollary}

\begin{proof}
GC removes only dead events. A dead Refuse event is one that marks a symbol
$x$ as inadmissible when $x$ is no longer needed by any reachable continuation.
Removing it restores $x$ to the admissibility set. Since no live continuation
needs $x$, this does not affect any future collapse certificate.
\end{proof}

\section{Three GC Strategies}

Different collapse rules for GC yield different strategies with different
cost-benefit profiles.

\begin{definition}[Mark-Scan GC]
\emph{Mark-Scan GC} applies the identity rule $c_I$ to the subhistory of live
events: it retains the full sequence of live events in order. This is the most
informative strategy and preserves all historical structure needed for replay.
\end{definition}

\begin{definition}[Quotient GC]
\emph{Quotient GC} applies a coarser rule $c_{LW}$ to the history before
collecting: it retains only the last commitment to each symbol. This is equivalent
to state-based GC with a history log and loses intermediate state information.
\end{definition}

\begin{definition}[Certificate-Only GC]
\emph{Certificate-Only GC} retains only Refuse events (inadmissibility
certificates) and Collapse events (observation certificates), discarding all Pop
and Bind events whose certificates have been sealed. This minimizes memory use
while preserving the admissibility structure needed for future type checking.
\end{definition}

\begin{theorem}[GC Strategy Tradeoff]
Let $M(c_{\mathrm{GC}})$ denote memory usage and $R(c_{\mathrm{GC}})$ denote
recovery power (the set of continuations that can be reconstructed after GC).
Then:
\[
M(c_I) \geq M(c_{LW}) \geq M(c_{\mathrm{cert}})
\quad \text{and} \quad
R(c_I) \geq R(c_{LW}) \geq R(c_{\mathrm{cert}}).
\]
\end{theorem}

\begin{proof}
The identity rule retains the most events and therefore uses the most memory but
also preserves the most recovery information. Each coarsening reduces both memory
usage and recovery power by identifying more history segments as equivalent.
\end{proof}

\section{Admissibility Certificates and GC Safety}

A key constraint on GC in Spherepop is that it must not invalidate pending type
judgements. If a collapse term $\mathsf{collapse}(t, c)$ is waiting on the
admissibility certificate $\Gamma \vdash t : \Adm(T)$, GC must retain the history
events that support that certificate.

\begin{definition}[GC Safety]
A GC operation is \emph{safe} if for every pending collapse judgement
$\Gamma \vdash \mathsf{collapse}(t, c) : \Col(T, c)$, the history events needed
to support $\Gamma \vdash t : \Adm(T)$ are retained after GC.
\end{definition}

\begin{theorem}[Safety of Certificate-Only GC]
Certificate-Only GC is safe: it retains all Refuse events that contribute to
any pending admissibility certificate.
\end{theorem}

\begin{proof}
The admissibility type $\Adm(T)$ is established by the sequence of Pop and Refuse
events that produced $t$. Certificate-Only GC retains all Refuse events
(inadmissibility certificates) unconditionally. The Pop events can be inferred
from the Collapse certificates that seal their results. Therefore all information
needed for pending collapse judgements is preserved.
\end{proof}

\chapter{Error Correction as Certified Refusal}

\section{The Standard Approach and Its Limits}

In conventional languages, error handling is an afterthought: exceptions are
thrown when something goes wrong, and error handling code is written separately
from the main computation. This approach has three problems. First, it treats
errors as exceptional when they are often structurally inevitable. Second, it
provides no information about why the error occurred. Third, it makes error
recovery expensive because the history of the computation is not available.

In Spherepop, errors are not exceptional. Every Refuse event is a documented
non-traversal of a possible path. Error handling is not separate from the main
computation; it is part of the computation's history.

\section{Error Correction Codes in History Space}

\begin{definition}[History Redundancy]
A history $H$ has \emph{redundancy factor} $k$ if there exist $k$ disjoint
subsequences $H_1, \ldots, H_k \subseteq H$ (as event subsequences) such that
each $H_i$ alone is sufficient to reconstruct the full admissibility certificate
for the terminal collapse:
\[
\calA(H_i) \supseteq \calA_{\mathrm{needed}}(H) \quad \text{for each } i.
\]
\end{definition}

\begin{definition}[Erasure-Correcting History]
An \emph{erasure-correcting history} with parameters $(n, k)$ encodes $k$
independent certificate paths through $n$ events, such that any $k$ of the $n$
events are sufficient to recover the full admissibility structure.
\end{definition}

\begin{theorem}[Error Correction via Redundant Refusal]
If a history $H$ with redundancy factor $k \geq 2$ has $m < k/2$ events
corrupted or erased, the correct admissibility structure can be recovered from
the remaining events.
\end{theorem}

\begin{proof}
With redundancy $k$ and fewer than $k/2$ erasures, at least $k - k/2 > k/2$
redundant paths remain intact. Majority vote over the surviving paths recovers
the correct admissibility judgement.
\end{proof}

\section{Typed Error Propagation}

\begin{definition}[Error Type]
An \emph{error type} is a type of the form $\Rfsd(T, r)$ where $r$ is a
proof-carrying reason. Error types carry the evidence of inadmissibility as
first-class type-level information.
\end{definition}

\begin{definition}[Error Propagation Rule]
\[
\frac{\Gamma \vdash t : \Rfsd(T_1, r) \quad \Gamma, x : T_1 \vdash b : T_2}{\Gamma \vdash \mathsf{handle}(t, x.b) : T_2 \vee \Rfsd(T_2, r)}
\quad [\mathrm{T\text{-}HANDLE}]
\]
The handler $b$ may produce a value of $T_2$ or propagate a refusal. The
propagated refusal carries the original reason $r$, preserving the full provenance
chain.
\end{definition}

\begin{theorem}[Error Preservation]
If $\Gamma \vdash t : \Rfsd(T, r)$ and there is no handler in scope for $r$,
then $r$ propagates to the nearest enclosing context that can handle
$\Rfsd(\cdot, r)$.
\end{theorem}

\begin{proof}
By the typing rules, $\Rfsd(T, r)$ types only propagate upward through
$[\mathrm{T\text{-}HANDLE}]$ if the handler does not fully resolve them. Since
every unresolved $\Rfsd(T, r)$ in a well-typed program must eventually reach a
handler or become the return type of the program, propagation terminates at the
program boundary with a documented reason.
\end{proof}

\section{Recovery as Certified Continuation}

\begin{definition}[Recovery Certificate]
A \emph{recovery certificate} for an error $\Rfsd(T, r)$ at world state $W$ is a
term $\mathsf{recover} : \Rfsd(T, r) \to \Adm(T')$ that transforms a documented
inadmissibility into an admissible continuation of type $T'$.
\end{definition}

\begin{theorem}[Soundness of Recovery]
If $\mathsf{recover} : \Rfsd(T, r) \to \Adm(T')$ is a recovery certificate and
$\Gamma \vdash t : \Rfsd(T, r)$, then $\Gamma \vdash \mathsf{recover}(t) : \Adm(T')$.
Furthermore, the history of $\mathsf{recover}(t)$ contains the full provenance
chain from the original refusal to the recovery, ensuring that the recovery is
not a silent erasure of the error.
\end{theorem}

\begin{proof}
By application of the recovery function type and $[\mathrm{T\text{-}POP}]$ applied
to the result. The history contains the original $\Refuse(x, r)$ event followed
by the recovery events. No history erasure is possible since histories grow
monotonically.
\end{proof}

\begin{remark}
Recovery in Spherepop is the computational analogue of repair theory: it restores
access to a continuation after documented inadmissibility, but unlike silent error
suppression, it preserves the full record of what went wrong and how it was
addressed. The recovery certificate is a note---specifically a Repair Note---in
the formal sense of the note taxonomy.
\end{remark}

%% ============================================================
%% PART VII: DEEP STRUCTURE
%% ============================================================
\part{Deep Mathematical Structure}

\chapter{The Observation Limit}

\section{The Reviewer's Challenge}

Spherepop does not abolish the state illusion. It makes the quotient explicit,
parameterised, and reversible whenever the underlying history is retained.

A careful reader will notice a tension. The document claims to resolve the state
illusion---the loss of historical information caused by projecting history to
state. Yet every output of a Spherepop program is a collapse: a value derived by
applying a rule to a history. The programmer observes the collapsed value, not
the history. Is Spherepop not simply relocating the state illusion?

This chapter addresses that challenge. We prove the Finite Observation
Non-injectivity Theorem, which shows the challenge follows from the specified
cardinalities rather than from the nature of observation as such, and then
distinguish precisely between state erasure (which Spherepop eliminates) and
state illusion (which bounded observation cannot eliminate).
\label{ch:no-direct-observation}

\section{The Finite Observation Non-injectivity Theorem}

\begin{definition}[History-Faithful Observation]
An observation map $\phi : \calH \to O$ is \emph{history-faithful} if it is
injective: $\phi(H_1) = \phi(H_2) \Rightarrow H_1 = H_2$.
\end{definition}

\begin{theorem}[Finite Observation Non-injectivity]
\label{thm:no-direct-observation}
Let $O$ be any set with $|O| < |\calH|$. Then every observation map
$\phi : \calH \to O$ is non-injective: there exist distinct $H_1 \neq H_2$ with
$\phi(H_1) = \phi(H_2)$. Consequently, no computational system whose observable
space has strictly smaller cardinality than its history domain can directly
verify the history it constructed from that observable alone.
\end{theorem}

\begin{proof}
Since $\calH = \bigcup_{n \geq 0} E^n$ is countably infinite and $|O| < |\calH|$,
by the pigeonhole principle $\phi$ cannot be injective. Therefore there exist
$H_1 \neq H_2$ with $\phi(H_1) = \phi(H_2)$.
\end{proof}

\begin{remark}
The qualified conclusion, and no stronger one, follows: every observation whose
codomain has lower cardinality than its history domain is non-injective, and
every non-injective observation induces a nontrivial quotient. This does
\emph{not} license the universal claim that ``every observation is a
quotient''---the Identity observation $c_I : \calH \to \calH$, introduced
elsewhere in this book, is injective by construction and induces the trivial
quotient. Theorem~\ref{thm:no-direct-observation} is a cardinality constraint on
bounded observables, not a claim about observation as such.
\end{remark}

\begin{remark}[Observational loss is not custodial loss]
This chapter's result concerns \emph{observational} loss: a lossy view need not
erase anything if $H$ remains available elsewhere. State illusion, in the sense
of Definition~\ref{def:state-illusion} below, is inevitable under bounded views;
state erasure is a separate, architectural choice about whether $H$ itself is
retained. Conflating the two would make the coming results about custody and
provenance (developed further in the discussion of physical inscription and
recoverability) look like restatements of a cardinality theorem, when they are
in fact claims about a different failure mode entirely.
\end{remark}

\section{Erasure versus Illusion}

\begin{definition}[State Erasure]
A computational system suffers \emph{state erasure} if it irrecoverably discards
the history $H$ after computing $\sigma(H)$. After execution, $H$ is unavailable
and only $\sigma(H)$ remains.
\end{definition}

\begin{definition}[State Illusion]
\label{def:state-illusion}
A computational system presents a \emph{state illusion} to its programmers if
they can only observe $\sigma(H)$ and not $H$ directly. By Theorem 7.1, some
state illusion is unavoidable for any non-trivial observable space.
\end{definition}

\begin{proposition}[What Spherepop Eliminates]
Spherepop eliminates state erasure: the history $H$ is retained in
\texttt{world.history} and is available to any observer. Spherepop cannot and
does not eliminate the state illusion: the programmer's primary observable is
$c(H)$, not $H$, so distinct histories that agree under $c$ remain
indistinguishable from the perspective of that observation.

The correct claim is therefore: Spherepop does not abolish the state illusion.
It makes the quotient explicit, parameterised, and history-preserving.
\end{proposition}

\section{Rule Refinement as Disambiguation}

\begin{definition}[Rule Refinement Order]
\label{def:refinement-order}
Collapse rule $c_1$ is \emph{finer than} $c_2$ (equivalently $c_2$ is
\emph{coarser than} $c_1$), written $c_1 \preceq c_2$, if $c_2$ factors through
$c_1$:
\[
c_1 \preceq c_2 \quad\Longleftrightarrow\quad \exists f \;\; c_2 = f \circ c_1.
\]
Factorization is used here rather than the equivalent but more error-prone
definition via reversed inclusion of induced equivalence relations, since the
latter is easy to state backward. The two are consistent: $c_1 \preceq c_2$ in
the sense above implies $H_1 \sim_{c_1} H_2 \Rightarrow H_1 \sim_{c_2} H_2$,
because $c_1(H_1) = c_1(H_2)$ gives $c_2(H_1) = f(c_1(H_1)) = f(c_1(H_2)) =
c_2(H_2)$. The Identity rule $c_I$ is the finest rule (every $c$ factors as
$c = c \circ c_I$); the trivial rule $c_\top$ mapping all histories to a single
point is the coarsest.
\end{definition}

\begin{proposition}[Disambiguation by Refinement]
Let $H_1 \sim_c H_2$ (histories indistinguishable under rule $c$) and $H_1 \neq H_2$.
Then there exists a finer rule $c' \preceq c$ such that $c'(H_1) \neq c'(H_2)$.
In the limit, $c_I$ always disambiguates since $c_I(H) = H$ exactly.
\end{proposition}

\begin{proof}
Take $c' = c_I$: since $H_1 \neq H_2$ and $c_I(H) = H$, we have
$c_I(H_1) = H_1 \neq H_2 = c_I(H_2)$.
\end{proof}

\chapter{Category Theory as Engine}

\begin{remark}[Rechecking after the VIEW/Collapse separation]
Every categorical claim in this chapter concerning ``observation'' was written
before Section~\ref{sec:view-vs-collapse} separated pure VIEW from committed
Collapse. The functors $\mathcal{O}_c$ and the adjunction below concern
projection or quotienting---VIEW, in the vocabulary of that section---not the
event that records the adoption of a quotient as constraining. Where a
left-adjoint or reflection is claimed, it is a reflection onto the observable
category $\Obs_c$ via the pure map $c$; it says nothing, by itself, about when
an observation is attested or collapsed into history. This chapter should be
read with that reattribution in mind throughout.
\end{remark}

\section{The History-Observation Adjunction}

\begin{remark}[A naming collision this repair must resolve first]
Before restating the adjunction, one more inconsistency has to be named. The
symbol $\Hist$ has been used for two incompatible categories in this book: the
\emph{History Category} of Chapter~\ref{ch:collapse-quotient} has option spaces
$\Omega \subseteq \Omega_0$ as objects and histories as morphisms between them,
whereas the History Construction Functor below treats $G(t)$ as a single
history, which only typechecks if the objects of the target category
\emph{are} histories. These are not the same category, and no functor can
target both at once. The repair introduces a second, explicitly named
category for histories-as-objects, rather than continuing to overload $\Hist$.
\end{remark}

\begin{definition}[Prefix Category of Histories]
$\Hist_\sqsubseteq$ is the category whose objects are histories $H \in \calH$
and whose morphisms are given by the prefix order: there is a (unique) morphism
$H \to H'$ iff $H$ is a prefix of $H'$, i.e. $H' = H \mathbin{++} K$ for some
$K \in \calH$. Composition is transitivity of the prefix relation; the identity
on $H$ is the trivial extension by $\varepsilon$. Since the prefix relation on
finite sequences is a genuine partial order (antisymmetric: $H \sqsubseteq H'$
and $H' \sqsubseteq H$ force $H = H'$), $\Hist_\sqsubseteq$ is a thin category
(a poset viewed as a category), distinct from the morphisms-are-histories
category $\Hist$ used elsewhere.
\end{definition}

\begin{definition}[History Construction Functor, Repaired]
Let $\mathbf{Terms}$ be the category whose objects are Spherepop terms and whose
morphisms are reduction sequences. Define $G : \mathbf{Terms} \to \Hist_\sqsubseteq$
by $G(t) = I(t).\mathsf{history}$ on objects. On morphisms, a reduction sequence
from $t$ to $t'$ is sent to the prefix morphism $G(t) \to G(t')$; this is
well-defined because interpretation only appends events along a reduction, so
$I(t).\mathsf{history}$ is always a prefix of $I(t').\mathsf{history})$ when $t$
reduces to $t'$.
\end{definition}

For the observation side, functoriality of $c$ with respect to prefix order is
not automatic---it is exactly as delicate as the composition-respecting
property already required for the Collapse Functor (Theorem~1.101), and for the
same reason: $c_{LW}$ is sensitive to what comes later in a history, so
extending a history can move its LastWrite value in either direction, not
monotonically.

\begin{definition}[Prefix-Monotone Collapse Rule]
A collapse rule $c : \calH \to O_c$ is \emph{prefix-monotone} if there is a
partial order $\preceq_c$ on $O_c$ such that $H \sqsubseteq H'$ implies
$c(H) \preceq_c c(H')$.
\end{definition}

\begin{example}
The Identity rule $c_I$ is prefix-monotone with $\preceq_{c_I} \;=\; \sqsubseteq$
itself. The Accumulate rule $c_{\mathrm{acc}}$, which counts occurrences of each
symbol, is prefix-monotone under the pointwise order on count vectors: appending
events can only increase counts. The LastWrite rule $c_{LW}$ is \emph{not}
prefix-monotone under any order, for the same reason it fails to respect
concatenation in Theorem~1.101: appending a write to $x$ can move $c_{LW}(H)$ to
an unrelated value, not one comparable to the old value under any fixed order.
\end{example}

\begin{definition}[Observation Functor, Repaired]
For a prefix-monotone collapse rule $c$, define
$\mathcal{O}_c : \Hist_\sqsubseteq \to (O_c, \preceq_c)$ by
$\mathcal{O}_c(H) = c(H)$. This is a functor between thin categories (i.e., a
monotone map of posets) precisely by the definition of prefix-monotonicity.
\end{definition}

The earlier ``History-Observation Adjunction,'' stated only for $c_I$, is now
visibly degenerate once this machinery is in place: since $c_I$ is the identity
map, $\mathcal{O}_{c_I}$ is the identity functor on $\Hist_\sqsubseteq$, and the
claimed bijection $\Hist(G(t), H) \cong \mathbf{Terms}(t, G^{-1}(H))$ reduces to
the tautology that $G(t) = H$ iff $t$ reduces to some term with that exact
history---which is simply Replay Equivalence restated, not a nontrivial
adjunction. Calling it an adjunction overstated what the identity-rule case
actually shows.

\begin{theorem}[Reflection for Prefix-Monotone Rules]
\label{thm:reflection-repaired}
Let $c$ be a prefix-monotone collapse rule equipped with a section
$s_c : O_c \to \calH$ satisfying $c(s_c(o)) = o$ for all $o$ (a chosen canonical
representative history for each observable value). Suppose further that $s_c$
is \emph{order-reflecting}: $H \sqsubseteq s_c(o) \iff c(H) \preceq_c o$ for all
$H \in \calH$ and $o \in O_c$. Then $\mathcal{O}_c \dashv s_c$: viewing $s_c$ as
a functor $(O_c, \preceq_c) \to \Hist_\sqsubseteq$, it is right adjoint to
$\mathcal{O}_c$, and $s_c$ exhibits $(O_c, \preceq_c)$ as a reflective
subcategory of $\Hist_\sqsubseteq$.
\end{theorem}

\begin{proof}
For thin categories, an adjunction is exactly a Galois connection: monotone
maps $F : P \to Q$ and $G : Q \to P$ between posets satisfy $F \dashv G$ iff
$F(p) \leq_Q q \iff p \leq_P G(q)$ for all $p, q$. Here $P = \calH$ (ordered by
$\sqsubseteq$), $Q = O_c$ (ordered by $\preceq_c$), $F = \mathcal{O}_c$, and
$G = s_c$. The hypothesis that $s_c$ is order-reflecting is precisely
$H \sqsubseteq s_c(o) \iff c(H) \preceq_c o$, which is the Galois-connection
condition read in one direction; the other direction,
$\mathcal{O}_c(H) \preceq_c o \iff H \sqsubseteq s_c(o)$, is the same statement.
Hence $\mathcal{O}_c \dashv s_c$.
\end{proof}

\begin{example}[The Accumulate rule satisfies the hypotheses]
Let $s_{\mathrm{acc}}$ send a count vector $o$ to the (essentially unique, up to
event ordering within a symbol) history obtained by popping each symbol the
number of times its count in $o$ specifies, in some fixed canonical order (say,
symbol index order). Then $c_{\mathrm{acc}}(s_{\mathrm{acc}}(o)) = o$ by
construction, and $H \sqsubseteq s_{\mathrm{acc}}(o)$ iff every count in $H$ is
$\leq$ the corresponding count in $o$ (since $s_{\mathrm{acc}}(o)$ realizes
exactly those counts and prefixes of it realize sub-counts in the canonical
order), which is exactly $c_{\mathrm{acc}}(H) \preceq_{\mathrm{acc}} o$. So
Theorem~\ref{thm:reflection-repaired} applies to $c_{\mathrm{acc}}$, giving a
genuine, nontrivial reflection: $(O_{\mathrm{acc}}, \preceq_{\mathrm{acc}})$ is
a reflective subcategory of $\Hist_\sqsubseteq$, with $\mathcal{O}_{\mathrm{acc}}$
as the reflector.
\end{example}

\begin{remark}[What this repairs and what it does not]
The theorem is stated for prefix-monotone rules with an order-reflecting
section---hypotheses satisfied by $c_I$ (trivially) and $c_{\mathrm{acc}}$
(shown above), but not by $c_{LW}$, matching its established failure to respect
composition in Theorem~1.101. The original ``Collapse as Reflection'' claim,
below in the derivations chapter, is restated to point back here rather than
re-deriving a weaker or inconsistent version. As with the Functoriality
theorem, this result concerns the pure functor $\mathcal{O}_c$---VIEW, in the
sense of Section~\ref{sec:view-vs-collapse}---and says nothing about when an
observation is attested or committed as a Collapse event; the reflective
subcategory is a fact about the geometry of quotienting, not about which
observations a program has chosen to adopt as binding.
\end{remark}

\section{Collapse Rules Form a Lattice}

\begin{proposition}[Observational Lattice]
The set of collapse rules ordered by refinement $c_1 \preceq c_2$ forms a complete
lattice. The meet of $\{c_i\}$ is the finest rule at least as fine as all; the
join is the rule generated by the union of equivalence relations.
\end{proposition}

\begin{proof}
The meet and join are defined by the standard lattice operations on equivalence
relations. Completeness follows because arbitrary intersections and generated
equivalence relations of equivalence relations are again equivalence relations.
\end{proof}

\section{Sheaf-Theoretic Semantics}

\begin{definition}[Observational Site]
Fix collapse rules $\{c_i\}$ and observable states $\{o_i\}$. The
\emph{observational site} $\mathcal{O}$ is the category whose objects are
observational regions $\{c_i^{-1}(o_i)\}$ and whose morphisms are inclusions
$c_j^{-1}(o_j) \hookrightarrow c_i^{-1}(o_i)$ when the former is contained in
the latter.
\end{definition}

\begin{proposition}[$\calH$ is a Sheaf]
The history presheaf $\calH$ on $\mathcal{O}$---defined by $\calH(U) = \{H \in \calH : H \in U\}$---satisfies the sheaf condition: compatible local histories glue
uniquely to global histories.
\end{proposition}

\begin{proof}
Compatibility means event sequences agree on their shared prefixes. Two compatible
event sequences agreeing on all overlaps determine a unique concatenation.
\end{proof}

\begin{proposition}[State Illusion as Non-Injectivity of Sections]
A collapse rule $c$ exhibits state illusion if and only if the observable sheaf
$\mathcal{O}_c$ fails to determine a unique global section of $\calH$. There
exist distinct $H_1, H_2 \in \calH(\calH)$ with the same section in
$\mathcal{O}_c(\calH)$.
\end{proposition}

The sheaf structure ties the local (what one observer sees under one rule) to the
global (the full history that all rules project from). Sheaf cohomology measures
the obstruction to lifting local observations to global histories---which is
precisely the observational entropy.

\chapter{CLIO Projections and Active Geodesic Inference}

\section{Collapse Rules as CLIO Projections}

\begin{definition}[CLIO Projection]
A CLIO projection $\pi : X \to M$ maps a representation space $X$ to a manifold
$M$ of observational strata. In the Spherepop setting: $X = \calH$ (history
space), $M = O_c$ (observable state space under rule $c$), and $\pi = F_c$
(the collapse functor).
\end{definition}

\begin{definition}[Observational Entropy]
For a collapse rule $c$ and observable state $o \in O_c$, the
\emph{observational entropy} is
\[
S_c(o) = \log |c^{-1}(o)|.
\]
\end{definition}

\begin{theorem}[Entropy Monotonicity Under Rule Refinement]
If $c_1 \preceq c_2$ ($c_1$ finer), then for every $o \in O_{c_1}$:
\[
S_{c_1}(o) \leq S_{c_2}(c_2(c_1^{-1}(o))).
\]
Finer rules have smaller or equal fibre sizes.
\end{theorem}

\begin{proof}
Since $c_1 \preceq c_2$, every $\sim_{c_1}$-class is contained in a
$\sim_{c_2}$-class. Therefore $|c_1^{-1}(o)| \leq |c_2^{-1}(c_2(H))|$ for
any $H \in c_1^{-1}(o)$.
\end{proof}

\begin{remark}
Observational entropy $S_c(o) = \log |c^{-1}(o)|$ is exactly the CLIO quantity
$S_\pi(m) = \log \mathrm{Vol}(\pi^{-1}(m))$ in the discrete Spherepop setting.
Spherepop is a discrete computational instance of the CLIO projection program.
\end{remark}

\section{The Santoro--Waghmare--Panaretos Unification}

\begin{theorem}[Representation Changes Topology]
Let $\rho_1, \rho_2 : X \to Y$ be two representation maps with $\tau_{\rho_1} \neq \tau_{\rho_2}$. Then there exist $x_1, x_2 \in X$ such that $x_1 \sim_{\rho_1} x_2$
but $x_1 \not\sim_{\rho_2} x_2$, or vice versa. The two representations disagree
on whether $x_1$ and $x_2$ are the same object.
\end{theorem}

\begin{proof}
If $\tau_{\rho_1} \neq \tau_{\rho_2}$, there exists an open set
$U \in \tau_{\rho_1} \setminus \tau_{\rho_2}$. Then $U = \rho_1^{-1}(V)$ for
some $V \subseteq Y$, but $U \neq \rho_2^{-1}(W)$ for any $W$. Therefore there
exist $x_1, x_2$ with $\rho_1(x_1) = \rho_1(x_2)$ but $\rho_2(x_1) \neq \rho_2(x_2)$.
\end{proof}

\begin{remark}
This theorem is the general form of the observation that representations
reorganise topology. It applies uniformly to:
(1) the Santoro--Waghmare--Panaretos kernel embedding, which converts metrically
close distributions into topologically singular Gaussian measures; and
(2) Spherepop collapse rules, where refinement moves history pairs between
equivalence classes.

In the kernel case, the topological change is forced by the Feldman--H\'ajek
theorem. In the collapse case, the topological change is chosen: the programmer
selects the rule, thereby choosing which distinctions are topologically visible.
Both are instances of the same structural operation.
\end{remark}

\section{Active Geodesic Inference}

The Spherepop framework connects to a broader theory of intelligent systems as
trajectory-maintaining structures.

\begin{definition}[Active Geodesic Inference]
An intelligent system performs \emph{active geodesic inference} if it actively
shapes its configuration space through energetic and entropic constraints, forming
low-action trajectories through possibility space rather than following
predetermined state transitions.
\end{definition}

\begin{proposition}[Spherepop as Active Geodesic Substrate]
Spherepop's irreversible, scope-based semantics enforce history-sensitivity and
prevent incoherent superposition by construction. Pop operations commit to
specific trajectories; Refuse operations document inadmissible branches; the
history preserves the full record of path selection. Spherepop is therefore a
natural substrate for implementing active geodesic inference: the history is the
geodesic; the admissibility set is the constraint field; collapse is the
observation that seals a trajectory as complete.
\end{proposition}

\begin{definition}[Formal Intelligence]
Intelligence is the capacity of a system to remain within a family of dynamically
admissible, low-action histories by actively reshaping its configuration space's
geometry. This definition extends across scales, encompassing learning within a
lifetime, evolution across generations, and reasoning within an episode.
\end{definition}

\begin{theorem}[Intelligence as Note Infrastructure]
An intelligent system with formal intelligence in the above sense is necessarily
a note-producing system: it constructs and maintains Capture, Prospective, and
Repair notes as part of the mechanism by which it reshapes its configuration space.
Notes are the mechanism by which intelligent systems install future reachability
from present positions.
\end{theorem}

\begin{proof}
A system that maintains low-action histories across time must preserve information
about its past trajectories (Capture Notes), commit to future trajectories
(Prospective Notes), and recover from trajectory interruptions (Repair Notes).
Each of these is a note class in the formal taxonomy. Therefore any formally
intelligent system produces notes as a structural necessity, not a contingent
feature.
\end{proof}

%% ============================================================
%% PART VIII: CIVILIZATION
%% ============================================================
\part{Civilization as Note Infrastructure}

\chapter{MEM\textbar8 and the Cognitive Substrate}

\section{Memory as Event Structure}

The MEM\textbar8 architecture treats memory not as isolated symbol storage but as
structured pattern access where retrieval depends on preserving relationships,
locality, and accessibility. Memory is not a database; it is a cognitive substrate
whose operations are structurally similar to Spherepop's history operations.

\begin{proposition}[MEM\textbar8 as Spherepop Instance]
The MEM\textbar8 architecture is a biological instance of the Spherepop framework:
\begin{center}
\begin{tabular}{ll}
\toprule
MEM\textbar8 Layer & Spherepop Note Class \\
\midrule
Episodic layer & Capture notes \\
Procedural layer & Repair notes \\
Semantic layer & Coordination notes \\
Indexical layer & Collapse notes \\
Theoretical layer & Generative notes \\
\bottomrule
\end{tabular}
\end{center}
\end{proposition}

\begin{remark}
A notebook is a miniature MEM\textbar8. A Git repository is a miniature
MEM\textbar8. An archive is a miniature MEM\textbar8. Each stores structured
trajectories organized to make reconstruction possible. The difference between
a good archive and a poor one is not quantity of material stored but quality
of access structure: how well the collapse notes and coordination notes that
organize the collection preserve the reachability of trajectories within it.
\end{remark}

\chapter{Repair Theory and Civilizational Collapse}

\section{Repair as Pre-emptive Refusal}

Repair Theory holds that repair is the restoration of access to a trajectory
that has been interrupted. The connection to Spherepop is not metaphorical but
structural: repair is the execution of a Repair Note that was created in
anticipation of the failure.

\begin{theorem}[Note Taxonomy as Failure Taxonomy]
The seven note classes correspond bijectively to seven classes of anticipated
trajectory failure:
\begin{center}
\begin{tabular}{ll}
\toprule
Note Class & Anticipated Failure \\
\midrule
Capture & Experiential collapse \\
Prospective & Intention-execution gap \\
Repair & Procedural interruption \\
Coordination & Coordination failure \\
Refusal & Distinction collapse \\
Generative & Inferential poverty \\
Collapse & Navigational failure \\
\bottomrule
\end{tabular}
\end{center}
\end{theorem}

\section{Civilizational Scale}

\begin{definition}[Civilizational Note Infrastructure]
Let civilization $Z$ possess note infrastructure $\calN = \{N_1, \ldots, N_k\}$.
The \emph{collective reachability} is
\[
R(Z) = \bigcup_{i=1}^{k} \calC(N_i),
\]
where $\calC(N_i)$ is the set of continuations preserved or enabled by $N_i$.
\end{definition}

\begin{theorem}[Civilizational Collapse]
The destruction of note infrastructure $\calN$ decreases collective reachability
even when underlying truths remain unchanged:
\[
R(Z \setminus \calN) < R(Z).
\]
A civilization that loses its note infrastructure does not become ignorant. It
becomes unable to reach its own knowledge.
\end{theorem}

\begin{proof}
By definition, $R(Z)$ includes all continuations preserved by elements of $\calN$.
If $\calN$ is destroyed, no element contributes to collective reachability. Since
some continuations are accessible only through $\calN$, it follows that
$R(Z \setminus \calN) < R(Z)$. The underlying facts are not destroyed, but the
paths that made them accessible are.
\end{proof}

\begin{center}
\begin{tabular}{lll}
\toprule
Technology & Reachability Horizon & Dominant Note Classes \\
\midrule
Speech & Minutes to hours & Coordination, Capture \\
Story & Decades & Capture, Generative \\
Writing & Centuries & All classes \\
Printing & Populations & Coordination, Refusal \\
Digital storage & Replication & All classes, at scale \\
Version control & Modification history & Repair, Capture \\
Network & Distribution & Coordination, Generative \\
\bottomrule
\end{tabular}
\end{center}

Each note technology expands the reachability horizon by addressing failure modes
the preceding technology could not prevent. Writing prevents the death of
witnesses. Printing prevents the destruction of individual copies. Version control
prevents the loss of developmental history through modification. The history of
civilization is the history of expanding note infrastructure.

%% ============================================================
%% PART IX: OPEN PROBLEMS
%% ============================================================

%% ============================================================
%% PART IX: RECONSTRUCTION, REFUSAL, AND REPAIR
%% ============================================================
\part{Reconstruction, Refusal, and Repair}

\chapter{Observational Entropy and Reconstruction Cost}

\section{From Observation to Reconstruction}

The preceding chapters established that every collapse rule
$c : \calH \to O_c$ induces an observational entropy
$S_c(o) = \log |c^{-1}(o)|$.
This quantity measures the size of the fibre associated with an
observation. A collapse with small observational entropy preserves many
distinctions between histories. A collapse with large observational entropy
identifies many histories and destroys information. The natural next question is
not merely how many histories are hidden, but how difficult it is to reconstruct
the correct one.

\begin{definition}[Reconstruction Cost]
Let $o \in O_c$. The \emph{reconstruction cost} of $o$ is
\[
\calR(o) = \mathbb{E}\bigl[C(H) \mid H \in c^{-1}(o)\bigr]
\]
where $C(H)$ denotes the computational effort required to identify history $H$
from among the candidates in the fibre.
\end{definition}

\begin{definition}[Uniform Reconstruction]
A reconstruction procedure is \emph{uniform} if every history in $c^{-1}(o)$
is considered equiprobable given only the observation $o$.
\end{definition}

\begin{theorem}[Entropy Lower Bound on Reconstruction]
For any uniform reconstruction procedure,
\[
\calR(o) \geq k \cdot S_c(o)
\]
for some positive constant $k$ depending only on the cost model $C$.
\end{theorem}

\begin{proof}
The fibre $c^{-1}(o)$ contains $N = |c^{-1}(o)|$ candidate histories. Any
reconstruction procedure must distinguish among these alternatives. A decision
tree distinguishing $N$ equally probable alternatives requires at least $\log_2 N$
binary decisions. Each decision has a positive cost bounded below by $k > 0$.
Since $S_c(o) = \log N$, the expected reconstruction effort satisfies
$\calR(o) \geq k \log N = k \cdot S_c(o)$.
\end{proof}

\begin{corollary}
Observational entropy is a lower bound on reconstruction difficulty. Notes that
reduce observational entropy are therefore notes that reduce reconstruction cost.
\end{corollary}

\begin{remark}
This theorem explains precisely why notes have value in terms the continuation
theory can measure. A note $N$ reduces $|c^{-1}(o)|$ by introducing additional
distinctions that collapse cannot erase. Reducing the fibre size reduces $S_c(o)$
and hence reduces $\calR(o)$. The note acts as a compression-resistant continuation
preserver: it encodes information that the observational collapse would otherwise
destroy, lowering the expected work required to recover the original history.
\end{remark}

\section{Continuation Value and Note Leverage}

\begin{definition}[Continuation Value of an Observation]
The \emph{continuation value} of observation $o$ is
\[
V(o) = \frac{1}{1 + \calR(o)}.
\]
\end{definition}

Observations with large fibres have low continuation value because the
original history is expensive to recover. Observations with small fibres have
high continuation value. The note taxonomy from Part~V can now be given a
quantitative characterization: a note $N$ is valuable to the extent that it
increases $V(o)$ for observations $o$ it affects, which is equivalent to
decreasing $\calR(o)$, which is equivalent to decreasing $S_c(o)$.

\begin{theorem}[Note Leverage as Entropy Reduction]
The reachability leverage $\Lambda(N)$ of a note (Definition 2.3) is bounded
below by the entropy reduction it induces:
\[
\Lambda(N) \geq \frac{k}{C_{\mathrm{create}}(N)} \sum_{o} [S_c(o) - S_c(o \mid N)],
\]
where $S_c(o \mid N)$ denotes the entropy of observation $o$ when the note $N$
is available.
\end{theorem}

\begin{proof}
By the Entropy Lower Bound theorem, entropy reduction lower-bounds reconstruction
cost reduction. Summing over affected observations and dividing by creation cost
gives the claimed bound on leverage.
\end{proof}

\section{The Reconstruction Complexity Hierarchy}

Different collapse rules induce different reconstruction complexity classes.

\begin{definition}[Tractable Reconstruction]
A collapse rule $c$ has \emph{tractable reconstruction} if $\calR(o)$ is bounded
by a polynomial in the size of the history alphabet for all $o \in O_c$.
\end{definition}

\begin{proposition}
The Identity rule $c_I$ has reconstruction complexity $O(1)$: given the
observation (which equals the history), reconstruction is immediate. The
LastWrite rule $c_{LW}$ has reconstruction complexity exponential in the
number of symbols, since the fibre contains all histories consistent with the
last-write state. The Accumulate rule $c_{\mathrm{acc}}$ has reconstruction
complexity equal to the multinomial coefficient $\binom{|\text{events}|}{n_1, \ldots, n_k}$
where $n_i$ is the count for symbol $i$.
\end{proposition}

\chapter{The Geometry of Refusal}

\section{Refusal as Geometric Excision}

Refusal was introduced in Part~II as a documented declaration of inadmissibility.
A deeper interpretation emerges when admissibility is viewed as a region of
possibility space. The admissibility set $\calA(H)$ is not merely a set of
available symbols; it is a geometric region in possibility space, and each refusal
excises a portion of that region.

\begin{definition}[Admissibility Region]
The \emph{admissibility region} after history $H$ is
$\calA(H) = \Omega_0 \setminus \{x \mid \exists r.\, \Refuse(x,r) \in H\}$.
\end{definition}

\begin{definition}[Refusal Region]
A refusal reason $r$ determines a \emph{forbidden subset} $U_r \subseteq \Omega_0$.
A refusal $\Refuse(x, r)$ excises $x$ from the admissibility region under
the constraint $r$.
\end{definition}

\begin{definition}[Refusal Action]
Applying $\Refuse(x, r)$ produces the new admissibility region
\[
\calA(H') = \calA(H) \setminus \{x\}.
\]
When $r$ is a constraint that applies to a family of symbols
$U_r = \{x : r(x) = \mathsf{true}\}$, the excision generalizes to
$\calA(H') = \calA(H) \setminus U_r$.
\end{definition}

Refusal is therefore a geometric excision operator: it removes regions from
possibility space. Multiple refusals compose to produce an increasingly constrained
admissibility region, a process that can be analyzed in terms of the topology of
the resulting space.

\section{Refusal Monotonicity}

\begin{theorem}[Refusal Monotonicity]
For every refusal event, $\calA(H') \subseteq \calA(H)$.
\end{theorem}

\begin{proof}
By construction, $\calA(H') = \calA(H) \setminus U_r$ for some $U_r \subseteq \Omega_0$.
Removing a subset cannot enlarge the region. Hence $\calA(H') \subseteq \calA(H)$.
\end{proof}

\begin{corollary}[Irreversibility of Refusal]
Once a symbol $x$ is refused, it cannot be re-admitted without replacing the
history. Since histories grow monotonically, refusal is irreversible within a
computation.
\end{corollary}

\section{Refusal Curvature}

\begin{definition}[Local Admissibility Curvature]
For a point $x$ in possibility space with neighbourhood $N(x)$,
the \emph{local admissibility curvature} is
\[
\kappa_A(x) = \frac{|N(x) \cap \calA|}{|N(x)|}.
\]
A curvature of $1$ indicates full local admissibility; a curvature of $0$
indicates a refusal singularity.
\end{definition}

\begin{definition}[Refusal Singularity]
A point $x$ is a \emph{refusal singularity} if $\kappa_A(x) = 0$: every
trajectory through $x$ is inadmissible.
\end{definition}

\begin{proposition}[Curvature Monotonicity Under Refusal]
Repeated refusals create regions of non-increasing admissibility curvature:
if $H' = H \mathbin{++} [\Refuse(x, r)]$ then $\kappa_{A'}(y) \leq \kappa_A(y)$
for all $y \in U_r$.
\end{proposition}

\begin{proof}
Each refusal removes admissible trajectories from $U_r$. Therefore
$|N(y) \cap \calA'| \leq |N(y) \cap \calA|$ for $y \in U_r$, which implies
$\kappa_{A'}(y) \leq \kappa_A(y)$.
\end{proof}

\begin{definition}[Admissibility Field]
The function $\kappa_A : \Omega_0 \to [0,1]$ is the \emph{admissibility field}
of the computation. Its level sets $\{x : \kappa_A(x) \geq \lambda\}$ for
threshold $\lambda$ define the $\lambda$-admissible regions of possibility space.
\end{definition}

\begin{remark}
The admissibility field $\kappa_A$ is the discrete computational analogue of
the smooth admissibility field in the RSVP framework. The xylomorphic criterion
$\lambda < 1$ from RSVP corresponds to $\kappa_A(x) < 1$ in the discrete
setting: a trajectory is admissible at level $\lambda$ when its local admissibility
curvature exceeds $\lambda$. The RSVP continuum limit of the Spherepop admissibility
field is therefore obtained by taking the lattice spacing to zero in the discrete
geometry.
\end{remark}

\section{The Galois Connection Between History Growth and Admissibility Contraction}

\begin{theorem}[History--Admissibility Galois Connection]
There is a Galois connection between the lattice of history extensions and the
lattice of admissibility regions, ordered by the relation:
\[
H \leq H' \iff H \text{ is a prefix of } H'
\]
on histories and $\calA \leq \calA'$ iff $\calA \subseteq \calA'$ on admissibility
regions. The Galois connection is given by:
\[
\Phi(H) = \calA(H) \quad \text{and} \quad \Psi(\calA) = \{H : \calA(H) \supseteq \calA\}.
\]
\end{theorem}

\begin{proof}
We verify the Galois condition: $H \leq H'$ implies $\Phi(H) \geq \Phi(H')$
(order reversal). By Refusal Monotonicity, if $H'$ extends $H$ by any sequence
of events, then $\calA(H') \subseteq \calA(H)$, that is $\Phi(H') \leq \Phi(H)$.
Conversely, $\calA \geq \calA'$ (i.e.\ $\calA \supseteq \calA'$) implies
$\Psi(\calA) \leq \Psi(\calA')$ (shorter histories needed to maintain the larger
admissibility set). The adjunction conditions $H \leq \Psi(\Phi(H))$ and
$\calA \leq \Phi(\Psi(\calA))$ follow from the definitions.
\end{proof}

\begin{remark}
The Galois connection formalizes the intuition that history growth and admissibility
contraction are dual processes. As a computation proceeds (history grows), the
space of admissible futures contracts. The Galois connection makes this duality
precise: every increase in historical commitment corresponds to a decrease in
future admissibility, and vice versa.
\end{remark}

\chapter{Notes as Externalized Adjoints}

\section{The Reconstruction Gap}

Collapse maps histories into observations, losing information. Notes partially
reverse this process by preserving information that collapse discards. The
question is: can this reversal be made precise in category-theoretic terms?

The answer is yes, and the result gives the note theory its deepest mathematical
formulation. Notes are partial right adjoints to collapse functors.

\section{The Note Reconstruction Functor}

\begin{definition}[Note Reconstruction Functor]
Given a collapse functor $\mathcal{O}_c : \Hist \to \Obs_c$, a
\emph{note reconstruction functor} is a functor
$\mathcal{N} : \Obs_c \to \Hist$
that assigns to each observable state $o$ a canonical history
$\mathcal{N}(o) \in c^{-1}(o)$.
\end{definition}

\begin{definition}[Adjoint Note]
A note is \emph{adjoint to collapse} if
\[
\Obs_c(\mathcal{O}_c(H), o) \cong \Hist(H, \mathcal{N}(o))
\]
naturally in $H$ and $o$.
\end{definition}

\begin{theorem}[Notes as Partial Right Adjoints]
\label{thm:notes-adjoints}
Every successful note acts as a partial right adjoint to a collapse operation,
in the sense that it provides a reconstruction map from observations back toward
histories on the subdomain where reconstruction succeeds.
\end{theorem}

\begin{proof}
Let $N$ be a note and $c$ a collapse rule such that $N$ preserves information
discarded by $c$ on some subdomain $D \subseteq O_c$. Define
$\mathcal{N}(o) = H_N(o)$ where $H_N(o)$ is the canonical history recovered
using note $N$ from observation $o$. For $o \in D$, the map
$\Hist(H, \mathcal{N}(o)) \to \Obs_c(\mathcal{O}_c(H), o)$ given by
$f \mapsto \mathcal{O}_c \circ f$ is an isomorphism: every history morphism
into $\mathcal{N}(o)$ projects to a unique observable morphism, and conversely,
since $N$ provides the additional data needed to lift observable morphisms to
history morphisms on $D$. This is the universal property of a right adjoint,
restricted to $D$.
\end{proof}

\begin{corollary}
The statement ``notes are anti-collapse artifacts'' is equivalent to saying that
notes approximate right adjoints to collapse functors. The richer the note's
continuation volume, the larger the subdomain $D$ on which the adjunction holds.
\end{corollary}

\section{Examples of Note Adjunctions}

The note-as-adjoint formulation unifies the taxonomy from Part~V.

A \emph{bookmark} is a right adjoint to the collapse of a URL into a browsing
context: it recovers the specific location from an otherwise unlabeled observation
of web content.

A \emph{mathematical proof} is a right adjoint to the collapse of a theorem
into a bare assertion: it recovers the inferential trajectory from the terminal
claim.

\emph{Source code} is a right adjoint to the collapse of executable behavior
into observable outputs: it recovers the computational trajectory from the
input-output specification.

A \emph{map} is a right adjoint to geographical collapse: it recovers spatial
structure from positional observations that contain no direction or distance
information.

In every case, the note $\mathcal{N}$ reconstructs distinctions erased by the
collapse $\mathcal{O}_c$, satisfying the adjunction on the subdomain where
reconstruction is possible.

\section{The Completeness Spectrum}

\begin{definition}[Note Completeness]
A note $N$ is \emph{complete} for collapse rule $c$ if $\mathcal{N}$ extends to a
global right adjoint: $\mathcal{O}_c \dashv \mathcal{N}$ on all of $\Obs_c$.
$N$ is \emph{partial} if the adjunction holds only on a proper subdomain.
\end{definition}

\begin{proposition}
No finite note is complete for a collapse rule that identifies infinitely many
distinct histories. Complete notes for finite history spaces exist and correspond
to observationally complete collapse rules (Proposition 19.3 of Part~VII).
\end{proposition}

\begin{proof}
If $c$ identifies infinitely many histories, then $|c^{-1}(o)| = \infty$ for
some $o$. A finite note contains finite information and can therefore eliminate
at most finitely many ambiguities per observation. Hence it cannot provide a
complete right adjoint on all of $\Obs_c$.
\end{proof}

\chapter{The Topology of Repair}

\section{Repair as a Morphism in the History Category}

Repair Theory holds that repair is the restoration of access to a trajectory
that has been interrupted. Within the Spherepop framework, repair can be
formalized as a special class of morphism in the history category.

\begin{definition}[Repair Morphism]
A \emph{repair morphism} from history $H$ to history $H'$ is a map
$\rho : H \to H'$ such that $c(H) = c(H')$ for some designated collapse rule
$c$: the observable behavior is preserved while the historical structure is altered.
\end{definition}

\begin{definition}[Repair Equivalence]
Two histories are \emph{repair equivalent}, written $H \sim_R H'$, if there
exists a collapse rule $c$ and a repair morphism $\rho : H \to H'$ such that
$c(H) = c(H')$.
\end{definition}

\begin{theorem}[Repair Equivalence Theorem]
Repair preserves observable reachability while modifying historical structure:
if $H \sim_R H'$ then every continuation observable under $c$ from $H$ is
also observable from $H'$, yet $H$ and $H' $ may differ as elements of $\calH$.
\end{theorem}

\begin{proof}
Since $c(H) = c(H')$, the two histories are in the same observational equivalence
class. Every continuation whose reachability depends only on the observable state
$c(H)$ is therefore equally reachable from $H'$. However, since $H \neq H'$ as
sequences of events, their internal structure differs. The repair has altered the
history while preserving its observable footprint.
\end{proof}

\section{Examples of Repair Classes}

The repair morphism formulation unifies several classes of real-world operation.

\emph{Software patches} replace a sequence of code events producing a buggy
observable behavior with a different sequence producing the correct observable
behavior. The patched and unpatched code share the same input-output specification
(the collapse rule is the black-box I/O function) but have different development
histories.

\emph{Biological healing} replaces a damaged tissue history with a regenerated
tissue history. The organism's observable function is preserved (the collapse rule
is the phenotypic function) while the internal molecular history changes.

\emph{Engineering maintenance} replaces a worn component history with a replacement
component history. The system's operational behavior is preserved (the collapse
rule is the performance specification) while the material history differs.

\emph{Institutional reform} replaces a history of governance decisions with a
revised sequence. The institution's external behavior and mandate are preserved
(the collapse rule is the institutional mandate) while the decision history changes.

All four are repair morphisms in the formal sense.

\section{Repair Cohomology}

Not every interrupted trajectory can be repaired. The obstruction to repair is
a topological invariant.

\begin{definition}[Repair Defect]
The \emph{repair defect} of histories $H$ and $H'$ is
$\Delta_R(H, H') = d_{\Hist}(H, H')$ where $d_{\Hist}$ measures the minimum
edit distance in the history category (minimum number of event replacements
needed to convert $H$ into a history with the same observable state as $H'$).
\end{definition}

\begin{definition}[Repair Cohomology Class]
When no repair morphism exists between $H$ and a target observable state $o$
(that is, $c^{-1}(o) \cap \mathsf{Reachable}(H) = \varnothing$), the
obstruction is a non-trivial element of the \emph{repair cohomology}
$H^1_R(\calH, c)$---the first cohomology group of the history space with
coefficients in the collapse sheaf.
\end{definition}

\begin{proposition}[Trivial Cohomology Implies Universal Repairability]
If $H^1_R(\calH, c) = 0$, then every history can be repaired to any
observationally equivalent target. The absence of cohomological obstructions
corresponds to a globally connected repair graph.
\end{proposition}

\begin{remark}
Repair cohomology provides a rigorous mathematical notion of irreparability.
A system with non-trivial repair cohomology has histories from which no
repair morphism can reach certain observable states, regardless of the effort
invested. In practice, this corresponds to: cryptographic irreversibility
(a hash cannot be repaired to recover the preimage), biological cell death
(a necrotic cell cannot be repaired to a living state), and data loss without
redundancy (a destroyed archive without backup admits no repair).
\end{remark}

\chapter{The Noun Fallacy Revisited}

\section{Objects as Stable Quotients}

Throughout this monograph we have argued that histories are primary and states
are derived. The present chapter makes this precise by giving a formal account
of how objects---the apparently primitive entities of everyday ontology---arise
from the collapse of histories.

\begin{definition}[Persistent Quotient]
An equivalence class $[H]_c$ under collapse rule $c$ is \emph{persistent} if
$c(H_t) = c(H_{t+\Delta})$ for all admissible extensions $H_{t+\Delta}$ of $H_t$
in a neighborhood of the current trajectory. That is, the observable state does
not change under small perturbations of the computation.
\end{definition}

\begin{definition}[Object]
An \emph{object} in the Spherepop ontology is a persistent quotient:
\[
\mathsf{Object} = [H]_c \quad \text{where } [H]_c \text{ is persistent.}
\]
\end{definition}

\section{Object Emergence}

\begin{theorem}[Object Emergence Theorem]
Stable observational equivalence classes emerge as objects to observers
who access histories only through collapse rules.
\end{theorem}

\begin{proof}
An observer accesses $\calH$ only through $c$. Repeated observation of the same
equivalence class produces observations $c(H_{t_1}) = c(H_{t_2}) = \cdots = o$.
The observer, receiving identical observations across time, infers an invariant
entity. This entity is $[H]_c$. The observer treats the persistent equivalence
class as an object because it is behaviorally indistinguishable from one: it
produces the same observable state under every observation the rule $c$ permits.
\end{proof}

\section{The Noun Fallacy}

\begin{definition}[Noun Fallacy]
The \emph{noun fallacy} is the error of treating a persistent quotient $[H]_c$
as an ontologically primitive object---as if $[H]_c$ existed independently of
the history $H$ and the collapse rule $c$ that produced it.
\end{definition}

\begin{corollary}
Every noun is a collapse note.
\end{corollary}

\begin{proof}
A noun compresses a family of trajectories into a single retrievable symbol.
It therefore performs the Collapse operation on the histories it subsumes. The
symbol is the collapsed observable; the suppressed histories are the fibres of the
collapse. A noun is precisely a Collapse Note in the sense of Chapter~13.
\end{proof}

\begin{example}
The word \emph{tree} is a collapse note. It collapses the family of all
developmental, ecological, metabolic, and evolutionary histories of tree-shaped
organisms into a single navigable token. The fibre $c_{\mathrm{tree}}^{-1}(\text{``tree''})$
is the set of all histories that produce observable states classifiable as trees.
The noun \emph{tree} makes this fibre navigable without requiring traversal of any
particular history within it.
\end{example}

\begin{example}
The word \emph{program} is a collapse note for computational history. It collapses
the set of all source texts, development decisions, debugging sessions, and
revision histories that produced a given executable into a single noun. The Spherepop
framework aims to reverse this collapse: to make the history that the noun
conceals explicitly accessible.
\end{example}

\section{Reality as Reachable History}

The preceding analysis permits a precise statement of the monograph's deepest claim.

\begin{theorem}[Reality as Reachability]
Under the Spherepop ontology:
\begin{enumerate}
\item Histories are primary: they are the fundamental objects of computation, cognition, and civilization.
\item Observations are quotients: they are equivalence classes of histories under collapse rules.
\item Objects are persistent quotients: stable equivalence classes that observers reify as things.
\item Notes preserve distinctions across quotients: they are artifacts that resist specific collapse operations.
\item Repair restores reachability after interruption: it constructs repair morphisms between interrupted and target histories.
\item Civilization constructs note infrastructure: it builds distributed systems for preserving continuations across generations.
\item Reality is a structured space of reachable histories, not a collection of primitive objects. Objects are what histories look like after sufficient collapse. 
\end{enumerate}
\end{theorem}

\begin{proof}
Items (1)--(6) follow from the formal definitions and theorems developed in
Parts II through VIII. Item (7) is the synthesis: given that objects are
persistent quotients (Definition 29.2) and quotients are derived from histories
(Proposition 2.2), objects depend ontologically on histories. The space of
objects is therefore a derived structure over the primary space of histories.
\end{proof}

%% ============================================================
%% PART X: COLLAPSE AS ORTHOGONAL PROJECTION
%% ============================================================
\part{Collapse as Orthogonal Projection}

\chapter{The Hilbert Space of Histories}

\section{Motivation}

Throughout the preceding chapters, collapse has been treated combinatorially:
a collapse rule $c : \calH \to O_c$ maps histories to equivalence classes, and
observational entropy counts fibre sizes. This combinatorial treatment is adequate
for finite history spaces, but it leaves the continuous and probabilistic dimensions
of collapse undeveloped.

The Hilbert space formulation of conditional expectation, developed in
probability theory through the work of Kolmogorov and later geometrized in
the $L^2$ framework, provides exactly the continuous analogue of Spherepop
collapse. The central thesis of this part is:

\begin{quote}
Every collapse rule introduced in previous chapters may be viewed as an orthogonal
projection operator. In finite settings, collapse appears combinatorial. In
probabilistic settings, collapse becomes orthogonal projection in a Hilbert space.
The two treatments are the same operation at different levels of abstraction.
\end{quote}

\section{The $L^2$ History Space}

\begin{definition}[$L^2$ History Space]
Let $(\Omega, \mathcal{F}, P)$ be a probability space where $\Omega$ is the
set of all histories, $\mathcal{F}$ is a $\sigma$-algebra over $\Omega$, and
$P$ is a probability measure over histories. The \emph{$L^2$ history space} is
\[
L^2(\mathcal{F}) = \{X : \Omega \to \mathbb{R} \mid \mathbb{E}[X^2] < \infty\}
\]
with inner product $\langle X, Y \rangle = \mathbb{E}[XY]$.
\end{definition}

\begin{remark}
Random variables $X \in L^2(\mathcal{F})$ are functions on the space of histories.
They measure observable quantities of histories---durations, counts of specific
event types, the time of the last Pop, the total weight of admissibility
certificates---in a way that respects the probability structure over histories.
\end{remark}

\begin{definition}[Observational Subspace]
Let $\mathcal{G} \subseteq \mathcal{F}$ be a sub-$\sigma$-algebra representing
an observational framework. The \emph{observational subspace} is
\[
L^2(\mathcal{G}) = \{X \in L^2(\mathcal{F}) \mid X \text{ is } \mathcal{G}\text{-measurable}\}.
\]
\end{definition}

$L^2(\mathcal{G})$ is a closed linear subspace of $L^2(\mathcal{F})$. It consists
of exactly those random variables whose values can be determined from the
information available in $\mathcal{G}$---the distinctions that the observational
framework $\mathcal{G}$ makes visible.

\section{Orthogonal Projection as Collapse}

\begin{theorem}[Conditional Expectation as Orthogonal Projection]
For any $X \in L^2(\mathcal{F})$ and any sub-$\sigma$-algebra
$\mathcal{G} \subseteq \mathcal{F}$, the conditional expectation
$\mathbb{E}[X \mid \mathcal{G}]$ is the unique element of $L^2(\mathcal{G})$
that minimizes $\mathbb{E}[(X - Z)^2]$ over all $Z \in L^2(\mathcal{G})$.
Equivalently, $\mathbb{E}[X \mid \mathcal{G}]$ is the orthogonal projection of
$X$ onto $L^2(\mathcal{G})$.
\end{theorem}

\begin{proof}
The closed subspace $L^2(\mathcal{G})$ has a unique nearest point to $X$
by the Hilbert space projection theorem. This nearest point $\hat{X}$ satisfies
the orthogonality condition: $\langle X - \hat{X}, Z \rangle = 0$ for all
$Z \in L^2(\mathcal{G})$. Expanding: $\mathbb{E}[(X - \hat{X})Z] = 0$ for all
$\mathcal{G}$-measurable $Z$. This is precisely Kolmogorov's characterization of
the conditional expectation. By the uniqueness of the projection, $\hat{X} =
\mathbb{E}[X \mid \mathcal{G}]$.
\end{proof}

\begin{definition}[Spherepop Interpretation of Conditional Expectation]
Under the Spherepop identification:
\begin{align*}
\text{History space } \calH &\longleftrightarrow L^2(\mathcal{F}) \\
\text{Observation space } O_c &\longleftrightarrow L^2(\mathcal{G}) \\
\text{Collapse functor } F_c &\longleftrightarrow \mathbb{E}[\cdot \mid \mathcal{G}] \\
\text{Collapse rule } c &\longleftrightarrow \text{Sub-}\sigma\text{-algebra } \mathcal{G} \\
\text{Observational entropy } S_c(o) &\longleftrightarrow \text{Variance of residual}
\end{align*}
\end{definition}

\section{The Orthogonality Condition as State Illusion}

\begin{theorem}[Orthogonality as State Illusion]
The residual $X - \mathbb{E}[X \mid \mathcal{G}]$ is orthogonal to every element
of $L^2(\mathcal{G})$:
\[
\bigl\langle X - \mathbb{E}[X \mid \mathcal{G}],\, Z \bigr\rangle = 0
\quad \text{for all } Z \in L^2(\mathcal{G}).
\]
The discarded component is completely invisible from within the observational
framework that produced the collapse.
\end{theorem}

\begin{proof}
By the projection theorem, $X - \mathbb{E}[X \mid \mathcal{G}]$ is perpendicular
to $L^2(\mathcal{G})$. Expanding: $\mathbb{E}[(X - \mathbb{E}[X\mid\mathcal{G}])Z] =
\mathbb{E}[XZ] - \mathbb{E}[\mathbb{E}[X\mid\mathcal{G}] \cdot Z]$.
By the definition of conditional expectation,
$\mathbb{E}[\mathbb{E}[X\mid\mathcal{G}] \cdot Z] = \mathbb{E}[XZ]$
for all $\mathcal{G}$-measurable $Z$. Therefore the difference is zero.
\end{proof}

\begin{remark}
This theorem is the Hilbert-space version of the state illusion. The observer
with access only to $L^2(\mathcal{G})$ cannot detect the residual component
$X - \mathbb{E}[X\mid\mathcal{G}]$: every measurement they can make is
orthogonal to it, yielding zero correlation. The observer perceives only
$\mathbb{E}[X\mid\mathcal{G}]$ and is unaware that a residual exists. This
is precisely the state illusion: the observer perceives a single observable
state while the underlying history contains information invisible to their
measurement apparatus.
\end{remark}

\chapter{The Algebraic Laws as Geometry}

\section{The Tower Property as Nested Collapse}

The Tower Property of conditional expectation---$\mathbb{E}[\mathbb{E}[X \mid
\mathcal{G}_2] \mid \mathcal{G}_1] = \mathbb{E}[X \mid \mathcal{G}_1]$ when
$\mathcal{G}_1 \subseteq \mathcal{G}_2$---is a geometric theorem about nested
projections.

\begin{theorem}[Tower Property as Nested Projection]
If $\mathcal{G}_1 \subseteq \mathcal{G}_2 \subseteq \mathcal{F}$, then
$L^2(\mathcal{G}_1) \subseteq L^2(\mathcal{G}_2)$, and
\[
\mathbb{E}\bigl[\mathbb{E}[X \mid \mathcal{G}_2] \mid \mathcal{G}_1\bigr]
= \mathbb{E}[X \mid \mathcal{G}_1].
\]
\end{theorem}

\begin{proof}
Since $\mathcal{G}_1 \subseteq \mathcal{G}_2$, every $\mathcal{G}_1$-measurable
function is $\mathcal{G}_2$-measurable. Therefore $L^2(\mathcal{G}_1) \subseteq
L^2(\mathcal{G}_2)$. Let $\hat{X}_2 = \mathbb{E}[X \mid \mathcal{G}_2]$ be
the projection of $X$ onto $L^2(\mathcal{G}_2)$. The residual
$X - \hat{X}_2$ is orthogonal to $L^2(\mathcal{G}_2)$ and hence to the subspace
$L^2(\mathcal{G}_1) \subseteq L^2(\mathcal{G}_2)$. Therefore
$\mathbb{E}[\hat{X}_2 \mid \mathcal{G}_1]$ is the projection of $\hat{X}_2$ onto
$L^2(\mathcal{G}_1)$. But the projection of $X$ onto $L^2(\mathcal{G}_1)$ equals
the projection of $\hat{X}_2 + (X - \hat{X}_2)$ onto $L^2(\mathcal{G}_1)$, and
the residual contributes zero (being orthogonal to $L^2(\mathcal{G}_1)$). Hence
$\mathbb{E}[X \mid \mathcal{G}_1] = \mathbb{E}[\hat{X}_2 \mid \mathcal{G}_1]$.
\end{proof}

\begin{remark}[Spherepop Interpretation]
The Tower Property is the Hilbert-space version of the Spherepop result on
Collapse Functor Composition (Theorem 20.2): applying a coarser collapse after
a finer one is equivalent to applying the coarser collapse directly. In both
settings, the coarser observational framework erases the distinctions preserved
by the finer one. The geometry makes the ``why'' transparent: projecting twice
onto nested subspaces is the same as projecting once onto the smaller subspace.
\end{remark}

\section{The Observational Lattice as Inclusion of Subspaces}

\begin{proposition}[Lattice-Subspace Correspondence]
The lattice of collapse rules ordered by refinement ($c_1 \preceq c_2$ iff
$c_1$ is finer) corresponds to the lattice of closed subspaces of $L^2(\mathcal{F})$
ordered by inclusion: finer rules correspond to larger (more informative)
subspaces.
\[
c_1 \preceq c_2 \quad \iff \quad L^2(\mathcal{G}_{c_1}) \supseteq L^2(\mathcal{G}_{c_2}).
\]
\end{proposition}

\begin{proof}
A finer rule $c_1$ preserves more distinctions, hence its observational subspace
$L^2(\mathcal{G}_{c_1})$ contains more measurable functions---it is a larger
subspace. A coarser rule $c_2$ preserves fewer distinctions---its subspace
$L^2(\mathcal{G}_{c_2})$ is smaller. The inclusion reversal corresponds to the
duality between refinement (more information) and subspace containment (larger
space of representable functions).
\end{proof}

\section{Eve's Law as the Pythagorean Theorem}

The Law of Total Variance (Eve's Law) is not an algebraic identity but a
Pythagorean theorem in $L^2(\mathcal{F})$.

\begin{theorem}[Eve's Law as Pythagorean Theorem]
\[
\mathrm{Var}(X) = \mathbb{E}[\mathrm{Var}(X \mid \mathcal{G})] + \mathrm{Var}(\mathbb{E}[X \mid \mathcal{G}]).
\]
\end{theorem}

\begin{proof}
Center $X$ by writing $\tilde{X} = X - \mathbb{E}[X]$. Decompose:
\[
\tilde{X} = \underbrace{(X - \mathbb{E}[X \mid \mathcal{G}])}_{\text{residual}} +
            \underbrace{(\mathbb{E}[X \mid \mathcal{G}] - \mathbb{E}[X])}_{\text{explained}}.
\]
The residual lies in $L^2(\mathcal{G})^\perp$ (it is the orthogonal complement
of the projection). The explained component lies in $L^2(\mathcal{G})$. These
two vectors are orthogonal. Therefore by the Pythagorean theorem:
\[
\|\tilde{X}\|^2 = \|X - \mathbb{E}[X\mid\mathcal{G}]\|^2 +
                  \|\mathbb{E}[X\mid\mathcal{G}] - \mathbb{E}[X]\|^2.
\]
Taking expectations: $\|\tilde{X}\|^2 = \mathrm{Var}(X)$;
$\|X - \mathbb{E}[X\mid\mathcal{G}]\|^2 = \mathbb{E}[\mathrm{Var}(X\mid\mathcal{G})]$
by the Tower Property; and $\|\mathbb{E}[X\mid\mathcal{G}] - \mathbb{E}[X]\|^2 =
\mathrm{Var}(\mathbb{E}[X\mid\mathcal{G}])$.
\end{proof}

\begin{remark}[Spherepop Interpretation]
Eve's Law decomposes total uncertainty into two orthogonal components:
\begin{itemize}
\item $\mathbb{E}[\mathrm{Var}(X\mid\mathcal{G})]$: \emph{hidden uncertainty}---the
expected residual uncertainty that remains after the observational collapse
$\mathcal{G}$ has occurred. This corresponds to the fibre entropy $S_c(o)$ summed
over observations: the irreducible ambiguity within each equivalence class.
\item $\mathrm{Var}(\mathbb{E}[X\mid\mathcal{G}])$: \emph{visible uncertainty}---the
variance of the observable itself. This corresponds to the variation between
equivalence classes.
\end{itemize}
The total variance $\mathrm{Var}(X)$ is the Pythagorean sum of these two orthogonal
sources. The state illusion hides the first component from the observer who
accesses the history only through $\mathcal{G}$.
\end{remark}

\chapter{Sigma-Algebras as Observational Frameworks}

\section{Information as Geometry}

The $L^2$ framework gives a geometric account of information. An observational
framework $\mathcal{G}$ is not merely a collection of measurable sets; it is a
geometric constraint on what can be known. The subspace $L^2(\mathcal{G})$
consists of exactly those history-valued functions whose values can be determined
given the information in $\mathcal{G}$.

\begin{theorem}[Sigma-Algebra as Information Constraint]
A sub-$\sigma$-algebra $\mathcal{G} \subseteq \mathcal{F}$ corresponds to an
observational framework that makes visible exactly the distinctions in
$L^2(\mathcal{G})$. The following are equivalent:
\begin{enumerate}
\item $X$ is $\mathcal{G}$-measurable.
\item $X \in L^2(\mathcal{G})$.
\item The value of $X$ is completely determined by the observation framework $\mathcal{G}$.
\item $\mathbb{E}[X \mid \mathcal{G}] = X$ (projecting $X$ onto its own subspace leaves it unchanged).
\end{enumerate}
\end{theorem}

\begin{proof}
(1) $\Leftrightarrow$ (2) by definition of $L^2(\mathcal{G})$. (2) $\Leftrightarrow$
(3) because $\mathcal{G}$-measurable functions are exactly those determined by
$\mathcal{G}$-information. (3) $\Leftrightarrow$ (4): if $X \in L^2(\mathcal{G})$,
projecting $X$ onto $L^2(\mathcal{G})$ yields $X$ itself (the projection of a
vector already in a subspace is the vector).
\end{proof}

\section{Collapse Rules and Sub-Sigma-Algebras}

\begin{theorem}[Collapse Rules Are Sub-Sigma-Algebras]
Every Spherepop collapse rule $c : \calH \to O_c$ determines a sub-$\sigma$-algebra
$\mathcal{G}_c$ of the history $\sigma$-algebra $\mathcal{F}$:
\[
\mathcal{G}_c = \{c^{-1}(U) \mid U \subseteq O_c\}.
\]
The collapse functor $F_c$ corresponds to the conditional expectation operator
$\mathbb{E}[\cdot \mid \mathcal{G}_c]$.
\end{theorem}

\begin{proof}
The collection $\{c^{-1}(U) \mid U \subseteq O_c\}$ is closed under
complementation and countable unions (since $c$ is measurable), hence forms a
$\sigma$-algebra. It is a sub-$\sigma$-algebra of $\mathcal{F}$ because every
$c^{-1}(U)$ is a union of fibers of $c$, hence measurable in $\mathcal{F}$.
The conditional expectation $\mathbb{E}[X \mid \mathcal{G}_c]$ then picks out
the component of $X$ that varies only across the equivalence classes of $c$,
which is exactly what the collapse functor $F_c$ does at the combinatorial level.
\end{proof}

\section{The Radon-Nikodym Theorem as Universal Note}

The Radon-Nikodym theorem, which grounds the measure-theoretic definition of
conditional expectation, has a natural interpretation in the Spherepop framework.

\begin{theorem}[Radon-Nikodym as Universal Reconstruction]
Let $\mu$ and $\nu$ be probability measures on $(\Omega, \mathcal{F})$ with
$\nu \ll \mu$ (absolute continuity). Then there exists a unique
$\mathcal{F}$-measurable function $f = d\nu/d\mu$ such that
$\nu(A) = \int_A f \, d\mu$ for all $A \in \mathcal{F}$.
\end{theorem}

\begin{remark}[Spherepop Interpretation]
The Radon-Nikodym derivative $d\nu/d\mu$ is a note in the Spherepop sense: it
preserves the information needed to reconstruct the measure $\nu$ from the
reference measure $\mu$. Without $d\nu/d\mu$, an agent who knows only $\mu$
cannot determine $\nu$. With $d\nu/d\mu$---a measurable function, hence a
Generative Note in the $L^2$ setting---the reconstruction is complete. The
Radon-Nikodym theorem is therefore a universal reconstruction theorem: it
guarantees the existence of the minimal note needed to recover one measure from
another, whenever such recovery is possible.
\end{remark}

\section{Variance as Continuous Observational Entropy}

\begin{theorem}[Variance as Continuous Fibre Entropy]
The conditional variance $\mathrm{Var}(X \mid \mathcal{G})$ is the continuous
analogue of the discrete observational entropy $S_c(o) = \log |c^{-1}(o)|$:
\[
\mathrm{Var}(X \mid \mathcal{G}) = \mathbb{E}\bigl[(X - \mathbb{E}[X \mid \mathcal{G}])^2 \mid \mathcal{G}\bigr]
\]
measures the expected squared deviation of $X$ from its projection onto
$\mathcal{G}$, which is the expected information content of the residual
component invisible to the observational framework.
\end{theorem}

\begin{proof}
The residual $X - \mathbb{E}[X\mid\mathcal{G}]$ is the component of $X$
that the framework $\mathcal{G}$ cannot see. Its conditional variance
$\mathbb{E}[(X - \mathbb{E}[X\mid\mathcal{G}])^2 \mid \mathcal{G}]$ measures
the expected squared magnitude of this invisible component, given the visible
component. This is the continuous analogue of counting the histories in a
fibre: instead of $\log |c^{-1}(o)|$ discrete histories, we have a continuous
measure of the spread of histories within the fibre.
\end{proof}

\begin{remark}[Unification]
This theorem completes the bridge between the discrete Spherepop collapse theory
and the continuous $L^2$ conditional expectation theory. The dictionary is:

\begin{center}
\begin{tabular}{ll}
\toprule
Discrete Spherepop & Continuous $L^2$ \\
\midrule
History space $\calH$ & $L^2(\mathcal{F})$ \\
Collapse rule $c$ & Sub-$\sigma$-algebra $\mathcal{G}$ \\
Collapse functor $F_c$ & Conditional expectation $\mathbb{E}[\cdot\mid\mathcal{G}]$ \\
Observational entropy $S_c(o)$ & Conditional variance $\mathrm{Var}(X\mid\mathcal{G})$ \\
Fibre $c^{-1}(o)$ & Fiber over $\mathcal{G}$-atom \\
State illusion & Orthogonality of residual \\
Note (anti-collapse) & Radon-Nikodym derivative \\
Tower Property & Nested projection theorem \\
Eve's Law & Pythagorean theorem in $L^2$ \\
\bottomrule
\end{tabular}
\end{center}
\end{remark}

\part{Open Problems and Future Directions}

\chapter{The Inverse Problem and Observational Completeness}

\section{When Does Observational Equivalence Imply History Equivalence?}

The deepest open question in the Spherepop framework is the converse of the
collapse functor: when does $c(H_1) = c(H_2)$ imply $H_1 = H_2$?

\begin{definition}[Observational Completeness]
A collapse rule $c$ is \emph{observationally complete} if it is injective:
$c(H_1) = c(H_2) \Rightarrow H_1 = H_2$.
\end{definition}

\begin{definition}[Collapse Kernel]
The \emph{kernel} of a collapse rule $c$ is
\[
\ker(c) = \{(H_1, H_2) \in \calH \times \calH \mid c(H_1) = c(H_2)\}.
\]
Observational completeness holds iff $\ker(c) = \{(H, H) \mid H \in \calH\}$.
\end{definition}

\begin{proposition}[Completeness Results for Canonical Rules]
The Identity rule $c_I$ is observationally complete. The LastWrite rule $c_{LW}$
is not: $[pop(x), pop(y)]$ and $[pop(y), pop(x), pop(y)]$ are
$c_{LW}$-equivalent but distinct. The Accumulate rule $c_{\mathrm{acc}}$ is not:
$[pop(x), pop(y)]$ and $[pop(y), pop(x)]$ have identical accumulate-states.
\end{proposition}

\begin{definition}[Inverse-Problem Fibre]
For $o \in O_c$, the \emph{fibre} of $o$ is
$c^{-1}(o) = \{H \in \calH \mid c(H) = o\}$.
\end{definition}

Characterising fibres for general rules, and determining when a fibre contains a
unique history, is the primary open mathematical problem in the framework.

\section{Open Problems}

\subsection{Distributed Spherepop Runtimes}

The current World is single-threaded. A distributed Spherepop runtime would
coordinate multiple Worlds, each with their own histories, under a global
consistency constraint. The consistency constraint must be expressed in terms of
history equivalence under a shared collapse rule---not value consistency. The
event sourcing and CRDT literatures provide partial models; the Spherepop-specific
challenge is the certified-refusal structure across distributed agents.

\subsection{Collapse Functor Composition}

The four canonical collapse rules are special cases of a more general family of
functors $F_c : \Hist \to \Obs_c$. A theory of functor composition would
characterise which combinations of rules are valid---when $F_{c_2} \circ F_{c_1}$
is itself a valid collapse functor---and would classify the algebraic structure
of the space of collapse rules.

\subsection{Dependent Process Types in Full Generality}

The current Process$(T_1 \rightsquigarrow T_2)$ type is first-order. A dependent
process type $\Pi(x : T_1).\, \mathrm{Process}(x, f(x))$ would allow the type
of the output to depend on the specific value consumed. This is the natural
extension of dependent type theory into the process calculus setting and would
enable precise typing of programs that branch on their inputs.

\subsection{Admissibility Lattice Completeness}

The current type system uses a Boolean admissibility structure. The full theory
requires a graded admissibility lattice with continuous admissibility levels,
connecting to the RSVP xylomorphic criterion $\lambda < 1$ and the admissibility
geometry program. Completeness theorems for the graded type system would
characterize exactly which trajectories are admissible at each level.

\subsection{Proof-Carrying Refusal in Full Generality}

The current $\mathsf{RefusalReason}$ is a vocabulary. A complete proof-carrying
refusal system would make $\mathsf{RefusalReason}$ a proof term: evidence that
can be type-checked independently of the computation that produced it. This
connects to the theory of certified refusal and would enable mechanized
verification of inadmissibility claims.

\chapter{The Coda: Implementation as Philosophical Argument}

The implementation did not merely realize the theory; it exposed which distinctions
the theory was still hiding.

The central claim of this monograph is not that Spherepop is a good programming
language, though we believe it is an interesting one. The central claim is that
taking a philosophical position seriously enough to implement it is a form of
philosophical argument.

Four implementation discoveries were not predicted by the initial framework. They
were forced by the attempt to mechanize it.

The original $\Refuse(\mathsf{Symbol})$ gestured at documented inadmissibility.
It took implementation pressure to reveal that a symbol alone documents nothing:
the reason is the document.

The original notion of observable state gestured at the quotient structure. It
took \texttt{World::observe} to force the question: observable under what rule?

The original type checker gestured at world-relative knowledge. It took the
free-variable conflict to force the question: which world are we assuming?

The original correctness criterion gestured at semantic preservation. It took the
compiler test to force the question: preservation of what, exactly?

Each discovery reveals a distinction the philosophical framework was gesturing at
without making precise. Only in the attempt to mechanize---to force the philosophy
to compile---did the gestures become definitions. This is the deepest lesson of
ontology-driven language design. Philosophy can identify the right questions. Only
implementation can force the right precision.

\medskip

Spherepop is the algebra of continuations.

Notes are the physical realization of that algebra.

Civilization is the attempt, always incomplete, always under threat, to build a
note infrastructure adequate to the continuation needs of the species that requires
it.

\newpage


%% ============================================================
%% PART X: CALCULUS OF CONSTRUCTIONS
%% ============================================================
\part{The Calculus of Constructions}

\chapter{Why Spherepop Needs the Calculus of Constructions}

\section{The Limits of Simple Types}

The type system developed in Part~IV is powerful enough to certify admissibility,
document refusals, and seal observations. But it is limited in one fundamental
respect: types cannot depend on values. The type of the output of a function cannot
depend on what the input actually is. This limitation becomes acute in several
situations that arise naturally in Spherepop.

Consider a function that reads a file and returns a proof that the file was
successfully read. The output type should mention the specific file that was read,
not just the abstract type of files. Or consider a GC pass that returns a history
with only live events: the output type should certify that specific inadmissibility
certificates from the input are preserved. Simple types cannot express either of
these conditions.

The Calculus of Constructions (CoC), introduced by Coquand and Huet~\cite{coquand88},
is the natural remedy. It is a typed lambda calculus in which types may depend on
terms, terms may be types, and the distinction between levels collapses into a
single universe hierarchy. In CoC, the type of a returned value can carry the full
specification of what was computed.

\section{The Pure Type System Perspective}

The Calculus of Constructions belongs to the family of \emph{pure type systems}
(PTSs), which unify the handling of terms and types through a single syntactic
category.

\begin{definition}[Sorts and Pure Type System]
A \emph{pure type system} is determined by a triple $(\mathcal{S}, \mathcal{A},
\mathcal{R})$ where $\mathcal{S}$ is a set of \emph{sorts}, $\mathcal{A}$ is a
set of \emph{axioms} $s_1 : s_2$, and $\mathcal{R}$ is a set of \emph{rules}
$(s_1, s_2, s_3)$ governing when $\Pi$-types can be formed.
\end{definition}

\begin{definition}[Calculus of Constructions]
The Calculus of Constructions is the PTS with:
\begin{align*}
\mathcal{S} &= \{\ast, \square\} \\
\mathcal{A} &= \{\ast : \square\} \\
\mathcal{R} &= \{(\ast, \ast, \ast),\; (\ast, \square, \square),\;
                 (\square, \ast, \ast),\; (\square, \square, \square)\}
\end{align*}
where $\ast$ is the sort of ordinary types and $\square$ is the sort of kinds
(types of types).
\end{definition}

The four rules in $\mathcal{R}$ permit:
\begin{enumerate}
\item $(\ast, \ast, \ast)$: functions from terms to terms (ordinary functions).
\item $(\ast, \square, \square)$: functions from terms to types (dependent types).
\item $(\square, \ast, \ast)$: functions from types to terms (polymorphism).
\item $(\square, \square, \square)$: functions from types to types (type operators).
\end{enumerate}

The full power of CoC is that all four are available simultaneously. This is what
makes it capable of expressing the entire type-theoretic apparatus needed for
Spherepop's admissibility certificates.

\section{The Spherepop-CoC Correspondence}

The Spherepop type system is naturally embedded into CoC.

\begin{definition}[CoC Embedding of Spherepop Types]
The embedding $\llbracket \cdot \rrbracket : \mathsf{SphType} \to \mathsf{CoCTerm}$
is defined by:
\begin{align*}
\llbracket \mathsf{Unit} \rrbracket &= \top_{\mathsf{CoC}} \\
\llbracket \mathsf{Never} \rrbracket &= \bot_{\mathsf{CoC}} \\
\llbracket \Adm(T) \rrbracket &= \Sigma(h : \mathsf{History}).\, \mathsf{Admissible}(h, \llbracket T \rrbracket) \\
\llbracket \Rfsd(T, r) \rrbracket &= \Sigma(h : \mathsf{History}).\, \mathsf{Refused}(h, \llbracket T \rrbracket, r) \\
\llbracket \Col(T, c) \rrbracket &= \Sigma(h : \mathsf{History}).\, \mathsf{Collapsed}(h, \llbracket T \rrbracket, c) \\
\llbracket \Pi(x : T_1).T_2 \rrbracket &= \Pi(x : \llbracket T_1 \rrbracket).\llbracket T_2 \rrbracket \\
\llbracket \Proc(T_1 \rightsquigarrow T_2) \rrbracket &=
  \Pi(h_1 : \llbracket T_1 \rrbracket).\, \Sigma(h_2 : \llbracket T_2 \rrbracket).\,
  \mathsf{BindDep}(h_1, h_2)
\end{align*}
\end{definition}

The $\Sigma$-types in the embedding are crucial: they pair the value with a
proof that the history supporting it has the required structure. The admissibility
certificate is not merely a type annotation but a proof term that can be
independently verified.

\begin{theorem}[Embedding Soundness]
If $\Gamma \vdash t : T$ in the Spherepop type system, then
$\llbracket \Gamma \rrbracket \vdash \llbracket t \rrbracket : \llbracket T \rrbracket$
in CoC.
\end{theorem}

\begin{proof}
By structural induction on the Spherepop typing derivation. Each typing rule maps
to a valid CoC derivation:

$[\mathrm{T\text{-}POP}]$: The embedding sends $\mathsf{pop}(t)$ to a pair
$\langle h, \pi \rangle$ where $\pi$ is a proof of $\mathsf{Admissible}(h, \llbracket T \rrbracket)$.
This is a valid $\Sigma$-introduction in CoC.

$[\mathrm{T\text{-}REFUSE}]$: The embedding sends $\mathsf{refuse}(t, r)$ to a
pair $\langle h, \pi_r \rangle$ where $\pi_r$ proves $\mathsf{Refused}(h, \llbracket T \rrbracket, r)$.
Again a valid $\Sigma$-introduction.

$[\mathrm{T\text{-}COLLAPSE}]$: Collapse requires a proof of admissibility as
input and produces a proof of the collapsed type. The CoC derivation uses
$\Sigma$-elimination on the admissibility proof followed by $\Sigma$-introduction
for the collapsed type.

The remaining cases follow by induction.
\end{proof}

\chapter{Implementing CoC for Spherepop}

\section{The Universe Hierarchy}

A full implementation of CoC for Spherepop requires a universe hierarchy to avoid
Russell-type paradoxes. We adopt a predicative hierarchy.

\begin{definition}[Universe Hierarchy]
Define sorts $\mathsf{Type}_0, \mathsf{Type}_1, \mathsf{Type}_2, \ldots$ with
axioms $\mathsf{Type}_i : \mathsf{Type}_{i+1}$ and cumulative coercions
$\mathsf{Type}_i \hookrightarrow \mathsf{Type}_{i+1}$.
\end{definition}

In practice, most Spherepop typing occurs at $\mathsf{Type}_0$ (the sort of
ordinary program types). The proof-carrying refusal reasons inhabit
$\mathsf{Type}_1$ (propositions about $\mathsf{Type}_0$ terms). The collapse
rules inhabit $\mathsf{Type}_1$ as functions from histories to observable states.

\begin{lstlisting}[style=rust, caption={Universe levels in the Spherepop type checker}]
pub enum Universe {
    Type(usize),   // Type_0, Type_1, ...
    Prop,          // proof-irrelevant propositions
}

pub enum Term {
    Sort(Universe),
    Var(Symbol),
    App(Box<Term>, Box<Term>),
    Lam(Symbol, Box<Term>, Box<Term>),  // λx:A.b
    Pi(Symbol, Box<Term>, Box<Term>),   // Πx:A.B
    Sigma(Symbol, Box<Term>, Box<Term>),// Σx:A.B
    Pair(Box<Term>, Box<Term>),         // (a, b)
    Fst(Box<Term>),                     // π₁
    Snd(Box<Term>),                     // π₂
    // Spherepop primitives
    Pop(Box<Term>),
    Refuse(Box<Term>, RefusalReason),
    Bind(Box<Term>, Box<Term>),
    Collapse(Box<Term>, CollapseRule),
    // Admissibility universe
    Admissible(Box<Term>),
    Refused(Box<Term>, RefusalReason),
    Collapsed(Box<Term>, CollapseRule),
}
\end{lstlisting}

\section{The Type Checker}

A CoC type checker for Spherepop must implement bidirectional type checking:
type inference for introduction forms and type checking for elimination forms.

\begin{definition}[Bidirectional Typing]
Bidirectional typing uses two judgements:
\begin{itemize}
\item $\Gamma \vdash t \Rightarrow T$ (\emph{synthesis}): the type $T$ of $t$
is inferred from $t$ alone.
\item $\Gamma \vdash t \Leftarrow T$ (\emph{checking}): $t$ is checked against
the expected type $T$.
\end{itemize}
\end{definition}

\begin{lstlisting}[style=rust, caption={Bidirectional type checker core}]
pub fn synthesize(ctx: &Context, term: &Term)
    -> Result<Term, TypeError>
{
    match term {
        Term::Var(x) => ctx.lookup(x),
        Term::App(f, a) => {
            let f_ty = synthesize(ctx, f)?;
            match whnf(&f_ty) {
                Term::Pi(x, a_ty, b_ty) => {
                    check(ctx, a, &a_ty)?;
                    Ok(subst(b_ty, x, a))
                }
                _ => Err(TypeError::NotAFunction)
            }
        }
        Term::Fst(p) => {
            let p_ty = synthesize(ctx, p)?;
            match whnf(&p_ty) {
                Term::Sigma(_, a, _) => Ok(*a),
                _ => Err(TypeError::NotAPair)
            }
        }
        Term::Pop(t) => {
            let t_ty = synthesize(ctx, t)?;
            Ok(Term::Admissible(Box::new(t_ty)))
        }
        Term::Refuse(t, r) => {
            let t_ty = synthesize(ctx, t)?;
            Ok(Term::Refused(Box::new(t_ty), r.clone()))
        }
        Term::Collapse(t, c) => {
            let t_ty = synthesize(ctx, t)?;
            match whnf(&t_ty) {
                Term::Admissible(inner) =>
                    Ok(Term::Collapsed(inner, c.clone())),
                _ => Err(TypeError::CollapseRequiresAdmissible)
            }
        }
        _ => Err(TypeError::CannotSynthesize)
    }
}

pub fn check(ctx: &Context, term: &Term, ty: &Term)
    -> Result<(), TypeError>
{
    match (term, whnf(ty)) {
        (Term::Lam(x, _, b), Term::Pi(y, a, b_ty)) => {
            let ctx2 = ctx.extend(x, &a);
            check(&ctx2, b, &subst(&b_ty, y, &Term::Var(x)))
        }
        (Term::Pair(a, b), Term::Sigma(x, a_ty, b_ty)) => {
            check(ctx, a, &a_ty)?;
            let a_val = eval(a);
            check(ctx, b, &subst(&b_ty, x, &a_val))
        }
        _ => {
            let t_ty = synthesize(ctx, term)?;
            if definitionally_equal(&t_ty, ty) { Ok(()) }
            else { Err(TypeError::TypeMismatch) }
        }
    }
}
\end{lstlisting}

\section{Definitional Equality and Normalization}

CoC type checking requires deciding definitional equality: two types are
definitionally equal if they reduce to the same normal form.

\begin{definition}[Weak Head Normal Form]
A term is in \emph{weak head normal form} (WHNF) if it is not a beta-redex at
the head: it is a lambda, a Pi, a Sigma, a sort, a variable, or an application
whose function is in WHNF.
\end{definition}

\begin{lstlisting}[style=rust, caption={WHNF reduction for Spherepop-CoC}]
pub fn whnf(term: &Term) -> Term {
    match term {
        Term::App(f, a) => {
            match whnf(f) {
                Term::Lam(x, _, b) => whnf(&subst(&b, &x, a)),
                f_whnf => Term::App(Box::new(f_whnf),
                                    a.clone())
            }
        }
        Term::Fst(p) => {
            match whnf(p) {
                Term::Pair(a, _) => whnf(&a),
                p_whnf => Term::Fst(Box::new(p_whnf))
            }
        }
        Term::Snd(p) => {
            match whnf(p) {
                Term::Pair(_, b) => whnf(&b),
                p_whnf => Term::Snd(Box::new(p_whnf))
            }
        }
        // Spherepop reductions
        Term::Pop(t) => Term::Pop(Box::new(whnf(t))),
        Term::Collapse(t, c) => {
            Term::Collapse(Box::new(whnf(t)), c.clone())
        }
        t => t.clone()
    }
}
\end{lstlisting}

\begin{theorem}[Normalization of Spherepop-CoC]
The Spherepop-CoC type system is strongly normalizing: every well-typed
reduction sequence terminates.
\end{theorem}

\begin{proof}[Proof sketch]
The Calculus of Constructions is strongly normalizing by the reducibility
candidates argument of Girard~\cite{girard89}. The Spherepop extensions Pop,
Refuse, Bind, and Collapse are all introduction forms that do not introduce
new beta-redexes. They reduce only by unwrapping their arguments to WHNF.
The combined system inherits strong normalization from CoC by a standard
extension argument: the new rules add no new reduction paths that could create
infinite sequences, since the underlying lambda calculus reductions are already
terminating.
\end{proof}

\section{Identity Types and History Equality}

CoC extended with identity types (Martin-L\"of type theory~\cite{martinlof84})
gives Spherepop a native notion of proof-level history equality.

\begin{definition}[Identity Type]
For any type $A$ and terms $a, b : A$, the \emph{identity type} $a =_A b$ is a
type whose inhabitants are proofs that $a$ and $b$ are definitionally equal.
The sole constructor is $\mathsf{refl}_a : a =_A a$.
\end{definition}

\begin{definition}[History Identity Type]
For histories $H_1, H_2 : \mathsf{History}$, the type $H_1 =_{\mathsf{History}} H_2$
is inhabited precisely when the histories are event-for-event identical. This is
the type-level statement of Replay Equivalence: a term of type
$I(p).\mathsf{history} =_{\mathsf{History}} V(C(p)).\mathsf{history}$
is a formal proof that the compiler is correct.
\end{definition}

\begin{lstlisting}[style=rust, caption={History identity in the type system}]
pub fn verify_replay_equivalence(
    prog: &Term,
    interp_world: &World,
    compiled_world: &World,
) -> Result<(), ProofError> {
    // Produce a proof of history identity
    if interp_world.history == compiled_world.history {
        Ok(())  // refl witnesses the identity
    } else {
        // Construct a counterexample witness
        let diff = history_diff(
            &interp_world.history,
            &compiled_world.history
        );
        Err(ProofError::HistoryMismatch(diff))
    }
}
\end{lstlisting}

\section{The Curry-Howard Correspondence for Spherepop}

The Curry-Howard correspondence---that proofs are programs and propositions are
types---takes a specific form in Spherepop.

\begin{proposition}[Spherepop Curry-Howard Correspondence]
The following identifications hold:
\begin{center}
\begin{tabular}{lll}
\toprule
Logical Notion & Program Notion & Spherepop Realization \\
\midrule
Proposition & Type & $T : \mathsf{Type}_0$ \\
Proof & Term & $t : T$ \\
Implication & Function type & $\Pi(x:A).B$ \\
Conjunction & Pair type & $\Sigma(x:A).B$ \\
Admissibility & Certificate type & $\Adm(T)$ \\
Refutation & Refusal type & $\Rfsd(T, r)$ \\
Observation & Collapsed type & $\Col(T, c)$ \\
History identity & Replay equivalence & $H_1 =_\mathcal{H} H_2$ \\
\bottomrule
\end{tabular}
\end{center}
\end{proposition}

%% ============================================================
%% PART XI: THE BNF GRAMMAR AS SUPERIOR NOTATION
%% ============================================================
\part{The BNF Grammar as Universal Notation}

\chapter{Why Grammar Beats Circuits and Code}

\section{The Problem of Notation}

Every computational formalism needs a notation: a way of writing down the objects
it studies. The choice of notation is not merely aesthetic. It determines what
can be expressed without ambiguity, what can be processed mechanically, what can
be read without specialized tools, and what can be implemented independently of
any particular hardware or software substrate.

Three dominant notations for computation exist: programming language syntax,
circuit diagrams, and formal grammars. The thesis of this chapter is that
Backus-Naur Form (BNF) grammar, or its extensions, is the superior notation for
specifying computational behavior, and that Spherepop's use of BNF is not a
convenience but a theoretical necessity.

\section{A Brief History of BNF}

Backus-Naur Form was introduced by Peter Naur in the ALGOL 60 report~\cite{naur60}.
Its purpose was to give a precise, implementation-independent specification of a
programming language's syntax. Before BNF, programming language grammars were
described in English prose: ambiguous, difficult to parse mechanically, and
dependent on shared human conventions.

BNF replaces prose with a minimal formal apparatus: nonterminals (enclosed in
angle brackets), terminals (written as themselves), a production arrow ($::=$),
and alternation ($\mid$). From these four elements, any context-free language can
be specified.

\begin{definition}[BNF Grammar]
A \emph{Backus-Naur Form grammar} is a tuple $G = (N, T, P, S)$ where:
\begin{itemize}
\item $N$ is a finite set of \emph{nonterminals}.
\item $T$ is a finite set of \emph{terminals} (disjoint from $N$).
\item $P$ is a finite set of \emph{productions} of the form $A ::= \alpha$
where $A \in N$ and $\alpha \in (N \cup T)^*$.
\item $S \in N$ is the \emph{start symbol}.
\end{itemize}
\end{definition}

Extended BNF (EBNF) adds optional elements ($[\cdot]$), repetition ($\{\cdot\}$),
and grouping, but adds no expressive power beyond context-free languages.

\section{BNF Versus Circuit Diagrams}

Circuit diagrams are the standard notation for hardware computation. They have
significant advantages: they are visually intuitive, they make data flow explicit,
and they can be directly compiled to hardware. But they have several critical
disadvantages relative to BNF.

\subsection{Implementation Dependence}

A circuit diagram is inherently tied to a specific level of abstraction. A gate
diagram for a full adder specifies AND, OR, and NOT gates in a specific topology.
This specification cannot be directly used at a different level of abstraction
without redrawing. A BNF grammar for the same operation specifies the structure
independently of implementation:

\begin{lstlisting}[style=rust, caption={BNF for binary addition (implementation-independent)}]
<bit>     ::= "0" | "1"
<addend>  ::= <bit> | <bit> <addend>
<sum>     ::= <addend> "+" <addend>
<carry>   ::= <bit>
<result>  ::= <carry> <sum>
\end{lstlisting}

This grammar specifies the structure of binary addition without committing to
gates, lookup tables, carry-lookahead circuits, or any other implementation.
Any implementation that accepts the grammar's language is correct by definition.

\subsection{Human Readability}

Circuit diagrams require specialized visual parsing. They cannot be read aloud,
stored in plain text, or processed by standard text tools. A BNF grammar is
written in ASCII (or Unicode), can be stored in a version control repository,
diffed, searched, and processed by any text processor. It is read left-to-right,
top-to-bottom, in the same direction as natural language.

\begin{theorem}[BNF Readability Advantage]
For any grammar $G$ specifying a language $L$, the BNF representation of $G$
requires only the characters of $N \cup T$ plus the BNF metasymbols
$\{::=, |, \langle, \rangle\}$. The circuit representation of the same computation
requires a graphical medium that cannot be stored or transmitted as plain text.
\end{theorem}

\begin{proof}
BNF is defined over a finite alphabet. Its productions are strings over that
alphabet. A circuit diagram, by contrast, represents spatial relationships (the
topology of the circuit) that cannot be serialized as a string without loss of
the spatial information that makes the diagram legible. The best-known lossless
serialization of a circuit diagram is a netlist, which is itself a grammar-like
notation---a point in favor of the grammar approach.
\end{proof}

\subsection{Compositional Structure}

BNF grammars compose naturally. If $G_1$ specifies language $L_1$ and $G_2$
specifies $L_2$, the grammar for $L_1 \cdot L_2$ (concatenation) is formed by
adding a production $S ::= S_1\, S_2$ and combining the production sets. This
compositional structure is the foundation of modular language design.

\begin{proposition}[Grammar Compositionality]
If $G_1 = (N_1, T_1, P_1, S_1)$ and $G_2 = (N_2, T_2, P_2, S_2)$ with
$N_1 \cap N_2 = \varnothing$, then:
\begin{align*}
G_1 \cdot G_2 &= (N_1 \cup N_2 \cup \{S\},\, T_1 \cup T_2,\,
                  P_1 \cup P_2 \cup \{S ::= S_1\, S_2\},\, S) \\
G_1 \mid G_2 &= (N_1 \cup N_2 \cup \{S\},\, T_1 \cup T_2,\,
                  P_1 \cup P_2 \cup \{S ::= S_1 \mid S_2\},\, S)
\end{align*}
These operations correspond to sequential composition and choice in Spherepop.
\end{proposition}

\section{BNF Versus Programming Language Syntax}

BNF is more fundamental than any particular programming language syntax because
it is what programming languages are defined in. A programming language's syntax
is a grammar; BNF is the metalanguage for grammars.

But the comparison is not merely definitional. BNF grammars have concrete
advantages as computational specifications over program text.

\subsection{Ambiguity is Detectable}

A BNF grammar can be checked for ambiguity (multiple parse trees for the same
string). An English prose description of a language cannot. A program in any
conventional language can be syntactically unambiguous but semantically ambiguous
in ways that BNF would surface if the semantics were given as a grammar.

\subsection{Separation of Concerns}

BNF cleanly separates the \emph{shape} of a program (its grammar) from its
\emph{meaning} (its semantics). This separation is enforced by the formalism:
a BNF grammar says nothing about what programs mean, only about what strings
are valid programs. This is not a limitation but a virtue: it allows the same
grammar to be given multiple semantics (operational, denotational, axiomatic)
without changing the specification of valid programs.

\subsection{Formal Verification Target}

A BNF grammar is a mathematical object that can serve as the specification for
formal verification. A Spherepop program can be verified to be within the language
of a grammar by a parser; verification at the semantic level then proceeds over
the parse tree. This two-stage verification is not possible when the language is
specified in prose or implicitly in a reference implementation.

\section{The Spherepop BNF as Computational Universal}

The Spherepop BNF grammar from Part~III is not merely a syntactic convenience.
It is a specification from which an implementation in any language, on any
substrate, can be mechanically derived.

\begin{definition}[Spherepop Core Grammar]
The Spherepop core language is given by the following BNF:
\begin{align*}
\langle\mathit{term}\rangle &::= \langle\mathit{var}\rangle
  \mid \langle\mathit{lambda}\rangle
  \mid \langle\mathit{app}\rangle
  \mid \langle\mathit{let}\rangle
  \mid \langle\mathit{pop}\rangle
  \mid \langle\mathit{refuse}\rangle
  \mid \langle\mathit{bind}\rangle
  \mid \langle\mathit{collapse}\rangle
  \mid \langle\mathit{seq}\rangle \\
\langle\mathit{var}\rangle &::= \langle\mathit{ident}\rangle \\
\langle\mathit{lambda}\rangle &::= (\lambda \mid \mathsf{fn})\;
  \langle\mathit{ident}\rangle\; [:\langle\mathit{type}\rangle]\;
  (\to \mid .)\; \langle\mathit{term}\rangle \\
\langle\mathit{app}\rangle &::= \langle\mathit{term}\rangle\;
  \langle\mathit{term}\rangle \\
\langle\mathit{let}\rangle &::= \mathsf{let}\; \langle\mathit{ident}\rangle\;
  = \langle\mathit{term}\rangle\; \mathsf{in}\; \langle\mathit{term}\rangle \\
\langle\mathit{pop}\rangle &::= \mathsf{pop}\; \langle\mathit{term}\rangle \\
\langle\mathit{refuse}\rangle &::= \mathsf{refuse}\; \langle\mathit{term}\rangle\;
  [\langle\mathit{reason}\rangle] \\
\langle\mathit{bind}\rangle &::= \mathsf{bind}\; \langle\mathit{term}\rangle\;
  \langle\mathit{term}\rangle \\
\langle\mathit{collapse}\rangle &::= \mathsf{collapse}\;
  \langle\mathit{term}\rangle\; [\langle\mathit{rule}\rangle] \\
\langle\mathit{seq}\rangle &::= \mathsf{seq}\;
  [\langle\mathit{term}\rangle\; \{;\langle\mathit{term}\rangle\}^*] \\
\langle\mathit{reason}\rangle &::= \mathsf{violation}(\langle\mathit{ident}\rangle)
  \mid \mathsf{explicit}(\langle\mathit{ident}\rangle) \\
\langle\mathit{rule}\rangle &::= \mathsf{id} \mid \mathsf{last\_write}
  \mid \mathsf{accumulate} \mid \mathsf{proj}(\langle\mathit{ident}\rangle) \\
\langle\mathit{type}\rangle &::= \mathsf{Unit} \mid \mathsf{Never}
  \mid \langle\mathit{type\_var}\rangle
  \mid \Adm(\langle\mathit{type}\rangle)
  \mid \langle\mathit{type}\rangle \to \langle\mathit{type}\rangle
\end{align*}
\end{definition}

\begin{theorem}[Grammar Determines Implementation]
Given the Spherepop core BNF, a parser, type checker, and evaluator are
mechanically derivable for any target platform.
\end{theorem}

\begin{proof}
The grammar is context-free and therefore parseable by any LALR(1) or PEG
parser generator. The type checker follows from the typing rules (which are
themselves a kind of grammar over typing contexts). The evaluator follows from
the operational semantics (which are also a grammar of world-state transitions).
Each of these derivations is mechanical and platform-independent: the same BNF
yields implementations in Rust, Python, Haskell, JavaScript, or any other host
language without modification to the grammar.
\end{proof}

\subsection{Grammar as Interface Contract}

When the Spherepop grammar is fixed, it constitutes an interface contract between:
(1) program authors, who write terms in the grammar's language;
(2) parsers, which accept exactly the grammar's language;
(3) type checkers, which verify admissibility of parsed terms;
(4) evaluators, which execute well-typed terms;
(5) compilers, which translate well-typed terms to Event IR.

Any implementation satisfying this contract is a correct Spherepop implementation.
The grammar is the single source of truth for what the language is, independent
of any particular reference implementation.

\section{BNF as Anti-Collapse Artifact}

The relationship between BNF grammars and the note theory of this monograph is
not accidental.

A BNF grammar is a Refusal Note of the highest generativity. It specifies which
strings are admissible (members of the language) and which are inadmissible. It
refuses the collapse of the language into any particular implementation. And it
generates the full continuation space of all programs in the language: any string
in the language is a reachable continuation, and the grammar specifies exactly
the set of reachable continuations.

\begin{theorem}[BNF Grammars as Generative Refusal Notes]
A BNF grammar $G$ is simultaneously:
\begin{enumerate}
\item A \emph{Generative Note}: it creates the continuation space $L(G)$ of all
programs in the language, which did not exist before the grammar was specified.
\item A \emph{Refusal Note}: it refuses all strings not in $L(G)$ as
inadmissible, with the implicit reason $\mathsf{ParseError}$.
\item A \emph{Coordination Note}: it couples the behaviors of all
implementations by binding them to the same accepted and rejected strings.
\item A \emph{Capture Note}: it preserves the structure of valid programs across
implementations, allowing programs written for one implementation to be run
on another.
\end{enumerate}
\end{theorem}

\begin{proof}
(1) The language $L(G)$ is the set of all strings derivable from $S$ in $G$.
Before $G$ exists, this set is undefined; after $G$ exists, it is precisely
specified. New programs in $L(G)$ are continuations that $G$ makes reachable.

(2) A string $w \notin L(G)$ has no parse tree under $G$. The parser refuses
it with a parse error. This is a Refuse operation with reason
$\mathsf{ParseError}$ applied to the inadmissible string.

(3) Any two implementations $I_1$ and $I_2$ of $G$ accept the same strings and
reject the same strings. They are bound by the grammar to the same
admissibility decisions. This is a Bind coupling their trajectory spaces.

(4) A program $p \in L(G)$ has an implementation-independent parse tree. The
parse tree captures the structure of $p$ independently of what evaluator
processes it. This is a Capture Note for the program's syntactic structure.
\end{proof}

\chapter{Grammar, Circuits, and the Efficiency Hierarchy}

\section{Measuring Expressiveness Per Symbol}

We want a precise way to compare notations for expressing computational structure.
Let the \emph{expressiveness per symbol} of a notation $\mathcal{N}$ for
specification $S$ be the ratio of the information content of $S$ to the number
of symbols required to express $S$ in $\mathcal{N}$.

\begin{definition}[Specification Complexity]
The \emph{specification complexity} $K_\mathcal{N}(S)$ of specification $S$ in
notation $\mathcal{N}$ is the minimum number of symbols in $\mathcal{N}$ required
to express $S$ unambiguously. The \emph{relative expressiveness} of $\mathcal{N}_1$
over $\mathcal{N}_2$ for $S$ is
\[
E(S, \mathcal{N}_1, \mathcal{N}_2) = \frac{K_{\mathcal{N}_2}(S)}{K_{\mathcal{N}_1}(S)}.
\]
$E > 1$ means $\mathcal{N}_1$ is more expressive per symbol.
\end{definition}

\begin{theorem}[BNF Specification Efficiency]
For context-free languages, BNF is asymptotically more symbol-efficient than
circuit diagrams and at least as efficient as any other notation that is
implementation-independent and mechanically parseable.
\end{theorem}

\begin{proof}[Proof sketch]
A circuit diagram for a function $f : \{0,1\}^n \to \{0,1\}^m$ specifies $f$
by encoding its topology as a directed acyclic graph of gates. The minimum size
of such a diagram is $\Omega(C(f))$ where $C(f)$ is the circuit complexity of
$f$. A BNF grammar for the same function, specified as a language recognizer,
can be written in $O(\log n + \log m + K(f))$ productions where $K(f)$ is the
Kolmogorov complexity of $f$'s description. For functions with short
descriptions but high circuit complexity (such as sorting networks), BNF can be
exponentially more compact.
\end{proof}

\section{The Readability Spectrum}

Different notations occupy different positions on the spectrum from machine-efficient
to human-efficient.

\begin{center}
\begin{tabular}{llll}
\toprule
Notation & Human Readable & Machine Parseable & Impl.-Independent \\
\midrule
Machine code & No & Yes & No \\
Assembly & Partial & Yes & No \\
Circuit diagram & Yes (visual) & Partial & No \\
C/Rust program & Yes & Yes & Partial \\
Lambda calculus & Yes & Yes & Yes \\
BNF grammar & Yes & Yes & Yes \\
Natural language & Yes & No & Yes \\
\bottomrule
\end{tabular}
\end{center}

BNF occupies the unique position of being simultaneously human-readable,
mechanically parseable, and fully implementation-independent. No other common
notation achieves all three.

\section{Grammar as the Missing Level of Abstraction}

Programming languages, circuit diagrams, and prose all operate at specific
levels of abstraction. Grammar operates at the level of \emph{specification}:
it describes the structure of all possible instances of a computational artifact
without committing to any particular instance.

\begin{definition}[Abstraction Level Hierarchy]
The abstraction levels of computational description form a hierarchy:
\begin{enumerate}
\item \emph{Physical}: transistor layouts, voltages, timing.
\item \emph{Circuit}: gate-level topology.
\item \emph{Machine}: instruction set architecture.
\item \emph{Language}: program syntax and semantics.
\item \emph{Grammar}: specification of the language itself.
\item \emph{Meta-grammar}: specification of the notation for grammars.
\end{enumerate}
\end{definition}

\begin{proposition}[Grammar as Implementation Firewall]
A grammar at level $n$ of the hierarchy is implementation-independent with
respect to all levels below $n$. A Spherepop BNF grammar is therefore
independent of all language implementations, all machine instruction sets, all
circuit designs, and all physical substrates.
\end{proposition}

This independence is exactly what makes BNF appropriate as the specification
language for Spherepop. The Spherepop grammar specifies what computations the
language can express without committing to how they are executed. The execution
details---Rust runtime, WASM compilation, hardware acceleration---are below the
abstraction level at which the grammar operates.

\section{Literate Grammar and Documentation}

Knuth's literate programming~\cite{knuth84} proposed that programs should be
written as literature: prose explanations interleaved with code. The Spherepop
approach inverts this slightly: grammar should be written as literature, with
prose explanations interleaved with BNF productions.

\begin{definition}[Literate Grammar]
A \emph{literate grammar} is a BNF grammar augmented with:
\begin{enumerate}
\item Prose explanations of the intent of each production.
\item Examples of derivations in the grammar.
\item Semantic notes indicating what well-formed strings mean.
\item Transformation rules mapping productions to type-checking and evaluation rules.
\end{enumerate}
\end{definition}

The Spherepop BNF in this monograph is a literate grammar in this sense: each
production is accompanied by the operational semantics rule it corresponds to
and the type-checking rule it enables. This makes the grammar simultaneously
a specification, a tutorial, and a formal verification target.

%% ============================================================
%% PART XII: SUPPLEMENTARY DERIVATIONS
%% ============================================================
\part{Supplementary Derivations}

\chapter{Derivation Compendium}

This chapter collects all supplementary derivations, filling the mathematical
gaps indicated in the preceding parts. Each derivation is labeled by the chapter
it supports.

\section{Chapter 1: Notes Increase Reachability}

For each continuation $c \in \calC$, define the \emph{reachability gain} of note
$N$ by
\[
G_A(N, c, t) = P_A(c, t \mid N) - P_A(c, t).
\]
Then $N$ is a note precisely when $\exists c \in \calC,\; G_A(N, c, t) > 0$.

\begin{proposition}[Nontrivial Notes Are Reachability-Increasing]
If $N$ is a note and all non-note effects are nonnegative, then the total
reachability gain
$G_A(N, t) = \sum_{c \in \calC} G_A(N, c, t)$
is strictly positive.
\end{proposition}

\begin{proof}
By definition, there exists at least one $c_0$ such that $G_A(N, c_0, t) > 0$.
If all other terms satisfy $G_A(N, c, t) \geq 0$, then the total sum contains
one strictly positive term and no negative terms. Hence $G_A(N, t) > 0$.
\end{proof}

\section{Chapter 2: State as Quotient}

\begin{proposition}[State Is a Quotient of History]
Every surjective state projection $\sigma : \calH \to S$ factors uniquely through
the quotient map $q : \calH \to \calH/{\sim_\sigma}$.
\end{proposition}

\begin{proof}
Define $\bar\sigma([H]) = \sigma(H)$. This is well-defined because
$H_1 \sim_\sigma H_2$ implies $\sigma(H_1) = \sigma(H_2)$. Surjectivity makes
$\bar\sigma$ onto; injectivity follows because distinct classes have distinct
$\sigma$-values. Thus $S \cong \calH/{\sim_\sigma}$.
\end{proof}

\section{Chapter 3: Operator Minimality}

\begin{theorem}[Operator Minimality]
No one of $\Pop, \Refuse, \Bind, \Collapse$ is definable from the other three
while preserving its effect on history, option space, admissibility, and
observation.
\end{theorem}

\begin{proof}
$\Pop$ uniquely decreases $\Omega$; the other three leave $\Omega$ unchanged.
$\Refuse$ uniquely changes the admissibility set $\calA(H)$ without decreasing
$\Omega$; $\Pop$ changes $\Omega$, while $\Bind$ and $\Collapse$ do not record
inadmissibility. $\Bind$ uniquely introduces a dependency relation between two
elements without consuming or refusing either; no combination of the other
operators produces a pure dependency edge. $\Collapse$ uniquely maps a history
through an observational rule into a quotient space; the other operators extend
history internally but produce no observation. Hence all four are irreducible.
\end{proof}

\section{Chapter 4: Free Category}

\begin{theorem}[Histories Form the Free Category on the Event Graph]
Let $G$ be the directed graph whose edges are primitive events. Then $\Hist$ is
the free category generated by $G$.
\end{theorem}

\begin{proof}
Objects are option spaces; morphisms are finite paths. Given any category
$\mathcal{D}$ and graph morphism $F : G \to U(\mathcal{D})$, the unique functor
$\bar{F} : \Hist \to \mathcal{D}$ extending $F$ sends $[e_1, \ldots, e_n]$ to
$F(e_n) \circ \cdots \circ F(e_1)$. This is the universal property of the free
category.
\end{proof}

\section{Chapter 5: Small-Step Determinacy}

\begin{proposition}[Determinacy Under Fixed Rule Selection]
If the next operator and its arguments are fixed, Spherepop small-step reduction
is deterministic.
\end{proposition}

\begin{proof}
Each transition rule has a unique conclusion for fixed inputs. No fixed operator
application can reduce to two distinct worlds.
\end{proof}

\section{Chapter 6: Collapse Entropy Monotonicity}

\begin{theorem}[Collapse Increases Observational Entropy Under Coarsening]
If $c_2$ is coarser than $c_1$, then every $c_2$-observation has entropy at
least as large as the entropy of any $c_1$-observation contained inside it.
\end{theorem}

\begin{proof}
Coarsening means each $c_1$-class is contained in some $c_2$-class. The fibre
of the coarser observation contains the fibre of the finer observation, so
$|c_2^{-1}(o_2)| \geq |c_1^{-1}(o_1)|$ for contained observations. Taking
logarithms gives the entropy inequality.
\end{proof}

\section{Chapter 7: Type Soundness}

\begin{theorem}[Spherepop Type Soundness]
If $\vdash t : T$ and $t \to^* t'$, then either $t'$ is a value of type $T$,
or there exists $t''$ such that $t' \to t''$.
\end{theorem}

\begin{proof}
By repeated application of Preservation and Progress. Preservation gives
$\vdash t' : T$ after any finite reduction sequence. Progress implies $t'$ is
either a value or can take another step.
\end{proof}

\section{Chapter 8: Lattice Collapse Threshold}

\begin{proposition}[Monotonicity of Collapse Permission]
If $\ell_1 \leq \ell_2$ and $\Gamma \vdash t : \Adm_{\ell_2}(T)$, then every
collapse permitted at level $\ell_1$ is also permitted at level $\ell_2$.
\end{proposition}

\begin{proof}
A higher admissibility level satisfies at least the constraints of a lower level.
Since $\ell_1 \leq \ell_2$, any rule whose threshold is met by $\ell_1$ is also
met by $\ell_2$.
\end{proof}

\section{Note Composition Theory}

Define note composition by
\[
\calC(N_1 \otimes N_2) = \calC(N_1) \cup \calC(N_2) \cup \calC_{\mathrm{bind}}(N_1, N_2),
\]
where $\calC_{\mathrm{bind}}(N_1, N_2)$ is the set of continuations reachable
only because both notes are jointly available.

\begin{theorem}[Coordination Surplus]
If $\calC_{\mathrm{bind}}(N_1, N_2) \neq \varnothing$, then
\[
V_C(N_1 \otimes N_2) > V_C(N_1) + V_C(N_2) - |\calC(N_1) \cap \calC(N_2)|.
\]
\end{theorem}

\begin{proof}
The continuation set of the composite contains the union of the individual sets
plus at least one additional bound continuation. The ordinary union has size
$V_C(N_1) + V_C(N_2) - |\calC(N_1) \cap \calC(N_2)|$. Adding the nonempty
bound set strictly increases this.
\end{proof}

\begin{theorem}[Note Entropy]
Define the \emph{note entropy} of $N$ relative to agent $A$ at time $t$ as
\[
H_A(N, t) = -\sum_{c \in \calC} P_A(c, t \mid N) \log P_A(c, t \mid N).
\]
A note $N'$ is more \emph{focused} than $N$ if $H_A(N', t) < H_A(N, t)$: it
concentrates reachability on a smaller set of continuations. A note $N'$ is more
\emph{generative} than $N$ if $V_C(N') > V_C(N)$.
\end{theorem}

\begin{remark}
Focus and generativity trade off: collapse notes are highly focused (low entropy)
while generative notes are high-entropy. The note taxonomy orders by this tradeoff.
\end{remark}

\section{Note Class Derivations}

\begin{proposition}[Capture Adequacy]
A capture note $N$ is adequate for task family $T$ iff for every continuation
$c \in T$ there exists recoverable data $d \subseteq N$ such that
$P_A(c, t \mid d) > P_A(c, t)$.
\end{proposition}

\begin{theorem}[Prospective Notes Are Deferred Pops]
Every prospective note determines a deferred commitment operator. At time $t$
the continuation is not yet executed (no immediate Pop occurs); at time $t'$ the
agent follows the note, selecting the future continuation via a Pop.
\end{theorem}

\begin{proof}
A prospective note binds present state $A_t$ to future continuation $c_{t'}$.
The selection at $t'$ is a Pop of $c_{t'}$ from the available alternatives.
The note stores a commitment whose execution is deferred. \end{proof}

\begin{theorem}[Repair Notes Minimize Reconstruction Cost]
A repair note $N$ is optimal in family $\mathcal{F}$ when
$C_r(c \mid N) = \min_{M \in \mathcal{F}} C_r(c \mid M)$.
\end{theorem}

\begin{proof}
Repair value is $C_r(c) - C_r(c \mid N)$. Since $C_r(c)$ is fixed for the
continuation, maximizing repair value is equivalent to minimizing $C_r(c \mid N)$.
\end{proof}

\begin{proposition}[Checksums as Pure Refusal Notes]
A checksum $\chi(h)$ refuses informational collapse by distinguishing corrupted
histories from admissible ones. If $\chi(h') \neq \chi(h)$, then $h'$ is refused
as a valid continuation of $h$. The checksum is a pure refusal note: it stores
only the boundary between admissible and inadmissible continuations.
\end{proposition}

\begin{theorem}[Generativity Is Strict Reachability Expansion]
A note $N$ is generative iff $\exists c : P_A(c,t) = 0$ and $P_A(c,t \mid N) > 0$.
\end{theorem}

\begin{proof}
If such $c$ exists, $N$ creates access to a previously unreachable trajectory.
Conversely, generativity requires opening at least one previously unavailable
trajectory, which is exactly this condition.
\end{proof}

\begin{theorem}[Collapse Notes Trade Detail for Access]
Let $\calN$ be a finite note collection. A collapse note $K$ partitioning $\calN$
into $m$ classes reduces search cost from $O(|\calN|)$ to $O(m + \max_i |\calN_i|)$.
\end{theorem}

\begin{proof}
Without organizing structure, worst-case search examines every note. With a
partition of $m$ classes, the agent selects one class ($m$ steps) then searches
that class ($\max_i |\calN_i|$ steps worst case).
\end{proof}

\section{Compilation and GC Derivations}

\begin{theorem}[Event IR Is History-Preserving]
For every primitive source event $e$, compilation followed by VM execution appends
the same event to history as interpretation.
\end{theorem}

\begin{proof}
By case analysis: $\Pop$ compiles to \texttt{IrEvent::Pop}, which the VM appends
as $\Pop$; $\Refuse(r)$ compiles to \texttt{IrEvent::Refuse\{r\}}, appended as
$\Refuse(r)$; analogously for $\Bind$ and $\Collapse(c)$.
\end{proof}

\begin{theorem}[GC Is Safe Exactly When It Preserves Live Fibres]
A GC operation $g : \calH \to \calH$ is safe iff for every live continuation $c$,
$P_A(c, t \mid g(H)) = P_A(c, t \mid H)$.
\end{theorem}

\begin{proof}
If the equality holds for all live continuations, GC has removed no history needed
for any reachable future---hence safe. Conversely, safety by definition preserves
all history required for live continuations.
\end{proof}

\section{Error Correction Derivation}

\begin{theorem}[Certified Refusal Prevents Silent Error Absorption]
If every refusal carries a proof term $\pi$, then no error can be erased without
either being handled or remaining visible in the output type.
\end{theorem}

\begin{proof}
A refusal has type $\Rfsd(T, \pi)$. The typing rules allow this type to be handled
by a recovery term or propagated. There is no rule converting $\Rfsd(T, \pi)$ into
$T$ while discarding $\pi$. Hence the error cannot be silently absorbed.
\end{proof}

\section{Category Theory Derivations}

\begin{proposition}[Collapse as Reflection, Corrected]
The version of this result stated in an earlier draft---``$F_c : \Hist \to
\Obs_c$ is left adjoint to the inclusion $\Obs_c \hookrightarrow \Hist$, where
$\Obs_c$ is the full subcategory of $\Hist$ of histories saturated under
$\sim_c$''---does not typecheck against either standing definition of $\Hist$
in this book: the History Category (Chapter~\ref{ch:collapse-quotient}) has
option spaces as objects and histories as morphisms, so histories cannot also
be objects of a subcategory of it; and even granting a second, histories-as-objects
reading of $\Hist$, an arbitrary $c$ need not respect the prefix structure
needed for $F_c$ to be a well-defined functor at all, exactly as an arbitrary
$c$ need not respect concatenation in Theorem~1.101.

The corrected statement, proved for prefix-monotone collapse rules with an
order-reflecting canonical-representative section, is
Theorem~\ref{thm:reflection-repaired} in Chapter~\ref{ch:collapse-quotient}'s
companion adjunction chapter: $\mathcal{O}_c \dashv s_c$ exhibits $(O_c,
\preceq_c)$ as a reflective subcategory of the prefix category
$\Hist_\sqsubseteq$, with the Accumulate rule as a worked example and LastWrite
as a worked non-example. Nothing here concerns attested Collapse; it is a
statement about the pure view $\mathcal{O}_c$, per
Section~\ref{sec:view-vs-collapse}.
\end{proposition}

\begin{corollary}[Perfect Observation Requires History]
An observation map is perfectly history-faithful iff it is isomorphic to the
identity observation on $\calH$.
\end{corollary}

\begin{proof}
Perfect history-faithfulness requires injectivity. A bijective observation
establishes an isomorphism between histories and observations, hence is isomorphic
to the identity.
\end{proof}

\section{Sheaf Derivations}

\begin{theorem}[State Illusion as Failure of Unique Gluing]
Given a cover of observations $\{U_i\}$, state illusion occurs exactly when
compatible local sections $s_i \in \calH(U_i)$ admit more than one global
section.
\end{theorem}

\begin{proof}
Compatible local sections agree on overlaps. If more than one global history
restricts to the same local observations, the observations do not determine a
unique history---precisely the state illusion.
\end{proof}

\section{CLIO Derivation}

\begin{theorem}[Spherepop Observational Entropy Is Discrete CLIO Entropy]
The CLIO fibre entropy $S_\pi(o) = \log \mathrm{Vol}(\pi^{-1}(o))$ reduces in
the discrete setting to $S_c(o) = \log |c^{-1}(o)|$.
\end{theorem}

\begin{proof}
In the discrete setting, volume is counting measure. Therefore
$\mathrm{Vol}(\pi^{-1}(o)) = |\pi^{-1}(o)| = |c^{-1}(o)|$.
\end{proof}

\section{Active Geodesic Derivation}

\begin{proposition}[Histories Are Discrete Geodesics]
Let $L(e_i)$ be the local cost of event $e_i$ and define the action
$\mathcal{S}(H) = \sum_i L(e_i)$. A Spherepop execution is a discrete geodesic
when it minimizes $\mathcal{S}(H)$ among admissible histories with the same
boundary conditions.
\end{proposition}

\begin{proof}
In discrete variational mechanics, a geodesic minimizes action between fixed
endpoints. A Spherepop history is a path in event space; restricting to admissible
histories and minimizing the additive cost functional gives the discrete geodesic
condition.
\end{proof}

\section{MEM\textbar8 Locality Derivation}

\begin{theorem}[Locality as Continuation Preservation]
If related memory components are physically or logically local, then the expected
cost of reconstructing their shared continuation is lower than under random
placement.
\end{theorem}

\begin{proof}
Let $d(x, y)$ be access distance between memory components. Reconstruction cost
is monotone in aggregate distance among components required by a continuation.
Local placement lowers expected pairwise distance relative to random placement,
hence lowers expected reconstruction cost.
\end{proof}

\section{Civilizational Collapse Derivation}

\begin{theorem}[Reachability Loss Under Archive Destruction]
Let $c \in R(Z)$ be reachable only through some $N \in \calN$. Destroying $\calN$
strictly decreases collective reachability.
\end{theorem}

\begin{proof}
Since $c$ is reachable only through $N$, removing $N$ removes $c$ from $R(Z)$.
The post-destruction reachability is a proper subset of the original.
\end{proof}

\section{Inverse Problem Derivation}

\begin{proposition}[Kernel Triviality Characterizes Complete Observation]
A collapse rule $c$ is observationally complete iff
$\ker(c) = \Delta_\calH = \{(H, H) : H \in \calH\}$.
\end{proposition}

\begin{proof}
Observational completeness means $c(H_1) = c(H_2) \Rightarrow H_1 = H_2$, so
the kernel contains only diagonal pairs. Conversely, a diagonal kernel gives
injectivity.
\end{proof}

\section{Final Unification Theorem}

\begin{theorem}[Notes as Reachability Artifacts]
\label{thm:unification}
An artifact $N$ is a note iff it is a reachability artifact: there exists an
agent $A$, a continuation $c$, and a future time $t$ such that
$P_A(c, t \mid N) > P_A(c, t)$.
Spherepop programs are notes because they preserve, restrict, bind, refuse, and
collapse continuations in history space.
\end{theorem}

\begin{proof}
The first claim is the definition of a note under the continuation account. For
the second: every Spherepop program constructs a history from primitive events.
$\Pop$ selects continuations; $\Refuse$ preserves inadmissibility distinctions;
$\Bind$ couples trajectories; $\Collapse$ makes histories observable under rules.
A Spherepop program is therefore an artifact that changes which continuations
remain reachable to future agents. Hence it is a note.
\end{proof}

\begin{corollary}[Universal Reachability Artifact Thesis]
Programs, proofs, archives, maps, protocols, libraries, institutions, languages,
and civilizations are all special cases of the same abstract object: a reachability
artifact. They differ in the medium through which they preserve continuations and
in the scale at which they operate. The abstract structure is identical in every
case.
\end{corollary}

\begin{proof}
By Theorem~\ref{thm:unification}, any artifact that increases the probability of
traversing at least one continuation is a note. Each of the listed artifacts does
this:

A \emph{program} preserves a computational continuation.
A \emph{proof} preserves an inferential continuation.
An \emph{archive} preserves historical continuations against decay.
A \emph{map} creates navigational continuations.
A \emph{protocol} binds communicative continuations between agents.
A \emph{library} is a repository of continuations across domains.
An \emph{institution} binds social and coordination continuations across
generations.
A \emph{language} is a system of semantic continuations shared by a community.
A \emph{civilization} is the totality of all of the above, organized into a
distributed hierarchy of note infrastructure.

In every case, the note increases $P_A(c, t)$ for some continuation $c$.
\end{proof}


%% ============================================================
%% PART XV: PROCESS ALGEBRA AND CONCURRENCY
%% ============================================================
\part{Process Algebra and Concurrency}

\chapter{Spherepop as a Process Algebra}

\section{Motivation}

The four Spherepop operators---Pop, Refuse, Bind, Collapse---were derived in
Part~II as the minimal complete set for managing possibility space in sequential
computation. But computation is rarely purely sequential. Processes run in
parallel, communicate, and synchronize. This chapter develops the concurrent
extension of Spherepop and shows that the history-primary ontology handles
concurrency more naturally than state-centric models.

The key insight is that concurrency is a source of state illusion. When two
processes $P$ and $Q$ run in parallel and access shared state, the observable
outcome depends on the interleaving order---the specific history of which events
occurred first. A state-centric model can only see the final state, losing the
interleaving information that explains race conditions, deadlocks, and
non-deterministic behavior. A history-primary model retains the interleaving
and makes it explicit.

\section{Parallel Composition}

\begin{definition}[Parallel World]
A \emph{parallel world} is a pair $W_\parallel = (H_1 \parallel H_2, \Omega_1
\cup \Omega_2)$ where $H_1$ and $H_2$ are histories of concurrent processes
with disjoint option spaces $\Omega_1 \cap \Omega_2 = \varnothing$.
\end{definition}

\begin{definition}[Interleaving Semantics]
The \emph{interleaving semantics} of $W_\parallel$ is the set of all histories
$H$ that are consistent interleavings of $H_1$ and $H_2$:
\[
\mathsf{IL}(H_1, H_2) = \{H \in \calH \mid H \text{ is a shuffle of } H_1
\text{ and } H_2\},
\]
where a \emph{shuffle} of $H_1$ and $H_2$ is any sequence containing all
events of both, with the relative order within each preserved.
\end{definition}

\begin{proposition}[Interleaving Count]
The number of distinct interleavings of histories of lengths $m$ and $n$ is
$\binom{m+n}{m}$.
\end{proposition}

\begin{proof}
An interleaving is determined by choosing which $m$ of the $m+n$ positions
are occupied by events from $H_1$; the remaining $n$ positions are occupied
by events from $H_2$ in order. There are $\binom{m+n}{m}$ such choices.
\end{proof}

\begin{remark}
The combinatorial explosion in $\binom{m+n}{m}$ is the formal content of the
observation that concurrent programs are hard to test. A state-centric test
observes only one interleaving; the Spherepop history preserves which one was
executed, making the interleaving explicit and auditable.
\end{remark}

\section{Synchronization via Bind}

\begin{definition}[Synchronization Event]
A \emph{synchronization event} between processes $P$ and $Q$ at symbols
$a \in \Omega_1$ and $b \in \Omega_2$ is the event $\Bind(a, b)$ appended to
the history when both processes commit to their respective symbols simultaneously.
\end{definition}

\begin{theorem}[Bind as Synchronization]
In a parallel world $W_\parallel$, the event $\Bind(a, b)$ with $a \in \Omega_1$
and $b \in \Omega_2$ is semantically equivalent to a channel communication in
the $\pi$-calculus: process $P$ sends $a$ on channel $ab$ while process $Q$
receives $b$.
\end{theorem}

\begin{proof}
In the $\pi$-calculus~\cite{milner80}, the communication $\bar{a}b.P \mid a(x).Q$
reduces to $P \mid Q[b/x]$, recording that $b$ was transmitted on channel $a$.
In Spherepop, $\Bind(a, b)$ records the dependency between $a$ and $b$ in the
history without consuming either from their option spaces. After the bind, any
future collapse rule can verify that $a$ and $b$ are coupled. The structural
correspondence holds: both formalisms record the communication event without
immediately collapsing its participants.
\end{proof}

\section{Race Conditions as Non-Deterministic Collapse}

\begin{definition}[Race Condition]
A \emph{race condition} occurs in a parallel world $W_\parallel$ when two or
more processes attempt to commit ($\Pop$) to the same symbol $x \in \Omega_1
\cap \Omega_2$ simultaneously, and the observable outcome depends on which
commits first.
\end{definition}

\begin{theorem}[Race Conditions Are Collapse Artifacts]
Race conditions are not intrinsic to the computation but arise from applying
a collapse rule that identifies interleavings differing only in which process
committed first.
\end{theorem}

\begin{proof}
Let $H_1$ be the history where process $P$ commits to $x$ first, and $H_2$ the
history where $Q$ commits first. The histories $H_1 \neq H_2$ are distinct.
However, under a collapse rule $c$ that discards ordering information
(such as $c_{\mathrm{acc}}$), $c(H_1) = c(H_2)$. The observable state is
identical despite the different histories. A programmer who observes only
$c(H_1)$ cannot determine whether $P$ or $Q$ committed first---the race is
invisible. But in the Spherepop world, $H_1 \neq H_2$ is a fact preserved in
the history, accessible by switching to the Identity rule $c_I$.
\end{proof}

\section{Deadlock as Inadmissibility Cycle}

\begin{definition}[Dependency Graph]
The \emph{dependency graph} $G_H$ of a history $H$ has one node for each
symbol in $\Omega_0$ and a directed edge $a \to b$ for each $\Bind(a, b) \in H$.
\end{definition}

\begin{theorem}[Deadlock as Graph Cycle]
A parallel computation is deadlocked at state $(H, \Omega_1 \cup \Omega_2)$
if and only if the dependency graph $G_H$ contains a directed cycle $a_1 \to
a_2 \to \cdots \to a_k \to a_1$ where each $a_i$ is waiting on the next in
its current admissibility set.
\end{theorem}

\begin{proof}
Deadlock occurs when no process can proceed: every process is waiting for a
resource held by another. In Spherepop, a process holds resource $x$ by having
executed $\Pop(x)$ (committing to $x$) and waits for resource $y$ when $y \notin
\calA(H)$ due to a dependency recorded by $\Bind(x, y)$. A cycle $a_1 \to a_2
\to \cdots \to a_k \to a_1$ in $G_H$ means $a_1$ waits for $a_2$, $a_2$ waits
for $a_3$, and so on cyclically, which is precisely the formal definition of
deadlock.
\end{proof}

\begin{corollary}[Deadlock Detection]
Deadlock detection in Spherepop reduces to cycle detection in the dependency
graph $G_H$, which is computable in $O(|V| + |E|)$ time by depth-first search.
\end{corollary}

\chapter{The Process Calculus of Spherepop}

\section{Terms and Their Meanings}

The Spherepop process calculus extends the core term language with parallel
composition, restriction, and replication.

\begin{definition}[Extended Term Language]
\[
\langle\mathit{proc}\rangle ::=
  \langle\mathit{term}\rangle
  \mid \langle\mathit{proc}\rangle \parallel \langle\mathit{proc}\rangle
  \mid \nu x.\langle\mathit{proc}\rangle
  \mid {!}\langle\mathit{proc}\rangle
  \mid \mathbf{0}
\]
where $P \parallel Q$ is parallel composition, $\nu x.P$ restricts the scope
of symbol $x$ to $P$, ${!}P$ is the replication of $P$ (an unbounded number
of copies), and $\mathbf{0}$ is the terminated process.
\end{definition}

\begin{definition}[Structural Congruence]
Processes are \emph{structurally congruent}, written $P \equiv Q$, if they
differ only in syntactic form:
\begin{align*}
P \parallel Q &\equiv Q \parallel P \\
P \parallel (Q \parallel R) &\equiv (P \parallel Q) \parallel R \\
P \parallel \mathbf{0} &\equiv P \\
\nu x.\mathbf{0} &\equiv \mathbf{0} \\
{!}P &\equiv P \parallel {!}P
\end{align*}
\end{definition}

\begin{theorem}[Structural Congruence Preserves History Equivalence]
If $P \equiv Q$, then for any initial world $W$, the sets of histories produced
by $P$ and $Q$ are equal:
\[
\mathsf{Hist}(P, W) = \mathsf{Hist}(Q, W).
\]
\end{theorem}

\begin{proof}
By induction on the structural congruence rules. Commutativity of $\parallel$
permutes the order of interleaved events but does not change which events occur.
Associativity regroups parallel composition without affecting the event set.
The identity law $P \parallel \mathbf{0} \equiv P$ holds because $\mathbf{0}$
contributes no events. The replication law ${!}P \equiv P \parallel {!}P$
reflects the coinductive nature of replication: history sets of both sides are
identical by the fixed-point equation.
\end{proof}

\section{Reduction Rules for Processes}

\begin{definition}[Process Reduction]
Process reduction $P \to P'$ is defined by:
\begin{align*}
&\mathsf{pop}(a) \parallel \mathsf{bind}(a, b).Q \to Q[\text{linked}(a,b)]
  \quad [\mathrm{SYNC}] \\
&\frac{P \to P'}{P \parallel Q \to P' \parallel Q} \quad [\mathrm{PAR}] \\
&\frac{P \to P'}{P \parallel Q \to P \parallel Q'} \quad [\mathrm{SYMM}] \\
&\frac{P \equiv P' \to Q' \equiv Q}{P \to Q} \quad [\mathrm{STRUCT}]
\end{align*}
\end{definition}

\begin{theorem}[Subject Reduction for Processes]
If $\Gamma \vdash P : T$ and $P \to P'$, then $\Gamma \vdash P' : T$.
\end{theorem}

\begin{proof}
By case analysis on the reduction rules. The key case is $[\mathrm{SYNC}]$:
$\mathsf{pop}(a)$ has type $\Adm(\mathsf{Unit})$ by $[\mathrm{T\text{-}POP}]$;
$\mathsf{bind}(a,b)$ has type $\Proc(T_a \rightsquigarrow T_b)$ by
$[\mathrm{T\text{-}BIND}]$. After synchronization, the linked state
$\text{linked}(a,b)$ is typed by the composition of these certificates.
\end{proof}

\section{Bisimulation and Behavioral Equivalence}

\begin{definition}[History Bisimulation]
A binary relation $\mathcal{R}$ on processes is a \emph{history bisimulation}
if whenever $(P, Q) \in \mathcal{R}$:
\begin{enumerate}
\item For every step $P \xrightarrow{H} P'$, there exists $Q'$ such that
$Q \xrightarrow{H} Q'$ and $(P', Q') \in \mathcal{R}$.
\item Symmetrically with $P$ and $Q$ swapped.
\end{enumerate}
where $P \xrightarrow{H} P'$ means $P$ produces history $H$ and reduces to $P'$.
\end{definition}

\begin{theorem}[Bisimulation Is Congruence]
History bisimulation $\sim$ is a congruence: if $P \sim Q$ then
$P \parallel R \sim Q \parallel R$ for all $R$.
\end{theorem}

\begin{proof}
Let $\mathcal{R} = \{(P \parallel R, Q \parallel R) \mid P \sim Q\}$. We show
$\mathcal{R}$ is a bisimulation. Any step of $P \parallel R$ is either a step
of $P$ alone, a step of $R$ alone, or a synchronization between $P$ and $R$.
For steps of $P$ alone: by $P \sim Q$, there is a matching step of $Q$,
giving $Q \parallel R$ a matching step. For steps of $R$: $R$ matches itself.
For synchronizations: the Bind event is recorded identically in both histories,
so $Q \parallel R$ can perform the same synchronization.
\end{proof}

\chapter{Distributed Spherepop}

\section{Multiple Worlds and Consistency}

A distributed Spherepop system consists of multiple worlds $(H_1, \Omega_1),
\ldots, (H_n, \Omega_n)$ running at different locations, each maintaining its
own history. Consistency requires that the histories agree on shared observations.

\begin{definition}[Distributed World System]
A \emph{distributed world system} is a tuple
$\mathcal{W} = (W_1, \ldots, W_n, \sim_{\mathsf{glob}})$
where each $W_i = (H_i, \Omega_i)$ is a local world and
$\sim_{\mathsf{glob}}$ is a global consistency relation specifying which
pairs of events across worlds must agree.
\end{definition}

\begin{definition}[Global Consistency]
A distributed world system $\mathcal{W}$ is \emph{globally consistent} under
collapse rule $c$ if
\[
c(H_i)|_{\Omega_i \cap \Omega_j} = c(H_j)|_{\Omega_i \cap \Omega_j}
\quad \text{for all } i \neq j,
\]
where $c(H)|_\Omega$ denotes the restriction of the observation to symbols in
$\Omega$.
\end{definition}

\begin{theorem}[CAP as Collapse Tradeoff]
In a distributed world system subject to network partitions, no system can
simultaneously guarantee:
\begin{enumerate}
\item \emph{Consistency}: all worlds observe the same collapse $c(H)$.
\item \emph{Availability}: every non-failed world responds to requests.
\item \emph{Partition tolerance}: the system continues operating during network failures.
\end{enumerate}
At most two of three can hold at any given time.
\end{theorem}

\begin{proof}[Proof sketch]
During a network partition, worlds $W_i$ and $W_j$ cannot communicate. If both
must remain available, they each process independent events, producing histories
$H_i$ and $H_j$ that diverge. At partition resolution, $c(H_i)$ and $c(H_j)$
may disagree on shared symbols, violating consistency. To maintain consistency,
one world must refuse to process events (sacrificing availability).
This is the formal Spherepop statement of the CAP theorem~\cite{shapiro11}.
\end{proof}

\section{CRDTs as Admissibility-Preserving Merges}

Conflict-free replicated data types (CRDTs) achieve eventual consistency by
ensuring that concurrent updates commute. In Spherepop terms, CRDTs are merge
operations that preserve admissibility.

\begin{definition}[History Merge]
A \emph{history merge} of $H_1$ and $H_2$ is any history $H_1 \sqcup H_2$
such that both $H_1$ and $H_2$ are subsequences of $H_1 \sqcup H_2$, preserving
the relative order of events within each.
\end{definition}

\begin{definition}[CRDT in Spherepop]
A \emph{Spherepop CRDT} is a data type whose update operations are Pop events
that commute under the collapse rule: for any two updates $P_x$ and $P_y$,
\[
c(H \mathbin{++} [P_x, P_y]) = c(H \mathbin{++} [P_y, P_x]).
\]
\end{definition}

\begin{theorem}[CRDTs Achieve Eventual Consistency]
A Spherepop CRDT system achieves eventual consistency: after all updates
propagate, all worlds reach the same observable state.
\end{theorem}

\begin{proof}
Since updates commute under $c$, any merge $H_1 \sqcup H_2$ produces the same
observable state $c(H_1 \sqcup H_2)$ regardless of the interleaving chosen.
After all updates have propagated to all worlds, each world's history is a
permutation of the same multiset of events. By commutativity, all produce the
same observable state.
\end{proof}

\begin{example}[Grow-Only Counter as Accumulate CRDT]
A grow-only counter in Spherepop is a collection of Pop events on a symbol $x$,
with the Accumulate collapse rule $c_{\mathrm{acc}}(H)(x) = |\{i : e_i = \Pop(x)\}|$.
Any two concurrent increments commute under $c_{\mathrm{acc}}$ because addition
is commutative: regardless of the order $\Pop(x), \Pop(x)$ are interleaved,
the count is the same. This is the formal Spherepop account of the G-Counter CRDT.
\end{example}


%% ============================================================
%% PART XVI: FULL TYPE-THEORETIC METATHEORY
%% ============================================================
\part{Full Type-Theoretic Metatheory}

\chapter{The Complete Proof System}

\section{Motivation}

The type system of Part~IV established Progress, Preservation, and Collapse
Soundness with proof sketches. A type theorist demands the full metatheory:
canonical forms, a weakening lemma, an exchange lemma, a substitution lemma,
and the complete proofs of Progress and Preservation deriving from them. This
chapter provides that apparatus following the syntactic approach of Wright and
Felleisen~\cite{wright94}.

\section{Canonical Forms}

\begin{lemma}[Canonical Forms]
\label{lem:canonical}
If $\vdash v : T$ (well-typed value in empty context) then:
\begin{enumerate}
\item If $T = \mathsf{Unit}$ then $v = ()$.
\item If $T = \Pi(x:T_1).T_2$ then $v = \lambda x:T_1.b$ for some $b$.
\item If $T = \Adm(T')$ then $v = \mathsf{pop}(v')$ for some $v'$ with
$\vdash v' : T'$.
\item If $T = \Rfsd(T', r)$ then $v = \mathsf{refuse}(v', r)$ for some $v'$
with $\vdash v' : T'$.
\item If $T = \Col(T', c)$ then $v = \mathsf{collapse}(v', c)$ for some $v'$
with $\vdash v' : \Adm(T')$.
\item If $T = \Proc(T_1 \rightsquigarrow T_2)$ then $v = \mathsf{bind}(v_1, v_2)$
for some $v_1 : T_1$ and $v_2 : T_2$.
\end{enumerate}
\end{lemma}

\begin{proof}
By induction on the typing derivation $\vdash v : T$ in the empty context.
Since $v$ is a value, it is not a redex. Each type is introduced by exactly
one rule, and values are the results of those introduction rules. We proceed
case by case.

\textit{Case $T = \mathsf{Unit}$}: The only introduction form for $\mathsf{Unit}$
is the unit value $()$, so $v = ()$.

\textit{Case $T = \Pi(x:T_1).T_2$}: The only introduction form for $\Pi$-types
is $[\mathrm{T\text{-}LAM}]$, which produces $v = \lambda x:T_1.b$.

\textit{Case $T = \Adm(T')$}: The only rule producing $\Adm(T')$ is
$[\mathrm{T\text{-}POP}]$, which requires $v = \mathsf{pop}(v')$.

\textit{Case $T = \Rfsd(T', r)$}: Only $[\mathrm{T\text{-}REFUSE}]$ produces
this type, requiring $v = \mathsf{refuse}(v', r)$.

\textit{Case $T = \Col(T', c)$}: Only $[\mathrm{T\text{-}COLLAPSE}]$ produces
this type, requiring the inner term to have type $\Adm(T')$.

\textit{Case $T = \Proc(T_1 \rightsquigarrow T_2)$}: Only $[\mathrm{T\text{-}BIND}]$
produces this type.
\end{proof}

\section{Structural Lemmas}

\begin{lemma}[Weakening]
\label{lem:weakening}
If $\Gamma \vdash t : T$ and $x \notin \mathrm{dom}(\Gamma)$, then
$\Gamma, x : S \vdash t : T$ for any type $S$.
\end{lemma}

\begin{proof}
By induction on the derivation $\Gamma \vdash t : T$.

\textit{Base case} $[\mathrm{T\text{-}VAR}]$: If $t = y$ and $(y : T) \in \Gamma$,
then $(y : T) \in \Gamma, x : S$ (since $x \neq y$). Apply $[\mathrm{T\text{-}VAR}]$.

\textit{Case} $[\mathrm{T\text{-}LAM}]$: $t = \lambda y:T_1.b$ with $\Gamma, y:T_1
\vdash b : T_2$. By induction hypothesis (choosing fresh $x$),
$\Gamma, x:S, y:T_1 \vdash b : T_2$. By exchange (Lemma~\ref{lem:exchange}),
$\Gamma, y:T_1, x:S \vdash b : T_2$. Apply $[\mathrm{T\text{-}LAM}]$.

\textit{Cases} $[\mathrm{T\text{-}POP}]$, $[\mathrm{T\text{-}REFUSE}]$,
$[\mathrm{T\text{-}BIND}]$, $[\mathrm{T\text{-}COLLAPSE}]$: Apply the induction
hypothesis to all premises, then re-apply the rule.

\textit{Case} $[\mathrm{T\text{-}APP}]$: Apply induction to both premises.
\end{proof}

\begin{lemma}[Exchange]
\label{lem:exchange}
If $\Gamma, x:S, y:T, \Delta \vdash t : U$, then
$\Gamma, y:T, x:S, \Delta \vdash t : U$.
\end{lemma}

\begin{proof}
By induction on the derivation. All rules are insensitive to the order of
entries in $\Gamma$ that do not match the current variable being looked up.
The only sensitive case is $[\mathrm{T\text{-}VAR}]$, which looks up exactly
one variable; swapping two different bindings leaves the lookup result unchanged.
\end{proof}

\begin{lemma}[Substitution]
\label{lem:substitution}
If $\Gamma, x:S \vdash t : T$ and $\Gamma \vdash s : S$, then
$\Gamma \vdash t[x := s] : T[x := s]$.
\end{lemma}

\begin{proof}
By induction on the derivation $\Gamma, x:S \vdash t : T$.

\textit{Case} $[\mathrm{T\text{-}VAR}]$ with $t = x$, $T = S$:
$t[x:=s] = s$ and $T[x:=s] = S[x:=s] = S$ (since $x \notin \mathrm{FV}(S)$
in well-formed contexts). We need $\Gamma \vdash s : S$, which is given.

\textit{Case} $[\mathrm{T\text{-}VAR}]$ with $t = y \neq x$:
$t[x:=s] = y$ and the binding $(y:T) \in \Gamma$ (not involving $x$).
Apply $[\mathrm{T\text{-}VAR}]$.

\textit{Case} $[\mathrm{T\text{-}LAM}]$: $t = \lambda y:T_1.b$ with
$\Gamma, x:S, y:T_1 \vdash b : T_2$. Choose $y$ fresh so $y \neq x$.
By induction, $\Gamma, y:T_1 \vdash b[x:=s] : T_2[x:=s]$.
Then $t[x:=s] = \lambda y:T_1[x:=s].b[x:=s]$ and apply $[\mathrm{T\text{-}LAM}]$.

\textit{Case} $[\mathrm{T\text{-}APP}]$: $t = f\, a$ with $\Gamma, x:S \vdash f :
\Pi(y:T_1).T_2$ and $\Gamma, x:S \vdash a : T_1$. By induction:
$\Gamma \vdash f[x:=s] : (\Pi(y:T_1).T_2)[x:=s]$ and
$\Gamma \vdash a[x:=s] : T_1[x:=s]$. Apply $[\mathrm{T\text{-}APP}]$.

\textit{Case} $[\mathrm{T\text{-}POP}]$: $t = \mathsf{pop}(t')$. By induction
$\Gamma \vdash t'[x:=s] : T'[x:=s]$. Apply $[\mathrm{T\text{-}POP}]$.

\textit{Cases} $[\mathrm{T\text{-}REFUSE}]$, $[\mathrm{T\text{-}BIND}]$,
$[\mathrm{T\text{-}COLLAPSE}]$: Analogous to the Pop case.
\end{proof}

\section{Full Progress and Preservation}

\begin{theorem}[Full Type Preservation]
\label{thm:preservation-full}
If $\Gamma \vdash t : T$ and $t \to t'$, then $\Gamma \vdash t' : T$.
\end{theorem}

\begin{proof}
By induction on the derivation of $\Gamma \vdash t : T$, with case analysis
on the reduction step $t \to t'$.

\textit{Case $[\mathrm{T\text{-}APP}]$ with $[\beta]$-reduction}:
$t = (\lambda x:T_1.b)\, a \to b[x:=a] = t'$.
From $[\mathrm{T\text{-}APP}]$: $\Gamma \vdash \lambda x:T_1.b : \Pi(x:T_1).T_2$
and $\Gamma \vdash a : T_1$.
From $[\mathrm{T\text{-}LAM}]$: $\Gamma, x:T_1 \vdash b : T_2$.
By Substitution (Lemma~\ref{lem:substitution}): $\Gamma \vdash b[x:=a] : T_2[x:=a] = T$.

\textit{Case $[\mathrm{T\text{-}POP}]$ with inner reduction}:
$t = \mathsf{pop}(t_0)$ with $\Gamma \vdash t_0 : T_0$ and $t_0 \to t_0'$.
By induction hypothesis, $\Gamma \vdash t_0' : T_0$.
Apply $[\mathrm{T\text{-}POP}]$ to obtain $\Gamma \vdash \mathsf{pop}(t_0') : \Adm(T_0)$.

\textit{Cases} $[\mathrm{T\text{-}REFUSE}]$, $[\mathrm{T\text{-}BIND}]$,
$[\mathrm{T\text{-}COLLAPSE}]$: Analogous, reducing the inner term.

\textit{Case $[\mathrm{T\text{-}COLLAPSE}]$ with inner Pop value}:
$t = \mathsf{collapse}(\mathsf{pop}(v), c)$ where $\mathsf{pop}(v)$ is a value
of type $\Adm(T')$. This reduces to a sealed observation value of type $\Col(T', c)$.
By $[\mathrm{T\text{-}COLLAPSE}]$, the type is preserved.
\end{proof}

\begin{theorem}[Full Progress]
\label{thm:progress-full}
If $\vdash t : T$ (well-typed in empty context), then either $t$ is a value,
or there exists $t'$ such that $t \to t'$.
\end{theorem}

\begin{proof}
By induction on the derivation $\vdash t : T$.

\textit{Case $[\mathrm{T\text{-}VAR}]$}: Impossible in empty context.

\textit{Case $[\mathrm{T\text{-}LAM}]$}: $t = \lambda x:T_1.b$ is a value.

\textit{Case $[\mathrm{T\text{-}APP}]$}: $t = f\, a$. By induction, $f$ is either
a value or can step. If $f$ can step to $f'$, then $f\, a \to f'\, a$.
If $f$ is a value, by Canonical Forms (Lemma~\ref{lem:canonical}) with type
$\Pi(x:T_1).T_2$, we have $f = \lambda x:T_1.b$. Now by induction on $a$:
if $a$ steps to $a'$, then $(\lambda x:T_1.b)\, a \to (\lambda x:T_1.b)\, a'$.
If $a$ is a value, then $(\lambda x:T_1.b)\, a \to b[x:=a]$ by $[\beta]$-reduction.

\textit{Case $[\mathrm{T\text{-}POP}]$}: $t = \mathsf{pop}(t_0)$. If $t_0$ is a
value, then $\mathsf{pop}(t_0)$ is a value. If $t_0 \to t_0'$, then
$\mathsf{pop}(t_0) \to \mathsf{pop}(t_0')$.

\textit{Case $[\mathrm{T\text{-}REFUSE}]$}: Analogous to Pop.

\textit{Case $[\mathrm{T\text{-}COLLAPSE}]$}: $t = \mathsf{collapse}(t_0, c)$
with $\vdash t_0 : \Adm(T')$. By induction on $t_0$: if $t_0$ is a value,
by Canonical Forms it is $\mathsf{pop}(v)$, and $\mathsf{collapse}(\mathsf{pop}(v), c)$
reduces to the sealed observation---a value. If $t_0$ steps, then $t$ steps
by reducing inside.

\textit{Case $[\mathrm{T\text{-}BIND}]$}: $t = \mathsf{bind}(t_1, t_2)$. By
induction, each component either steps or is a value. If any component steps,
$t$ steps. If both are values, $\mathsf{bind}(v_1, v_2)$ is a value.
\end{proof}

\begin{corollary}[Type Soundness]
If $\vdash t : T$, then evaluation of $t$ either produces a value of type $T$
or diverges. It does not get stuck.
\end{corollary}

\begin{proof}
By repeated application of Progress and Preservation. At each step, either $t$
is a value (done) or $t \to t'$ with $\vdash t' : T$ (continue). The process
cannot get stuck because Progress guarantees a step is always available if the
term is not yet a value.
\end{proof}

\section{Admissibility Preservation Theorem}

\begin{theorem}[Admissibility Preservation Across Reduction]
If $\Gamma \vdash t : \Adm(T)$ and $t \to^* v$ where $v$ is a value, then
$\vdash v : \Adm(T)$ and the history produced by evaluating $t$ contains no
refusal events that would invalidate the admissibility certificate.
\end{theorem}

\begin{proof}
By Full Preservation (Theorem~\ref{thm:preservation-full}), $\Adm(T)$ is
preserved throughout reduction. When $v$ is reached, it has type $\Adm(T)$
and by Canonical Forms, $v = \mathsf{pop}(v')$ with $\vdash v' : T$. The
admissibility type $\Adm(T)$ certifies that the path from $t$ to $v$ was
admissible: no $\Refuse$ event along the way was applied to the same symbol
as any $\Pop$ in the path, because such a conflict would create a type error
at the $\Pop$ introduction site.
\end{proof}

\chapter{Extended Type Examples and Applications}

\section{Typing File System Operations}

A concrete example grounds the type theory. File system operations provide a
natural domain because they involve admissibility (file must exist), refusal
(access denied), and collapse (reading produces an observable state).

\begin{example}[File Read as Typed Spherepop Term]
\begin{lstlisting}[style=rust, caption={File read in typed Spherepop}]
let read_file : FilePath -> Adm(FileContents) =
  fn path : FilePath.
    if file_exists(path) then
      pop(file_contents(path))   -- Adm(FileContents)
    else
      refuse(path, violation("FileNotFound"))
      -- Refused(FilePath, violation("FileNotFound"))

-- The return type is: Adm(FileContents) | Refused(FilePath, r)
-- The caller must handle both branches
\end{lstlisting}
The type signature forces the caller to acknowledge that reading can fail.
Unlike an exception system that hides the failure mode in runtime behavior,
the Spherepop type makes it explicit: the reason for refusal ($\mathsf{FileNotFound}$)
is part of the type, enabling callers to pattern-match on specific failure modes.
\end{example}

\section{Typing Memory Safety}

\begin{example}[Pointer Dereference with Admissibility Certificate]
\begin{lstlisting}[style=rust, caption={Safe pointer dereference}]
type Ptr(a : Type) = Adm(Reference(a))
-- A Ptr is an admissibility certificate that
-- the reference is valid and live.

let deref : Ptr(a) -> Adm(a) =
  fn p : Ptr(a).
    match p with
    | pop(ref) ->
        -- ref : Reference(a), known admissible
        pop(load(ref))    -- Adm(a)

let free_ptr : Ptr(a) -> Refused(Reference(a), "Deallocated") =
  fn p : Ptr(a).
    match p with
    | pop(ref) ->
        refuse(ref, explicit("Deallocated"))
-- After free_ptr, the ptr's type becomes Refused,
-- and any further deref is a type error:
-- deref requires Adm, not Refused.
\end{lstlisting}
Memory safety is enforced by types: after \texttt{free\_ptr}, the pointer's
type is $\Rfsd(\mathsf{Reference}(a), \mathsf{Deallocated})$, and
$\mathsf{deref}$ requires $\Adm(\mathsf{Reference}(a))$. A use-after-free is
a type error caught at compile time.
\end{example}

\section{Typing Cryptographic Operations}

\begin{example}[Hash as Collapse with Fixed Rule]
\begin{lstlisting}[style=rust, caption={Cryptographic hash as collapse}]
type SHA256Rule = CollapseRule::Proj("sha256")

let hash : Adm(Bytes) -> Collapsed(Bytes, SHA256Rule) =
  fn b : Adm(Bytes).
    collapse(b, SHA256Rule)
-- The output type Collapsed(Bytes, SHA256Rule) certifies:
-- (1) the input was admissible (had type Adm(Bytes)),
-- (2) the specific collapse rule SHA256 was applied,
-- (3) the observation is sealed: no further pops can unwrap it.

-- Preimage resistance is expressed by the absence of a term:
-- there is no well-typed term of type
--   Collapsed(Bytes, SHA256Rule) -> Adm(Bytes)
-- without additional information.
\end{lstlisting}
The irreversibility of cryptographic hashing is expressed by the type system:
$\Col(\mathsf{Bytes}, \mathsf{SHA256})$ has no inverse in the type theory
(since $[\mathrm{T\text{-}COLLAPSE}]$ has no elimination form that restores
$\Adm$). This makes preimage resistance a type-level property.
\end{example}


%% ============================================================
%% PART XVII: VARIATIONAL MECHANICS AND INFORMATION DYNAMICS
%% ============================================================
\part{Variational Mechanics and Information Dynamics}

\chapter{Histories as Paths Through Possibility Space}

\section{The Action Functional}

A Spherepop computation is a path through event space. Like a physical path in
classical mechanics, we can assign an action to each history and ask which
histories are stationary points of this action---the computational analogues of
geodesics.

\begin{definition}[Event Cost Function]
An \emph{event cost function} is a map $L : E \to \mathbb{R}_{\geq 0}$ assigning
a non-negative cost to each primitive event. Common choices:
\begin{align*}
L(\Pop(x)) &= c_{\mathrm{commit}} > 0 \quad\text{(cost of commitment)} \\
L(\Refuse(x, r)) &= c_{\mathrm{refuse}} \geq 0 \quad\text{(cost of refusal)} \\
L(\Bind(a, b)) &= c_{\mathrm{bind}} \geq 0 \quad\text{(cost of dependency)} \\
L(\Collapse(x, c)) &= c_{\mathrm{obs}} > 0 \quad\text{(cost of observation)}
\end{align*}
\end{definition}

\begin{definition}[History Action]
The \emph{action} of a history $H = [e_1, \ldots, e_n]$ is
\[
\mathcal{S}(H) = \sum_{i=1}^{n} L(e_i).
\]
\end{definition}

\begin{definition}[Admissible Path Space]
For boundary conditions $(H_0, H_f)$ (initial and final histories), the
\emph{admissible path space} is
\[
\mathcal{P}(H_0, H_f) = \{H \in \calH \mid H_0 \text{ is a prefix of } H,
H_f \text{ is a suffix of } H,\, \calA(H) \neq \varnothing\}.
\]
\end{definition}

\section{Discrete Euler-Lagrange Equations}

\begin{definition}[History Variation]
A \emph{variation} of history $H$ is a modified history $H + \delta H$ obtained
by inserting, deleting, or replacing a single event while keeping boundary
conditions fixed.
\end{definition}

\begin{definition}[Stationary History]
A history $H$ is \emph{stationary} in $\mathcal{P}(H_0, H_f)$ if for all
admissible variations $\delta H$,
\[
\mathcal{S}(H + \delta H) - \mathcal{S}(H) \geq 0.
\]
That is, $H$ is a local minimum of the action functional.
\end{definition}

\begin{theorem}[Discrete Euler-Lagrange Conditions]
A history $H = [e_1, \ldots, e_n]$ is stationary if and only if for each
interior event $e_k$ ($1 < k < n$), replacing $e_k$ with any admissible
event $e'_k$ satisfies $L(e'_k) \geq L(e_k)$.
\end{theorem}

\begin{proof}
In the discrete setting, a variation that replaces $e_k$ with $e'_k$ changes
the action by $L(e'_k) - L(e_k)$. Stationarity requires this to be non-negative
for all admissible $e'_k$, which means $e_k$ is the minimum-cost admissible event
at position $k$. This is the discrete analogue of the Euler-Lagrange equation:
the optimal trajectory selects, at each step, the event that minimizes the local
cost subject to the admissibility constraint.
\end{proof}

\begin{remark}[Connection to Admissibility Geometry]
The admissibility constraint $\calA(H) \neq \varnothing$ plays the role of the
geometric constraint in constrained variational mechanics. A history that exhausts
all admissible options ($\calA(H) = \varnothing$) is a terminal trajectory---the
computation has committed to all available possibilities. The xylomorphic criterion
$\lambda < 1$ from the RSVP framework is the continuum version of this constraint:
a trajectory remains in the admissible region as long as the local admissibility
curvature $\kappa_A(x) > \lambda$.
\end{remark}

\section{Least-Action Computation}

\begin{theorem}[Least-Action Principle for Computation]
A Spherepop computation that minimizes the total action $\mathcal{S}(H)$ over
all admissible paths between boundary conditions performs the minimum-cost
sequence of commitments and observations needed to reach the terminal state.
\end{theorem}

\begin{proof}
By the Euler-Lagrange conditions, minimizing $\mathcal{S}(H)$ requires selecting
the minimum-cost admissible event at each step. This is precisely greedy optimal
computation: choose the cheapest admissible operation at each decision point.
The resulting history is the discrete geodesic in event space connecting the
boundary conditions.
\end{proof}

\begin{example}[Parsing as Least-Action Computation]
A parser that recognizes a string against a grammar can be viewed as a
Spherepop computation minimizing the action. Each match of a terminal symbol
is a Pop event; each failed match is a Refuse event. The parser's job is to
find the minimum-action path through the grammar's event space that accepts
the input string. The optimal parser (in the sense of minimal backtracking)
corresponds to the least-action history. PEG parsers implement this by trying
alternatives in order and committing (Pop) on the first match, rather than
exploring all possibilities.
\end{example}

\chapter{Information-Theoretic Properties of Histories}

\section{History Entropy}

\begin{definition}[History Distribution]
A \emph{history distribution} over event alphabet $E$ is a probability
distribution $\mu$ over $\calH$ such that $\mu(H)$ represents the probability
of a specific execution path.
\end{definition}

\begin{definition}[History Entropy]
The \emph{Shannon entropy} of a history distribution $\mu$ is
\[
\mathbb{H}[\mu] = -\sum_{H \in \calH} \mu(H) \log \mu(H).
\]
\end{definition}

\begin{theorem}[Commitment Reduces Entropy]
Each Pop event strictly reduces the history entropy: if $\mu'$ is the
distribution over extensions of $H$ after $\Pop(x)$, then
$\mathbb{H}[\mu'] < \mathbb{H}[\mu]$.
\end{theorem}

\begin{proof}
The Pop event $\Pop(x)$ eliminates all histories in $\mu$ that do not have
$x$ committed. The surviving histories are conditioned on $x$ having been
committed, reducing the support of the distribution. A distribution over a
strictly smaller support has strictly lower entropy (since the eliminated
histories had positive probability under $\mu$).
\end{proof}

\begin{theorem}[Refusal Does Not Reduce History Entropy]
A Refuse event $\Refuse(x, r)$ does not reduce the entropy of the history
distribution: the option space $\Omega$ is unchanged, and the support of
the future-history distribution is unchanged (any future history compatible
with $\calA(H')$ was already compatible with $\calA(H) \setminus \{x\}$).
\end{theorem}

\begin{proof}
$\Refuse(x, r)$ appends an event to the history (increasing its complexity)
but does not remove $x$ from $\Omega$. Future histories can still attempt
$\Pop(x)$ (it will be caught by the admissibility check), but they can also
avoid $x$ entirely. The future history space is therefore unchanged in size,
and entropy is preserved.
\end{proof}

\begin{remark}[The Commitment-Refusal Asymmetry]
Theorem 36.2 and 36.3 together formalize the fundamental asymmetry between
Pop and Refuse. Pop is entropy-reducing: it commits to a possibility and
forecloses alternatives, shrinking the future history space. Refuse is
entropy-neutral: it documents that a path was inadmissible without shrinking
the space of future histories. This asymmetry is why refusals are cheaper
to make than commitments and why documented refusal (rather than silent
avoidance) is a genuine contribution to the history record.
\end{remark}

\section{Hamming Distance on Histories}

\begin{definition}[History Edit Distance]
The \emph{edit distance} between histories $H_1$ and $H_2$ is the minimum
number of event insertions, deletions, and substitutions needed to transform
$H_1$ into $H_2$:
\[
d_{\mathrm{edit}}(H_1, H_2) = \min\{k \mid H_1 \xrightarrow{k\text{ edits}} H_2\}.
\]
\end{definition}

\begin{definition}[History Hamming Distance]
For histories of equal length $n$, the \emph{Hamming distance} is
\[
d_H(H_1, H_2) = |\{i \mid e_{1,i} \neq e_{2,i}\}|.
\]
\end{definition}

\begin{theorem}[Error-Correcting Bounds for Histories]
A set of histories $\mathcal{C} \subseteq \calH^n$ (a history code) can
correct up to $t$ event errors (substitutions) if and only if
\[
\min_{H_1 \neq H_2 \in \mathcal{C}} d_H(H_1, H_2) \geq 2t + 1.
\]
\end{theorem}

\begin{proof}
This is the standard Hamming bound applied to the history alphabet. If the
minimum pairwise Hamming distance is $2t+1$, then any two codeword-histories
differ in at least $2t+1$ positions. An error pattern affecting at most $t$
positions cannot move any codeword to within Hamming distance $t$ of any
other codeword, so the nearest codeword is always the correct one.
\end{proof}

\begin{remark}[Connection to Error Correction Chapter]
The history error-correcting bound connects the abstract history theory to the
Error Correction as Certified Refusal chapter of Part~VI. A set of redundant
histories (encoding the same computation multiple times) forms a history code.
Event corruption (bit errors in the history log) can be corrected if the
minimum Hamming distance of the code exceeds twice the error rate. The certified
refusal system then serves as the decoder: when a Refuse event indicates that
a committed symbol's certificate does not match the redundant record, the
decoder identifies and corrects the corrupted event.
\end{remark}

\section{Kolmogorov Complexity and History Compression}

\begin{definition}[Kolmogorov Complexity of a History]
The \emph{Kolmogorov complexity} of a history $H$ is
\[
K(H) = \min\{|p| \mid U(p) = H\}
\]
where $U$ is a universal Turing machine and $|p|$ is the length of the shortest
program $p$ that outputs $H$.
\end{definition}

\begin{proposition}[Incompressible Histories Exist]
For each length $n$, there exist histories $H$ of length $n$ with $K(H) \geq n$.
These are \emph{random histories} whose event sequences contain no exploitable
regularity.
\end{proposition}

\begin{proof}
There are $|E|^n$ histories of length $n$ but at most $2^{n-1}$ programs of
length less than $n$ bits. By counting, most histories cannot be compressed.
\end{proof}

\begin{theorem}[Generative Notes Reduce Kolmogorov Complexity]
A generative note $N$ that enables a new inferential framework reduces the
Kolmogorov complexity of the histories generated within that framework:
$K(H \mid N) < K(H)$ for histories $H$ expressible concisely within $N$.
\end{theorem}

\begin{proof}
Given $N$ as a reference, the shortest program for $H$ can invoke $N$ as a
subroutine. If $N$ provides a compression for the structure of $H$ (for
example, $H$ is a theorem in a formal system captured by $N$), then
$|p| < K(H)$ where $p$ invokes $N$ and provides only the theorem-specific
details. Hence $K(H \mid N) \leq |p| < K(H)$.
\end{proof}

\begin{remark}
This theorem gives a Kolmogorov complexity account of why generative notes
have high reachability leverage. A programming language, a mathematical theory,
or a grammar compresses the Kolmogorov complexity of all programs, theorems,
or strings expressible within it. The note is a reference object that reduces
the minimum description length of the histories it enables.
\end{remark}


%% ============================================================
%% PART XVIII: APPLICATIONS AND CONNECTIONS
%% ============================================================
\part{Applications and Connections to Existing Frameworks}

\chapter{Event Sourcing and Database Theory}

\section{Event Sourcing as Spherepop}

Event sourcing~\cite{fowler05} is an architectural pattern where application
state is derived from an append-only log of events rather than a mutable state
store. Spherepop provides a formal vocabulary in which event sourcing can be
described precisely, though it should be read as an architectural analogue or
instance of the history-first pattern rather than a claim that ordinary
event-sourced systems already implement the full calculus: they do not, in
general, implement Refuse with reasons, Bind, admissibility typing, or attested
Collapse.

\begin{proposition}[Event Sourcing as an Instance of the History-First Pattern]
An event-sourced application using an append-only event log can be modeled as
a Spherepop world $(H, \Omega)$ with the LastWrite collapse rule $c_{LW}$
applied to derive the current application state, provided the additional
Spherepop structure (documented refusal, dependency binding, admissibility
typing, attested Collapse) is added on top of the bare log.
\end{proposition}

\begin{proof}[Sketch]
The event log is the history $H$: each event appended to the log corresponds to
a Pop event in the Spherepop world. The current application state, in the
ordinary event-sourcing sense, is $\operatorname{View}_{c_{LW}}(H) = c_{LW}(H)$:
the most recent value for each entity, derived by \emph{viewing} the history
through the LastWrite rule---this is a VIEW operation, not by itself an
attested Collapse, since a read model is not thereby adopted as constraining
future computation. The event sourcing invariant---that the current state can
always be reconstructed from the log---parallels Spherepop's Replay Equivalence
(Theorem~10.3): $I(H) = V(C(H))$, meaning the interpreter and VM produce
identical histories and hence identical LastWrite views.

Event sourcing's ``projection'' (computing a read model from the event log) and
CQRS read models are likewise VIEW operations: a specific collapse rule $c$
applied to the history to produce $c(H)$. They become Collapse events only if
and when their outputs are attested and accepted as constraints binding later
computation---for instance, a read model cached and then relied upon by a
downstream process that treats it as authoritative would, at that point, need
an attested Collapse to record the dependency, rather than remaining a
disposable, re-derivable VIEW.
\end{proof}

\begin{example}[Bank Account as Spherepop World]
A bank account event log in event sourcing corresponds to the following Spherepop
world:
\begin{lstlisting}[style=rust, caption={Bank account as Spherepop world}]
// Each transaction is a Pop event
let history = vec![
    Event::Pop(Symbol::new("deposit:100")),
    Event::Pop(Symbol::new("withdraw:30")),
    Event::Pop(Symbol::new("deposit:50")),
    Event::Refuse(Symbol::new("withdraw:200"),
                  RefusalReason::ConstraintViolation(
                    "InsufficientFunds")),
];

// Balance is the Accumulate collapse
// deposits positive, withdrawals negative
let balance = apply_collapse(&history, &acc_rule);
// balance = 120
\end{lstlisting}
The refusal event for the insufficient-funds withdrawal is preserved in the
history with its reason, making the audit trail complete. A state-centric
database would only show the final balance; Spherepop shows every decision.
\end{example}

\section{Provenance Semirings}

Provenance semirings~\cite{shapiro11} are algebraic structures that track the
derivation of database query results from input tuples. They generalize boolean
(does this tuple appear?), counting (how many times?), and probability
(with what likelihood?) provenance.

\begin{theorem}[Collapse Rules Are Provenance Semirings]
Each Spherepop collapse rule $c : \calH \to O_c$ induces a provenance semiring
$(O_c, \oplus, \otimes)$ where:
\begin{align*}
o_1 \oplus o_2 &= c(H_1 \sqcup H_2) \quad\text{(union provenance)} \\
o_1 \otimes o_2 &= c(H_1 \mathbin{++} H_2) \quad\text{(join provenance)}
\end{align*}
\end{theorem}

\begin{proof}
The Accumulate rule $c_{\mathrm{acc}}$ gives the counting semiring
$(\mathbb{N}, +, \times)$: the count of how many derivations produced each
fact is the sum of counts from alternative derivations (union) and the product
of counts from joined derivations. The LastWrite rule gives the boolean semiring
$(\{0,1\}, \vee, \wedge)$: a fact holds if any derivation produces it (union is
OR) and if all required sub-derivations succeed (join is AND). Different collapse
rules give different provenance semantics, all arising from the same Spherepop
history.
\end{proof}

\section{The Differential Privacy Connection}

Differential privacy requires that no single individual's data significantly
affects the observable output of a computation. In Spherepop terms, this is
a constraint on the collapse rule's sensitivity to individual Pop events.

\begin{definition}[$\varepsilon$-Differential Privacy in Spherepop]
A collapse rule $c$ satisfies \emph{$\varepsilon$-differential privacy} if
for any two histories $H$ and $H'$ differing in exactly one Pop event,
and for any observable set $O \subseteq O_c$:
\[
P(c(H) \in O) \leq e^\varepsilon \cdot P(c(H') \in O).
\]
\end{definition}

\begin{proposition}[Noisy Accumulate Is Differentially Private]
The Accumulate rule with Laplace noise $\mathsf{Lap}(1/\varepsilon)$ added to
each count satisfies $\varepsilon$-differential privacy.
\end{proposition}

\begin{proof}
The sensitivity of $c_{\mathrm{acc}}$ to a single Pop event is 1 (one event
changes exactly one count by exactly 1). Adding $\mathsf{Lap}(1/\varepsilon)$
noise calibrated to this sensitivity produces a differentially private mechanism
by the standard Laplace mechanism analysis.
\end{proof}

\chapter{Programming Language Theory Connections}

\section{Denotational Semantics as VIEW}

Denotational semantics assigns mathematical meanings (denotations) to programs.
In Spherepop's vocabulary, denotational semantics corresponds to applying a
collapse rule as a pure \emph{view} of the operational history---semantic
folding in the sense of Section~\ref{sec:layered-world-model}---and not to a
committed Collapse. A denotation does not ordinarily mutate the program's
history merely by being calculated; computing $\llbracket P \rrbracket$ is a
read, not an attested adoption of that meaning as a constraint on later
computation. The functorial structure below therefore concerns the map
$F_c : \Hist \to \Obs_c$ as a VIEW-level construction; it should not be read as
implying that denotation semantics themselves collapse.

\begin{theorem}[Denotational Semantics Corresponds to a Collapse-Rule View]
Every compositional denotational semantics $\llbracket \cdot \rrbracket :
\mathsf{Terms} \to \mathsf{Meanings}$ is a collapse functor
$F_c : \Hist \to \Obs_c$ for an appropriate collapse rule $c$.
\end{theorem}

\begin{proof}
Compositionality requires that $\llbracket P \parallel Q \rrbracket =
\llbracket P \rrbracket \oplus \llbracket Q \rrbracket$ for some meaning
composition operator $\oplus$. This is precisely the functor composition law
$F_c(H_1 \mathbin{++} H_2) = F_c(H_1) \circ F_c(H_2)$ of the collapse functor.
The denotation is the observable state produced by applying $c$ to the
operational history. Different denotational semantics correspond to different
collapse rules: the trace semantics uses the Identity rule; the collecting
semantics uses the Accumulate rule; the big-step semantics uses the LastWrite
rule.
\end{proof}

\section{Hoare Logic as Admissibility Typing}

Hoare logic~\cite{hoare85} provides pre- and postconditions for program
statements: $\{P\}\, C\, \{Q\}$ asserts that if $P$ holds before executing
$C$, then $Q$ holds after. In Spherepop, Hoare triples are admissibility typing
judgements.

\begin{theorem}[Hoare Triple as Spherepop Typing]
The Hoare triple $\{P\}\, C\, \{Q\}$ is equivalent to the Spherepop typing
judgement
\[
P \vdash C : \Adm(Q)
\]
where $P$ is a predicate on admissibility sets and $Q$ is the postcondition
certified by the admissibility type.
\end{theorem}

\begin{proof}
The precondition $P$ specifies the admissibility set before $C$ executes:
$\calA(H) \supseteq \mathrm{Required}(P)$. The program $C$ corresponds to a
sequence of Spherepop operations that, when typed in context $P$, produces a
result of type $\Adm(Q)$. The admissibility certificate $\Adm(Q)$ certifies
that the postcondition $Q$ holds after execution. The standard Hoare logic rules
(consequence, composition, conditional, loop) correspond to the standard
Spherepop typing rules under this identification.
\end{proof}

\section{Separation Logic and Bind}

Separation logic extends Hoare logic with a spatial connective: $P * Q$ asserts
that $P$ and $Q$ hold over disjoint portions of memory. In Spherepop, this
corresponds to independent Bind structures over disjoint option spaces.

\begin{theorem}[Separating Conjunction as Disjoint Bind]
The separating conjunction $P * Q$ in separation logic corresponds to the
Spherepop assertion that $P$ and $Q$ hold over disjoint option spaces
$\Omega_P \cap \Omega_Q = \varnothing$, with no Bind events coupling $\Omega_P$
and $\Omega_Q$.
\end{theorem}

\begin{proof}
The frame rule of separation logic---$\{P\}\, C\, \{Q\}$ implies
$\{P * R\}\, C\, \{Q * R\}$ for any $R$ not modified by $C$---corresponds
to the Spherepop property that operations on $\Omega_C$ do not affect the
admissibility certificates over $\Omega_R$ when $\Omega_C \cap \Omega_R = \varnothing$
and no Bind events cross the boundary. The spatial separation in memory
corresponds to option-space separation in Spherepop.
\end{proof}

\chapter{Connections to Cognitive Science and MEM\textbar8}

\section{The Spherepop Model of Memory Retrieval}

The MEM\textbar8 architecture~\cite{chenoweth25} treats memory as wave-based
pattern activation rather than symbol lookup. This section develops the formal
connection between Spherepop histories and MEM\textbar8 memory dynamics.

\begin{definition}[Memory as History Distribution]
In the MEM\textbar8 framework, a memory trace corresponds to a history
distribution $\mu_M$ over the history space $\calH$: the distribution represents
the set of histories compatible with the stored memory and their relative
activation strengths.
\end{definition}

\begin{theorem}[Memory Retrieval as Conditional Collapse]
Memory retrieval in MEM\textbar8 corresponds to conditioning the history
distribution $\mu_M$ on an observational cue $o$:
\[
\mu_{\mathrm{retrieved}} = \mu_M(\cdot \mid c^{-1}(o))
\]
where $c$ is the collapse rule determining what aspects of the memory are
visible to the cue.
\end{theorem}

\begin{proof}
MEM\textbar8 retrieval works by pattern completion: a partial cue activates the
wave patterns associated with the complete memory. In Spherepop terms, the cue
$o$ specifies an observable state, and retrieval reconstructs the most probable
history in the fibre $c^{-1}(o)$. This is Bayesian inference over the history
distribution conditioned on the cue observation---precisely the conditional
collapse $\mu_M(\cdot \mid c^{-1}(o))$.
\end{proof}

\begin{example}[Episodic Memory as Identity-Rule Collapse]
Episodic memory (memory for specific events, including their temporal and
contextual details) corresponds to near-Identity collapse: the memory trace
preserves the full event sequence $H$, and retrieval under Identity rule $c_I$
returns the original history. The wave-based pattern in MEM\textbar8 stores
the temporal ordering and contextual bindings of events---equivalent to preserving
$H$ up to small perturbations.
\end{example}

\begin{example}[Semantic Memory as Accumulate-Rule Collapse]
Semantic memory (knowledge about concepts and categories, abstracted from
specific episodes) corresponds to Accumulate collapse: the memory trace
stores the count of how many times each concept was encountered, losing the
specific episodic context. MEM\textbar8's categorical wave patterns represent
this statistical summary across episodes.
\end{example}

\section{Forgetting as Controlled Collapse}

\begin{definition}[Forgetting Rule]
A \emph{forgetting rule} $c_{\mathrm{forget}}(t)$ is a time-parameterized
collapse rule where the coarseness of the collapse increases with time $t$:
for $t_1 < t_2$, $c_{\mathrm{forget}}(t_1)$ is finer than
$c_{\mathrm{forget}}(t_2)$.
\end{definition}

\begin{theorem}[Ebbinghaus Forgetting as Entropy Growth]
Under a forgetting rule $c_{\mathrm{forget}}(t)$, the observational entropy
$S_{c(t)}(o)$ of any fixed observation $o$ is non-decreasing in $t$:
\[
t_1 < t_2 \implies S_{c(t_1)}(o) \leq S_{c(t_2)}(o).
\]
\end{theorem}

\begin{proof}
By Entropy Monotonicity Under Rule Refinement (Theorem of Part~IX): since
$c_{\mathrm{forget}}(t_2)$ is coarser than $c_{\mathrm{forget}}(t_1)$, the
fibre of $c_{\mathrm{forget}}(t_2)$ is larger, hence its log is larger.
Forgetting coarsens the collapse rule, increasing the number of histories
consistent with the retained observation, which is precisely what the
Ebbinghaus forgetting curve describes: more historical detail is lost, more
alternative histories become consistent with the remaining memory trace.
\end{proof}

\section{Priming as Pre-emptive Admissibility Adjustment}

\begin{definition}[Priming Event]
A \emph{priming event} is a Refuse event that marks certain continuations as
temporarily less admissible and a Bind event that increases the admissibility
curvature of related continuations.
\end{definition}

\begin{theorem}[Semantic Priming as Admissibility Field Reshaping]
Semantic priming (the facilitated retrieval of concepts related to a recently
activated concept) corresponds to a temporary reshaping of the admissibility
field $\kappa_A$ in the direction of the primed concept: symbols related to
the prime have increased $\kappa_A$ (higher admissibility) while unrelated
symbols are temporarily Refused.
\end{theorem}

\begin{proof}
Priming in cognitive science increases the activation of related concepts by
lowering their retrieval threshold. In Spherepop, this corresponds to
temporarily increasing $\kappa_A(x)$ for symbols $x$ related to the prime
(making them more admissible, i.e., lower-cost to Pop), and decreasing
$\kappa_A(y)$ for unrelated symbols (making them relatively less admissible).
The decay of the priming effect over time corresponds to $\kappa_A$ returning
to its baseline, modeled by the forgetting rule.
\end{proof}

\chapter{Connections to Physics and Complex Systems}

\section{Thermodynamic Interpretation of Collapse}

The relationship between information-theoretic entropy and thermodynamic entropy
runs deep. The Spherepop collapse framework provides a computational analogue
of thermodynamic processes.

\begin{theorem}[Second Law as Entropy Non-Decrease Under Coarsening]
Under any physical measurement (coarsening of observational framework from
$c_1$ to $c_2$ with $c_2$ coarser), the observational entropy of any history
does not decrease:
\[
S_{c_2}(o) \geq S_{c_1}(o').
\]
This is the computational Second Law: every measurement loses information.
\end{theorem}

\begin{proof}
By Entropy Monotonicity Under Rule Refinement (Part~IX Chapter~22): coarsening
increases fibre size and hence entropy. Physical measurement is a collapse
(projecting the quantum state onto a classical observable), and each measurement
can only coarsen the observational framework.
\end{proof}

\begin{remark}[Maxwell's Demon as Refusal Without Cost]
Maxwell's Demon in thermodynamics is an entity that observes individual molecules
and opens or closes a gate based on their velocities, apparently violating the
Second Law. In Spherepop terms, the Demon performs Refuse operations: marking
slow molecules as inadmissible for passage through the gate. The resolution
(Landauer's principle) is that the Demon's memory erasure has thermodynamic
cost. In Spherepop, this corresponds to GC: erasing the Demon's history of
observations (applying Certificate-Only GC to the Refuse records) has a minimum
computational cost of $k_B T \ln 2$ per bit erased. The Refuse events cannot
be erased for free.
\end{remark}

\section{Phase Transitions in Note Infrastructure}

\begin{definition}[Note Density]
The \emph{note density} of a civilization at time $t$ is
\[
\rho_N(t) = \frac{|R(Z(t))|}{|\calC_{\mathrm{total}}|}
\]
the fraction of all possible continuations that are reachable given the
civilization's note infrastructure $\calN(t)$.
\end{definition}

\begin{theorem}[Percolation Threshold in Note Networks]
Note infrastructure undergoes a percolation phase transition: for note density
$\rho_N$ below a critical threshold $\rho_c$, the reachability graph of the
civilization's knowledge is fragmented into isolated components. Above $\rho_c$,
a giant connected component emerges, and most knowledge becomes mutually reachable.
\end{theorem}

\begin{proof}[Proof sketch]
The reachability graph has nodes for continuations and edges for notes that
enable traversal between them. As note density increases from $0$ to $1$, this
random graph undergoes the Erd\H{o}s-R\'enyi percolation transition: for
$\rho_N < \rho_c$, the largest connected component has size $O(\log n)$; for
$\rho_N > \rho_c$, it has size $\Theta(n)$. The critical density $\rho_c$
depends on the structure of the continuation space.
\end{proof}

\begin{remark}
The invention of writing corresponds to a transition through $\rho_c$ in human
civilizational history: before writing, knowledge was stored in individual
memories with high loss rates (note density below threshold). Writing created
persistent notes with much lower loss rates, enabling the civilization's
reachability graph to transition into its giant-component phase. The emergence
of the internet represents another such transition, this time increasing not
just persistence but connectivity between note systems.
\end{remark}


%% ============================================================
%% PART XIX: PHILOSOPHICAL FOUNDATIONS
%% ============================================================
\part{Philosophical Foundations}

\chapter{Ontology of Process Versus Substance}

\section{The Substance Tradition}

Western metaphysics from Aristotle through Kant has predominantly held that
the primary entities of reality are \emph{substances}: things that persist
through time and bear properties. On this view, a rock, a person, a number,
a program are all substances---entities that exist, that have attributes, and
whose identity persists across change.

The substance tradition generates characteristic puzzles. What makes a ship
that has had all its planks replaced the same ship (the Ship of Theseus)?
What makes a person who has changed all their beliefs the same person? What
makes a program that has been refactored the same program? These puzzles arise
because the substance view treats identity as primitive and change as something
that happens to an otherwise stable entity.

The Spherepop ontology inverts this. Histories are primary; substances are
derived. Change is not something that happens to an entity; entities are
persistent patterns in a history of events. The puzzles of identity through
change dissolve: the ship is identified by its history, not by its material
composition at any moment. The person is identified by their developmental
history, not by their current beliefs. The program is identified by its
development history (the Git repository), not by its current source text.

\section{Process Philosophy Formalized}

Process philosophy, developed by Whitehead and later by Rescher, holds that
processes rather than substances are the fundamental entities of reality.
Spherepop provides a formal framework for this tradition.

\begin{definition}[Process]
In the Spherepop ontology, a \emph{process} is a history $H \in \calH$:
a finite sequence of events that constitutes a trajectory through possibility
space. Processes are the fundamental entities. Substances are derived.
\end{definition}

\begin{definition}[Substance as Persistent Process]
A \emph{substance} is an equivalence class $[H]_c$ under a stable collapse
rule $c$: a family of histories whose observable state $c(H)$ remains constant
over admissible extensions. Substances are persistent processes---processes
whose observable character does not change as they evolve.
\end{definition}

\begin{theorem}[Process-Substance Duality]
There is a categorical duality between the category $\Hist$ (processes as
morphisms) and the category $\Obs_c$ (substances as objects): the collapse
functor $F_c : \Hist \to \Obs_c$ exhibits $\Obs_c$ as the category of
stable orbits of $\Hist$ under the equivalence $\sim_c$.
\end{theorem}

\begin{proof}
Objects of $\Obs_c$ are equivalence classes $[H]_c$---precisely the persistent
processes (substances). Morphisms of $\Obs_c$ are observable transitions between
these classes. The functor $F_c$ maps each process (morphism of $\Hist$) to its
observable substance (object of $\Obs_c$). This is the formal statement of the
Noun Fallacy: objects of $\Obs_c$ appear to be the primary entities, but they
are images of morphisms of $\Hist$ under $F_c$.
\end{proof}

\section{Identity Through Change}

\begin{theorem}[Ship of Theseus Resolution]
In the Spherepop ontology, the Ship of Theseus has a determinate identity at
every point in its history. Two configurations of the ship are the same ship
if and only if they appear in the same continuous history.
\end{theorem}

\begin{proof}
The ship is identified by its history $H$, not by its material composition
at any moment. Two configurations $H_t$ and $H_{t'}$ are configurations of
the same ship if and only if $H_t$ is a prefix of $H_{t'}$ (or vice versa)
in the continuous developmental history. The gradual replacement of planks
is a sequence of events in the history; the ship's identity is the continuous
history, not any particular state $c(H_t)$ at a moment.
\end{proof}

\chapter{Language, Meaning, and Semantic Collapse}

\section{Meaning as Reachability}

The semantics of natural language has been understood primarily in terms of
truth conditions (what makes a sentence true) or use conditions (what role
a sentence plays in a practice). The Spherepop framework suggests a third
account: meaning as reachability.

\begin{definition}[Meaning as Continuation Set]
The \emph{meaning} of a linguistic expression $E$ is the set of continuations
$\calC(E)$ that $E$ makes reachable for a competent interpreter:
\[
\llbracket E \rrbracket = \calC(E) = \{c \in \calC \mid P_A(c, t \mid E) > P_A(c, t)\}.
\]
\end{definition}

\begin{theorem}[Synonymy as Continuation Equivalence]
Two expressions $E_1$ and $E_2$ are synonymous if and only if
$\calC(E_1) = \calC(E_2)$: they make the same set of continuations reachable.
\end{theorem}

\begin{proof}
If $\calC(E_1) = \calC(E_2)$, then for any agent $A$ and continuation $c$,
$P_A(c, t \mid E_1) = P_A(c, t \mid E_2)$. The two expressions have identical
effects on the reachability of all future actions, which is the strongest sense
in which two expressions can mean the same thing.
\end{proof}

\begin{theorem}[Semantic Drift as Collapse Rule Change]
Semantic change over time---a word acquiring a new meaning or losing an old
one---corresponds to a change in the collapse rule applied to the word's
continuation set.
\end{theorem}

\begin{proof}
A word $w$ at time $t_1$ has meaning $\llbracket w \rrbracket_{c_1} = F_{c_1}(H_w)$
where $H_w$ is the history of uses of $w$ and $c_1$ is the community's collapse
rule for interpreting uses. At time $t_2 > t_1$, the community may apply a
different rule $c_2$ to the same history $H_w$, producing
$\llbracket w \rrbracket_{c_2} = F_{c_2}(H_w)$. The word's meaning has changed
not because the word itself changed but because the observational framework for
interpreting it changed.
\end{proof}

\section{Translation as Approximate Adjoint}

\begin{theorem}[Translation as Partial Adjoint]
Translation between languages $L_1$ and $L_2$ is a partial adjoint between
the semantic categories $\mathsf{Sem}(L_1)$ and $\mathsf{Sem}(L_2)$: it
preserves continuation sets as closely as possible given the structural
differences between the languages.
\end{theorem}

\begin{proof}
A translation $T_{12} : L_1 \to L_2$ maps expressions in $L_1$ to expressions
in $L_2$. It is an approximate right adjoint to the meaning functor if
$\calC(T_{12}(E)) \approx \calC(E)$ for all $E \in L_1$. Perfect translation
would require $\calC(T_{12}(E)) = \calC(E)$ exactly, which fails when $L_2$
lacks the grammatical resources to express the continuation set of some
$E \in L_1$ (a translation gap). The quality of a translation is measured by
the closeness of the induced adjunction.
\end{proof}

\chapter{Ethics, Institutions, and Civilizational Responsibility}

\section{Moral Obligations as Coordination Notes}

\begin{theorem}[Moral Norms as Coordination Notes]
Moral norms, insofar as they enable or constrain coordination between agents,
are coordination notes in the Spherepop sense: they bind the continuation
spaces of multiple agents, creating coupled trajectories from what would
otherwise be independent evolutions.
\end{theorem}

\begin{proof}
A moral norm $N$ that prohibits action $a$ (such as theft or deception) is a
Refuse event applied to all agents: it marks action $a$ as inadmissible with
reason $\mathsf{MoralProhibition}(N)$. A moral norm that requires action $a$
(such as keeping promises) is a Bind event: it couples the agent's current
commitment to the future execution of $a$. Both types of norm are coordination
notes: they synchronize agent trajectories, enabling cooperation and joint
reachability that would be impossible without the shared normative framework.
\end{proof}

\begin{remark}
The formalization of moral norms as Spherepop operations has implications for
moral epistemology. The claim that a norm is valid is the claim that the
corresponding coordination note is genuine: that it actually increases collective
reachability for the community that adopts it. Evaluating moral norms by their
reachability consequences---rather than by their intrinsic properties or their
conformity to ideal principles---is a version of moral consequentialism expressed
in the Spherepop framework.
\end{remark}

\section{Institutional Failure as Note Collapse}

\begin{definition}[Institutional Integrity]
An institution has \emph{integrity} with respect to its mandate $M$ if its
history $H_{\mathrm{inst}}$ preserves the admissibility certificates needed
to continue executing $M$: $\calA(H_{\mathrm{inst}}) \supseteq
\mathrm{Required}(M)$.
\end{definition}

\begin{theorem}[Institutional Failure as Admissibility Collapse]
An institution fails when its history accumulates refusal events that make
its mandate inadmissible: $\calA(H_{\mathrm{inst}}) \cap \mathrm{Required}(M)
= \varnothing$. Institutional reform is a repair morphism that restores
$\calA$ to include $\mathrm{Required}(M)$.
\end{theorem}

\begin{proof}
By Refusal Monotonicity, the admissibility region $\calA(H)$ can only shrink
as the institution's history grows. If the history accumulates refusals of
actions required by the mandate $M$, then eventually $\calA(H) \cap
\mathrm{Required}(M) = \varnothing$: the institution can no longer execute
its mandate. This is institutional failure. Reform corresponds to a repair
morphism that replaces the inadmissible history with an alternative history
$H'$ in which the refused actions are no longer refused---restoring the
admissibility of the mandate.
\end{proof}

\begin{example}[Democratic Institutions]
A democratic institution's mandate includes the requirement that elections be
held, results certified, and power transferred peacefully. If the institution's
history accumulates refusals of these requirements (e.g., failure to certify
elections, refusal to transfer power), then $\calA(H) \cap \{\mathsf{Election},
\mathsf{Certification}, \mathsf{Transfer}\} = \varnothing$, and the institution
has failed its democratic mandate. Restoration requires a repair morphism---new
elections, constitutional reform, external oversight---that restores these
requirements to the admissibility set.
\end{example}


%% ============================================================
%% PART XX: IMPLEMENTATION ARCHITECTURE
%% ============================================================
\part{Implementation Architecture}

\chapter{The Spherepop Interpreter in Full}

\section{Architecture Overview}

The Spherepop interpreter is organized around the principle that the World
structure---containing history and option space---is the primary runtime
object. Observable state is never stored; it is always derived on demand.
This section documents the full module architecture and its relationship to
the theoretical results developed in preceding parts.

The single \texttt{observe} method of earlier drafts conflated the pure lookup
that the formal $[\mathrm{VIEW}]$ judgment describes with the committed,
event-appending act that the formal $[\mathrm{COLLAPSE}]$ judgment describes---an
inconsistency the operational semantics never actually had (Section
\ref{sec:view-vs-collapse} draws the formal line; the code below now draws the
same line). The interface is split into two methods accordingly. \texttt{view}
is pure. \texttt{collapse} verifies that the supplied observation matches the
value the world would currently derive, checks admissibility, and only then
appends a Collapse event.

\begin{lstlisting}[style=rust, caption={Complete World structure with view/collapse split}]
pub struct World {
    pub history: History,
    pub options: OptionSpace,
    // Operational state (options, admissibility index, dependency graph)
    // is the primary stored structure. Semantic/observable state is never
    // stored; it is always derived on demand via `view`.
}

/// Records why a quotient was adopted, so replay can reconstruct the
/// reasoning, not just the resulting value. A bare (x, observation) pair
/// is insufficient if rule definitions can change over the program's
/// lifetime, so the rule itself is identified by a stable id plus version
/// (or, more strongly, a content-addressed rule specification).
pub struct Provenance {
    pub rule_id: RuleId,
    pub rule_version: u64,
    pub source: ProvenanceSource, // e.g. Computed, ExternalInput, ReplayedFrom(EventId)
}

pub struct TaskBoundary {
    pub task_id: TaskId,
    pub constraint: String,
}

pub struct AttestedObservation {
    pub value: ObservableState,
    pub provenance: Provenance,
    pub task: TaskBoundary,
}

pub enum CollapseError {
    NotAdmissible(AdmissibilityError),
    ObservationMismatch { expected: ObservableState, supplied: ObservableState },
}

impl World {
    pub fn new(options: OptionSpace) -> Self {
        World { history: History::new(), options }
    }

    // Pop: commits to symbol, reduces option space
    pub fn pop(&mut self, sym: Symbol)
        -> Result<(), AdmissibilityError>
    {
        if !self.options.contains(&sym) {
            return Err(AdmissibilityError::NotInOptionSpace(sym));
        }
        if !self.admissible(&sym) {
            return Err(AdmissibilityError::PreviouslyRefused(sym));
        }
        self.history.append(Event::Pop(sym.clone()));
        self.options.remove(&sym);
        Ok(())
    }

    // Refuse: documents inadmissibility with reason
    pub fn refuse(&mut self, sym: Symbol, reason: RefusalReason) {
        self.history.append(Event::Refuse(sym, reason));
        // Option space unchanged: Refuse != Pop
    }

    // Bind: records dependency between symbols
    pub fn bind(&mut self, a: Symbol, b: Symbol) {
        self.history.append(Event::Bind(a, b));
    }

    // VIEW: pure projection. Appends nothing, mutates nothing.
    pub fn view(&self, rule: &CollapseRule) -> ObservableState {
        ObservableState::derive(&self.history, rule)
    }

    // COLLAPSE: attested adoption of a quotient. Verifies the supplied
    // observation against what the world would currently derive, checks
    // admissibility, and appends a Collapse event carrying full provenance.
    pub fn collapse(
        &mut self,
        subject: Symbol,
        rule: CollapseRule,
        observation: ObservableState,
        provenance: Provenance,
        task: TaskBoundary,
    ) -> Result<AttestedObservation, CollapseError> {
        let expected = self.view(&rule);
        if expected != observation {
            return Err(CollapseError::ObservationMismatch {
                expected,
                supplied: observation,
            });
        }
        if !self.admissible(&subject) {
            return Err(CollapseError::NotAdmissible(
                AdmissibilityError::PreviouslyRefused(subject),
            ));
        }
        self.history.append(Event::Collapse(
            subject, rule, observation.clone(), provenance.clone(), task.clone(),
        ));
        Ok(AttestedObservation { value: observation, provenance, task })
    }

    // Admissibility check
    fn admissible(&self, sym: &Symbol) -> bool {
        !self.history.events.iter().any(|e| matches!(
            e, Event::Refuse(s, _) if s == sym
        ))
    }

    // Conservation law invariant
    pub fn check_conservation(&self, initial_size: usize) -> bool {
        let history_weight: usize = self.history.events.iter()
            .filter(|e| matches!(e, Event::Pop(_)))
            .count();
        self.options.len() + history_weight == initial_size
    }
}
\end{lstlisting}

The event carries enough information to replay \emph{why} the quotient was
adopted, not merely which value resulted: \texttt{Provenance} identifies the
rule by a stable id and version (or a content hash of the rule specification),
since storing only the pair $(x, c)$ is insufficient once rule definitions are
allowed to change over the program's lifetime. \texttt{OperationalState}
(options, admissibility index, dependency graph) is represented explicitly
rather than left implicit beside history, and its materialization is checked
against the fold it is required to equal:

\begin{lstlisting}[style=rust, caption={Replay test for operational state}]
assert_eq!(world.materialized, fold_operational(&world.history));
\end{lstlisting}

Compiler correctness (Chapter~\ref{ch:compiler}) should therefore compare
\emph{both} the authoritative history and the derived operational state
between two candidate implementations, while history equality remains the
strictly stronger of the two criteria: two histories that are equal always
produce equal operational states under a deterministic fold, but two
implementations could in principle agree on operational state while disagreeing
on history (e.g., by different, non-replay-equivalent event sequences that fold
to the same cache), which history equality alone rules out and operational-state
equality alone would miss.

\section{The Evaluation Engine}

\begin{lstlisting}[style=rust, caption={Core evaluation loop}]
pub struct EvalResult {
    pub value: Value,
    pub world: World,
    pub admissible: bool,
    pub steps: usize,
}

pub fn eval(term: &Term, env: &Env, world: &mut World)
    -> Result<EvalResult, EvalError>
{
    let mut steps = 0;
    let result = eval_inner(term, env, world, &mut steps)?;
    Ok(EvalResult {
        value: result,
        world: world.clone(),
        admissible: world.history.is_admissible(),
        steps,
    })
}

fn eval_inner(
    term: &Term,
    env: &Env,
    world: &mut World,
    steps: &mut usize,
) -> Result<Value, EvalError> {
    *steps += 1;
    match term {
        Term::Var(x) => env.lookup(x).ok_or(EvalError::UnboundVar),
        Term::Lam(x, ty, b) => Ok(Value::Closure {
            param: x.clone(),
            param_ty: *ty.clone(),
            body: b.clone(),
            env: env.clone(),
        }),
        Term::App(f, a) => {
            let fv = eval_inner(f, env, world, steps)?;
            let av = eval_inner(a, env, world, steps)?;
            match fv {
                Value::Closure { param, body, env: cenv, .. } => {
                    let env2 = cenv.extend(&param, av);
                    eval_inner(&body, &env2, world, steps)
                }
                _ => Err(EvalError::NotAFunction),
            }
        }
        Term::Pop(t) => {
            let v = eval_inner(t, env, world, steps)?;
            let sym = v.to_symbol()?;
            world.pop(sym.clone())?;
            Ok(Value::Admitted(Box::new(Value::Symbol(sym))))
        }
        Term::Refuse(t, r) => {
            let v = eval_inner(t, env, world, steps)?;
            let sym = v.to_symbol()?;
            world.refuse(sym.clone(), r.clone());
            Ok(Value::Refused(Box::new(Value::Symbol(sym)), r.clone()))
        }
        Term::Bind(a, b) => {
            let av = eval_inner(a, env, world, steps)?;
            let bv = eval_inner(b, env, world, steps)?;
            let asym = av.to_symbol()?;
            let bsym = bv.to_symbol()?;
            world.bind(asym.clone(), bsym.clone());
            Ok(Value::Bound(Box::new(av), Box::new(bv)))
        }
        Term::Collapse(t, rule) => {
            let v = eval_inner(t, env, world, steps)?;
            match v {
                Value::Admitted(inner) => {
                    let obs = world.view(rule); // VIEW: pure lookup; a real Collapse commit would call world.collapse(...) with provenance and task
                    Ok(Value::Observed(inner, rule.clone(), obs))
                }
                _ => Err(EvalError::CollapseRequiresAdmissible),
            }
        }
        Term::Seq(terms) => {
            let mut last = Value::Unit;
            for t in terms {
                last = eval_inner(t, env, world, steps)?;
            }
            Ok(last)
        }
    }
}
\end{lstlisting}

\section{The Observable State Derivation}

\begin{lstlisting}[style=rust, caption={Collapse rules as observable state derivation}]
#[derive(Clone, Debug)]
pub enum CollapseRule {
    Identity,
    LastWrite,
    Accumulate,
    Projection(String),
}

#[derive(Clone, Debug, PartialEq)]
pub struct ObservableState {
    pub entries: HashMap<String, ObsValue>,
    pub rule: CollapseRule,
}

impl ObservableState {
    pub fn derive(history: &History, rule: &CollapseRule)
        -> ObservableState
    {
        let mut entries = HashMap::new();
        match rule {
            CollapseRule::Identity => {
                // Full history: every event is visible
                for (i, event) in history.events.iter().enumerate() {
                    entries.insert(
                        format!("event_{}", i),
                        ObsValue::Event(event.clone())
                    );
                }
            }
            CollapseRule::LastWrite => {
                // Most recent pop for each symbol
                for event in &history.events {
                    if let Event::Pop(sym) = event {
                        entries.insert(
                            sym.name.clone(),
                            ObsValue::Present
                        );
                    }
                    // Refusals after last pop remove symbol
                    if let Event::Refuse(sym, _) = event {
                        entries.remove(&sym.name);
                    }
                }
            }
            CollapseRule::Accumulate => {
                // Count of pops per symbol
                for event in &history.events {
                    if let Event::Pop(sym) = event {
                        let count = entries
                            .entry(sym.name.clone())
                            .or_insert(ObsValue::Count(0));
                        if let ObsValue::Count(n) = count {
                            *n += 1;
                        }
                    }
                }
            }
            CollapseRule::Projection(prefix) => {
                // Only events matching prefix
                for event in &history.events {
                    let name = event.symbol_name();
                    if name.starts_with(prefix.as_str()) {
                        entries.insert(
                            name.to_string(),
                            ObsValue::Event(event.clone())
                        );
                    }
                }
            }
        }
        ObservableState { entries, rule: rule.clone() }
    }

    pub fn equivalent(h1: &History, h2: &History,
                      rule: &CollapseRule) -> bool {
        Self::derive(h1, rule) == Self::derive(h2, rule)
    }
}
\end{lstlisting}

\chapter{The Compiler and Virtual Machine}

\section{The Event IR}

\begin{lstlisting}[style=rust, caption={Complete Event IR definition}]
#[derive(Clone, Debug, PartialEq)]
pub enum IrEvent {
    Pop,
    Refuse { reason: RefusalReason },
    Collapse { rule: CollapseRule },
    Bind,
    LoadSym(String),    // push symbol onto stack
    LoadVar(String),    // push variable binding
    Apply,              // function application
    LamIntro(String, Box<Type>),  // lambda introduction
    Seq(usize),         // sequence of n events
}

pub struct IrBlock {
    pub events: Vec<IrEvent>,
    pub source_map: Vec<usize>,  // event index -> source line
}
\end{lstlisting}

\section{The Compiler}

\begin{lstlisting}[style=rust, caption={The Spherepop compiler}]
pub struct Compiler {
    pub output: IrBlock,
}

impl Compiler {
    pub fn compile(&mut self, term: &Term) {
        match term {
            Term::Var(x) =>
                self.output.emit(IrEvent::LoadVar(x.clone())),
            Term::Lam(x, ty, body) => {
                self.output.emit(IrEvent::LamIntro(
                    x.clone(), ty.clone()
                ));
                self.compile(body);
            }
            Term::App(f, a) => {
                self.compile(f);
                self.compile(a);
                self.output.emit(IrEvent::Apply);
            }
            Term::Pop(t) => {
                self.compile(t);
                self.output.emit(IrEvent::Pop);
            }
            Term::Refuse(t, r) => {
                self.compile(t);
                self.output.emit(IrEvent::Refuse {
                    reason: r.clone()
                });
            }
            Term::Bind(a, b) => {
                self.compile(a);
                self.compile(b);
                self.output.emit(IrEvent::Bind);
            }
            Term::Collapse(t, rule) => {
                self.compile(t);
                self.output.emit(IrEvent::Collapse {
                    rule: rule.clone()
                });
            }
            Term::Seq(terms) => {
                let n = terms.len();
                for t in terms {
                    self.compile(t);
                }
                self.output.emit(IrEvent::Seq(n));
            }
        }
    }
}
\end{lstlisting}

\section{The Virtual Machine}

\begin{lstlisting}[style=rust, caption={The Spherepop VM}]
pub struct Machine {
    pub world: World,
    pub stack: Vec<Value>,
    pub call_stack: Vec<(String, Box<Term>, Env)>,
}

impl Machine {
    pub fn run(&mut self, block: &IrBlock)
        -> Result<Value, MachineError>
    {
        for event in &block.events {
            match event {
                IrEvent::LoadSym(s) =>
                    self.stack.push(Value::Symbol(Symbol::new(s))),
                IrEvent::LoadVar(x) => {
                    let v = self.current_env().lookup(x)
                        .ok_or(MachineError::UnboundVar)?;
                    self.stack.push(v);
                }
                IrEvent::Pop => {
                    let sym = self.stack.pop()
                        .ok_or(MachineError::StackUnderflow)?
                        .to_symbol()?;
                    self.world.pop(sym.clone())?;
                    self.stack.push(
                        Value::Admitted(Box::new(Value::Symbol(sym)))
                    );
                }
                IrEvent::Refuse { reason } => {
                    let sym = self.stack.pop()
                        .ok_or(MachineError::StackUnderflow)?
                        .to_symbol()?;
                    self.world.refuse(sym.clone(), reason.clone());
                    self.stack.push(Value::Refused(
                        Box::new(Value::Symbol(sym)), reason.clone()
                    ));
                }
                IrEvent::Bind => {
                    let b = self.stack.pop()
                        .ok_or(MachineError::StackUnderflow)?;
                    let a = self.stack.pop()
                        .ok_or(MachineError::StackUnderflow)?;
                    let asym = a.to_symbol()?;
                    let bsym = b.to_symbol()?;
                    self.world.bind(asym.clone(), bsym.clone());
                    self.stack.push(Value::Bound(
                        Box::new(a), Box::new(b)
                    ));
                }
                IrEvent::Collapse { rule } => {
                    let top = self.stack.pop()
                        .ok_or(MachineError::StackUnderflow)?;
                    match top {
                        Value::Admitted(inner) => {
                            let obs = self.world.view(rule); // VIEW: pure lookup, see world.collapse for the committed form
                            self.stack.push(Value::Observed(
                                inner, rule.clone(), obs
                            ));
                        }
                        _ => return Err(
                            MachineError::CollapseRequiresAdmissible
                        ),
                    }
                }
                IrEvent::Apply => {
                    let arg = self.stack.pop()
                        .ok_or(MachineError::StackUnderflow)?;
                    let func = self.stack.pop()
                        .ok_or(MachineError::StackUnderflow)?;
                    match func {
                        Value::Closure { param, body, env, .. } => {
                            let env2 = env.extend(&param, arg);
                            let result = eval_inner(
                                &body, &env2, &mut self.world,
                                &mut 0
                            )?;
                            self.stack.push(result);
                        }
                        _ => return Err(MachineError::NotAFunction),
                    }
                }
                _ => {} // LamIntro, Seq handled by context
            }
        }
        self.stack.pop().ok_or(MachineError::EmptyStack)
    }
}
\end{lstlisting}

\section{Replay Equivalence Test Suite}

\begin{lstlisting}[style=rust, caption={Replay equivalence verification}]
#[cfg(test)]
mod replay_tests {
    use super::*;

    fn assert_replay_equivalent(term: &Term) {
        let mut world_i = World::new(test_option_space());
        let result_i = eval(term, &Env::empty(), &mut world_i)
            .expect("interpretation failed");

        let mut compiler = Compiler::new();
        compiler.compile(term);
        let block = compiler.output;

        let mut machine = Machine::new(test_option_space());
        machine.run(&block).expect("VM execution failed");

        // THE CORRECTNESS CRITERION: history must be identical
        assert_eq!(
            world_i.history, machine.world.history,
            "Replay Equivalence violated: \
             I(p).history != V(C(p)).history"
        );
    }

    #[test]
    fn replay_pop() {
        assert_replay_equivalent(&Term::Pop(
            Box::new(Term::Var("x".into()))
        ));
    }

    #[test]
    fn replay_refuse_with_reason() {
        assert_replay_equivalent(&Term::Refuse(
            Box::new(Term::Var("x".into())),
            RefusalReason::ConstraintViolation("c".into())
        ));
    }

    #[test]
    fn replay_seq_bind_pop() {
        assert_replay_equivalent(&Term::Seq(vec![
            Term::Bind(
                Box::new(Term::Var("a".into())),
                Box::new(Term::Var("b".into()))
            ),
            Term::Pop(Box::new(Term::Var("a".into()))),
        ]));
    }

    #[test]
    fn replay_collapse_after_pop() {
        assert_replay_equivalent(&Term::Collapse(
            Box::new(Term::Pop(
                Box::new(Term::Var("x".into()))
            )),
            CollapseRule::LastWrite,
        ));
    }

    #[test]
    fn conservation_invariant_maintained() {
        let initial_size = 5;
        let mut world = World::new(OptionSpace::of_size(initial_size));
        world.pop(Symbol::new("a")).unwrap();
        world.pop(Symbol::new("b")).unwrap();
        world.refuse(Symbol::new("c"),
                     RefusalReason::NotAvailable);
        // Conservation: |H_pop| + |Omega| = initial_size
        assert!(world.check_conservation(initial_size));
    }
}
\end{lstlisting}


%% ============================================================
%% PART XXI: EXTENDED EXAMPLES AND CASE STUDIES
%% ============================================================
\part{Extended Examples and Case Studies}

\chapter{The Spherepop Type Checker as Its Own Subject}

\section{Self-Application}

The most striking demonstration of a type system's coherence is its ability to
type-check itself. The Spherepop type checker is written in a language that
Spherepop can reason about, and the type-checking process is itself a Spherepop
computation.

\begin{example}[Type Checking as Collapse]
The type checker applies a specific collapse rule $c_{\mathrm{type}}$ to the
history of term formation: each introduction rule ($[\mathrm{T\text{-}POP}]$,
$[\mathrm{T\text{-}REFUSE}]$, etc.) is a Pop event in the type-derivation
history, and the resulting type is the observable state under $c_{\mathrm{type}}$.

Type errors are Refuse events: the type checker refuses a term with a reason
(``collapse requires Admissible type, found Refused'') that becomes part of
the derivation history. The history of a type-checking run is itself a Spherepop
history, auditable by applying different collapse rules to extract different
aspects of the derivation.
\end{example}

\section{Proof Checking as Admissibility Verification}

\begin{example}[Mathematical Proof as Spherepop Computation]
A formal mathematical proof in a system like Lean 4 can be expressed as a
Spherepop computation. Each lemma application is a Pop event (committing to
a specific inference step). Each assumption is a Bind event (recording the
dependency on an axiom or hypothesis). Each application of modus ponens is a
Collapse event (collapsing the conditional and its antecedent to the consequent
under the modus-ponens rule).

The correctness of the proof corresponds to the Spherepop type:
$\vdash \mathit{proof} : \Adm(T)$ where $T$ is the statement being proved.
A proof error is $\Rfsd(T, r)$ where $r$ names the specific inference error
(wrong lemma applied, type mismatch in substitution, etc.).

The history of the proof-checking computation is the complete audit trail of
every inference step, every assumption used, and every definition invoked.
This history is more informative than the proof term alone: it records not
just \emph{what} was proved but \emph{how} the proof was discovered.
\end{example}

\chapter{Historical Case Studies}

\section{The Library of Alexandria as Note Collapse}

\begin{example}[The Library of Alexandria]
The Library of Alexandria (founded approximately 3rd century BCE, burned in
successive incidents between 48 BCE and 642 CE) was the ancient world's
largest note repository. At its peak, it held an estimated 400,000--700,000
scrolls, representing a significant fraction of the written knowledge of the
Mediterranean world.

In Spherepop terms, the Library was a civilizational note infrastructure
$\calN_{\mathrm{Alex}}$ whose destruction was a catastrophic note collapse.
By the Civilizational Collapse Theorem:
\[
R(Z_{\mathrm{after}}) < R(Z_{\mathrm{before}}).
\]

The lost trajectories include works of Aristotle's that are known only through
references in surviving texts (Capture Notes lost), mathematical texts by
Euclid's contemporaries whose results are inaccessible (Refusal Notes lost---we
cannot verify the admissibility of contemporary results against these benchmarks),
and technical treatises on engineering and medicine that would have enabled
trajectories in applied science that instead had to be rediscovered centuries
later (Generative Notes lost).

The repair morphisms that partially addressed the loss---the preservation of
Arabic translations of Greek texts, the transmission of some works through
Byzantine copies, the reconstruction of some mathematical results by independent
derivation---each restore specific fibers of the destroyed history, but none
restores the full reachability of the original infrastructure.
\end{example}

\section{Git as a Note Infrastructure}

\begin{example}[Git as Spherepop Implementation]
Git, the distributed version control system, is the most widely deployed
instance of Spherepop principles in existing software infrastructure.

The Git commit graph is a history $H_{\mathrm{git}}$: each commit is a Pop
event that records a specific state of the codebase. Branches are parallel
worlds $(H_1, H_2)$ that diverge from a common ancestor. Merges are history
merges $H_1 \sqcup H_2$. Tags are Collapse events that seal a specific state
under a named rule.

Git's admissibility structure is enforced by the branch protection rules:
a Refuse event is a rejected pull request, with a reason (failing CI, code
review rejection, merge conflict). The reason is preserved in the Git history
as a closed pull request with comments.

The collapse rules in Git correspond to:
\begin{itemize}
\item $c_I$ (Identity): \texttt{git log --all} showing the full history.
\item $c_{LW}$ (LastWrite): the current working tree state.
\item $c_{\mathrm{acc}}$ (Accumulate): \texttt{git shortlog} counting commits per author.
\item $\pi_L$ (Projection): \texttt{git log -- path/to/file} showing only events affecting a specific path.
\end{itemize}

The Replay Equivalence theorem (Theorem~10.3) is implemented in Git's architecture:
checking out any commit and applying all subsequent commits must produce the
same working tree state as checking out the latest commit directly. Git's
content-addressable storage ensures this invariant by making commit hashes
depend on the full history.
\end{example}

\chapter{Spherepop Programs as Executable Specifications}

\section{The Dining Philosophers in Spherepop}

\begin{example}[Dining Philosophers]
The classic Dining Philosophers problem (five philosophers, five forks, each
philosopher needs two adjacent forks to eat) is a deadlock scenario expressible
in the Spherepop process calculus.
\begin{lstlisting}[style=rust, caption={Dining philosophers as Spherepop}]
// Five philosophers and five forks
let fork_options = OptionSpace::new([
    Symbol::new("fork0"), Symbol::new("fork1"),
    Symbol::new("fork2"), Symbol::new("fork3"),
    Symbol::new("fork4"),
]);

// Philosopher i picks up left fork (i) then right fork (i+1 mod 5)
fn philosopher(i: usize, world: &mut World)
    -> Result<(), AdmissibilityError>
{
    let left = Symbol::new(&format!("fork{}", i));
    let right = Symbol::new(&format!("fork{}", (i + 1) % 5));
    // Try to acquire left fork
    world.pop(left.clone())?;
    // Try to acquire right fork
    match world.pop(right.clone()) {
        Ok(_) => {
            // Eat, then release
            world.refuse(left, RefusalReason::Explicit(
                "released".into()
            ));
            world.refuse(right, RefusalReason::Explicit(
                "released".into()
            ));
            Ok(())
        }
        Err(e) => {
            // Deadlock: right fork unavailable
            // Must release left and retry
            Err(e)
        }
    }
}
\end{lstlisting}
Deadlock is detected by the cycle detection algorithm (Theorem~34.2): when all
five philosophers hold their left fork and wait for their right fork, the
dependency graph $G_H$ has the cycle $0 \to 1 \to 2 \to 3 \to 4 \to 0$.
The Spherepop history preserves which philosopher acquired which fork in which
order, making the deadlock state fully auditable.
\end{example}

\section{A Blockchain as Spherepop History}

\begin{example}[Blockchain as Distributed Spherepop]
A blockchain is a distributed Spherepop world system where each block is a
Collapse event that seals a batch of transactions (Pop events) under the
Accumulate collapse rule, producing a cryptographic hash (the block hash)
as the observable state.

The blockchain's consensus mechanism corresponds to global consistency (Definition~33.2):
all nodes in the network must agree on the same history $H$ and hence the same
block hashes. A blockchain fork is a network partition that allows two sets of
nodes to produce divergent histories $H_1 \neq H_2$. The longest-chain rule
(in Bitcoin) is a specific merge strategy: when the partition heals, adopt
the history with more accumulated proof-of-work, corresponding to the history
with higher total action $\mathcal{S}(H)$ in the variational sense.

Smart contracts are Spherepop programs executed within the blockchain's history.
They commit to outcomes (Pop), document failures (Refuse), record dependencies
between contracts (Bind), and produce verifiable states (Collapse). The
immutability of the blockchain history corresponds to the monotonicity of
Spherepop history growth: events can be added but never removed.
\end{example}


%% ============================================================
%% PART XXII: POSSIBILITY DYNAMICS AND ENTROPY
%% ============================================================
\part{Possibility Dynamics and Entropy}

\chapter{Possibility as a Computational Resource}

\section{Possibility Entropy}

The conservation law counts possibility discretely. A more informative invariant
is obtained by passing from cardinality to entropy.

\begin{definition}[Possibility Entropy]
The \emph{possibility entropy} of a history $H$ is
\[
S_\Omega(H) = \log |\Omega(H)|.
\]
When $\Omega(H) = \varnothing$, we set $S_\Omega(H) = -\infty$.
\end{definition}

\begin{theorem}[Pop Entropy Reduction]
If $H' = H \mathbin{++} [\Pop(x)]$ and $x \in \Omega(H)$ with $|\Omega(H)| > 1$, then
\[
S_\Omega(H') < S_\Omega(H).
\]
\end{theorem}

\begin{proof}
Since Pop removes exactly one available symbol,
$|\Omega(H')| = |\Omega(H)| - 1$.
If $|\Omega(H)| > 1$, then $0 < |\Omega(H)| - 1 < |\Omega(H)|$.
The logarithm is strictly increasing, so
$\log(|\Omega(H)| - 1) < \log|\Omega(H)|$.
Therefore $S_\Omega(H') < S_\Omega(H)$.
\end{proof}

\begin{corollary}[Refuse Does Not Reduce Structural Entropy]
If $H' = H \mathbin{++} [\Refuse(x,r)]$, then $S_\Omega(H') = S_\Omega(H)$.
\end{corollary}

\begin{proof}
Refuse appends an event but does not remove $x$ from $\Omega$.
Thus $\Omega(H') = \Omega(H)$, so $S_\Omega(H') = \log|\Omega(H')| = \log|\Omega(H)| = S_\Omega(H)$.
\end{proof}

\begin{remark}
This separates two kinds of possibility loss. Pop reduces the option space:
it commits to a specific future by removing alternatives. Refuse does not reduce
structural availability; it records that a structurally available path is
inadmissible. The two entropies---structural ($S_\Omega$) and admissibility
($S_{\calA}$)---track these independent dimensions.
\end{remark}

\begin{example}[Choosing from a Menu]
A diner faces a menu with $n$ options. Each choice is a Pop event. Before
choosing: $S_\Omega = \log n$. After choosing the main course: $S_\Omega = \log(n-1)$.
The dinner's preference rules (vegetarian, allergies) are Refuse events: they
reduce $S_{\calA}$ without affecting $S_\Omega$. The restaurant still offers the
refused dishes; the diner simply will not order them. This models real decisions
precisely: structural and admissibility constraints are independent.
\end{example}

\begin{example}[Compiler Optimization as Possibility Reduction]
A compiler transforms high-level source (high $S_\Omega$: many possible
execution paths) into optimized machine code (low $S_\Omega$: fewer execution
branches). Each optimization pass is a sequence of Pop events---selecting specific
implementations for abstract operations---reducing possibility entropy while
preserving the observable behavior (the collapse rule $c$ applied to source
and output produces the same observable state).
\end{example}

\begin{example}[Biological Differentiation]
A pluripotent stem cell has high possibility entropy: it can differentiate into
many cell types. Differentiation is a sequence of Pop events, each committing
to specific gene expression patterns. A fully differentiated neuron has
$S_\Omega \approx 0$: it has committed to its cell type. Apoptosis signals
correspond to Refuse events: they mark developmental paths as inadmissible
without immediately committing the cell to an alternative.
\end{example}

\begin{example}[Legal Precedent as Pop]
A legal precedent is a Pop event in a jurisdiction's possibility space. Before
the ruling, multiple legal interpretations are structurally available
($S_\Omega$ high). After the precedent, the court has committed to one
interpretation, reducing $S_\Omega$. Dissenting opinions are Refuse events:
they document that alternative interpretations were considered and found
inadmissible (in the view of dissenting justices), without reducing the
structural option space (which is still available for future courts to revisit).
\end{example}

\section{Admissibility Entropy}

Structural possibility counts what remains selectable. Admissibility counts
what remains justifiable.

\begin{definition}[Admissibility Entropy]
The \emph{admissibility entropy} of $H$ is
\[
S_{\calA}(H) = \log |\calA(H)|.
\]
\end{definition}

\begin{theorem}[Refuse Entropy Reduction]
If $H' = H \mathbin{++} [\Refuse(x,r)]$ and $x \in \calA(H)$ with $|\calA(H)| > 1$, then
\[
S_{\calA}(H') < S_{\calA}(H).
\]
\end{theorem}

\begin{proof}
The refusal of $x$ removes $x$ from the admissibility set:
$\calA(H') = \calA(H) \setminus \{x\}$, so $|\calA(H')| = |\calA(H)| - 1$.
Since the logarithm is strictly increasing, $\log|\calA(H')| < \log|\calA(H)|$.
\end{proof}

\begin{corollary}[Pop and Refuse Act on Different Entropies]
Pop strictly reduces $S_\Omega$ but need not reduce $S_{\calA}$.
Refuse strictly reduces $S_{\calA}$ but does not reduce $S_\Omega$.
\end{corollary}

\begin{proof}
Pop removes a symbol from $\Omega$ by definition. Refuse removes a symbol from
$\calA(H)$ by definition but leaves $\Omega$ unchanged.
\end{proof}

\begin{remark}[The Entropy Dual]
The two entropies form a dual pair:
\[
S_\Omega(H) \geq S_{\calA}(H) \quad \text{for all } H,
\]
because $\calA(H) \subseteq \Omega(H)$ always holds (admissible futures are
a subset of structurally available futures). The gap $S_\Omega(H) - S_{\calA}(H) = \log|\Omega(H)/\calA(H)|$
measures the admissibility burden: the proportion of structurally available
futures that are inadmissible given the current refusal history.
\end{remark}

\section{Possibility Debt}

A computation that accumulates unresolved possibilities incurs a burden: every
undecided symbol must eventually be committed or refused, and the cost of that
decision is deferred but not eliminated.

\begin{definition}[Decision Cost Function]
Let $C_s(x, H)$ denote the future cost of resolving whether symbol
$x \in \Omega(H)$ should be committed, refused, bound, or collapsed.
\end{definition}

\begin{definition}[Possibility Debt]
The \emph{possibility debt} of history $H$ is
\[
D_\Omega(H) = \sum_{x \in \Omega(H)} C_s(x, H).
\]
\end{definition}

\begin{proposition}[Pop Reduces Local Possibility Debt]
If $H' = H \mathbin{++} [\Pop(x)]$, then
\[
D_\Omega(H') = D_\Omega(H) - C_s(x, H) + \Delta,
\]
where $\Delta = \sum_{y \in \Omega(H')} (C_s(y, H') - C_s(y, H))$ is the
change in decision cost for the remaining options.
\end{proposition}

\begin{proof}
Before Pop, the debt sum includes $C_s(x, H)$ and all remaining terms.
After Pop, $x$ is removed from $\Omega$. Let $\Delta$ be as defined.
Then $D_\Omega(H') = D_\Omega(H) - C_s(x, H) + \Delta$.
\end{proof}

\begin{remark}
A Pop operation is \emph{locally simplifying} when $\Delta < C_s(x, H)$: the
remaining options become cheaper after the commitment. It is \emph{globally
complicating} when $\Delta > C_s(x, H)$: committing to $x$ increases the
downstream complexity of the remaining decisions. This models the phenomenon
of premature commitment: early choices can constrain later choices in costly
ways, even when they simplify the immediate decision.
\end{remark}

\begin{example}[Technical Debt in Software]
Technical debt is possibility debt in software development. A codebase with
many unresolved design decisions (high $D_\Omega$) imposes costs on every
future change: the developer must navigate around or resolve these decisions.
Refactoring is a sequence of Pop and Refuse events that resolves deferred
decisions, reducing $D_\Omega$ at the cost of immediate effort. The observation
that technical debt compounds corresponds to $\Delta > C_s(x, H)$ in many
software contexts.
\end{example}

\section{Reversible and Irreversible Computation}

\begin{definition}[Reversible Computation]
A computation is \emph{reversible} if every operation has an inverse that
restores the prior history. In Spherepop, an operation $e$ is reversible if
there exists an event $e^{-1}$ such that
$H \mathbin{++} [e] \mathbin{++} [e^{-1}] \sim_c H$ for all collapse rules $c$.
\end{definition}

\begin{theorem}[Pop Is Irreversible]
No event $e^{-1}$ can reverse a Pop operation $\Pop(x)$ while preserving
history monotonicity.
\end{theorem}

\begin{proof}
By the Conservation Law, $|\Omega(H \mathbin{++} [\Pop(x)])| = |\Omega(H)| - 1$.
Any subsequent event can only decrease or maintain $|\Omega|$; no event increases
it (since $\Omega$ only shrinks). Therefore the option space after $\Pop(x)$
can never recover to $|\Omega(H)|$, and the Pop is irreversible.
\end{proof}

\begin{remark}[Connection to Thermodynamics]
The irreversibility of Pop corresponds to the thermodynamic irreversibility of
measurement and commitment. In reversible computing (Bennett's erasure principle),
reversible operations have zero thermodynamic cost; irreversible operations
(like Pop) dissipate energy equal to $k_B T \ln 2$ per bit erased. The
Spherepop Conservation Law is the computational analogue of the thermodynamic
entropy production law: commitment converts possibility into history at a fixed
rate, with no way to convert history back into possibility without external
work.
\end{remark}

\chapter{Collapse Curvature and Fiber Geometry}

\section{Collapse Curvature}

\begin{definition}[Collapse Fiber]
For collapse rule $c : \calH \to O_c$ and observable state $o \in O_c$,
the \emph{collapse fiber} over $o$ is $F_o = c^{-1}(o)$.
\end{definition}

\begin{definition}[Collapse Curvature]
The \emph{collapse curvature} of history $H$ under rule $c$ is
\[
\kappa_c(H) = \log |F_{c(H)}| = \log |c^{-1}(c(H))|.
\]
\end{definition}

\begin{theorem}[Identity Collapse Has Zero Curvature]
For the identity rule, $\kappa_{c_I}(H) = 0$ for all $H$.
\end{theorem}

\begin{proof}
Under $c_I$, the fiber over $H$ is $\{H\}$, so $|F_H| = 1$ and
$\kappa_{c_I}(H) = \log 1 = 0$.
\end{proof}

\begin{theorem}[Coarser Collapse Increases Curvature]
If $c_2$ is coarser than $c_1$ (i.e.\ $H \sim_{c_1} H' \Rightarrow H \sim_{c_2} H'$), then
$\kappa_{c_2}(H) \geq \kappa_{c_1}(H)$ for all $H$.
\end{theorem}

\begin{proof}
Coarsening means every history identified with $H$ by $c_1$ is also identified
by $c_2$, so $F^{c_1}_{c_1(H)} \subseteq F^{c_2}_{c_2(H)}$. Therefore
$|F^{c_1}_{c_1(H)}| \leq |F^{c_2}_{c_2(H)}|$, and since log is monotone,
$\kappa_{c_1}(H) \leq \kappa_{c_2}(H)$.
\end{proof}

\begin{corollary}[Maximum Collapse Theorem]
State-only representations maximize collapse curvature among all collapse rules
with finite observable space. That is, the LastWrite rule $c_{LW}$ achieves
the maximum curvature over the canonical rules.
\end{corollary}

\begin{proof}
$c_{LW}$ identifies histories that differ only in intermediate states, retaining
only the final committed value per symbol. This is coarser than the Identity rule
(which retains all history) and coarser than the Accumulate rule (which retains
counts). By Coarser Collapse Increases Curvature, $c_{LW}$ has higher curvature
than these alternatives.
\end{proof}

\begin{remark}
The collapse curvature $\kappa_c(H) = \log|c^{-1}(c(H))|$ is precisely the
observational entropy $S_c(o)$ from Part~IX, now interpreted geometrically as
the curvature of the quotient space. A high-curvature collapse rule flattens
history space into a few large equivalence classes; a low-curvature rule
preserves most historical distinctions. The connection to CLIO is direct:
CLIO's representational entropy is the collapse curvature of the projection
functor.
\end{remark}

\section{Fiber Structure of Histories}

\begin{definition}[History Fiber]
Given a state projection $\sigma : \calH \to S$, the \emph{fiber} over state
$s \in S$ is $F_s = \sigma^{-1}(s) = \{H \in \calH \mid \sigma(H) = s\}$.
\end{definition}

\begin{theorem}[Average State Degeneracy Growth]
Let $S$ be a finite state space with $|S|$ elements and let
$\sigma_n : \calH_n \to S$ be any state projection from length-$n$ histories.
Then the average fiber size is
\[
\bar{D}_n = \frac{|E|^n}{|S|},
\]
which grows exponentially in $n$ when $|E| > 1$.
\end{theorem}

\begin{proof}
The number of length-$n$ histories over alphabet $E$ is $|\calH_n| = |E|^n$.
The projection $\sigma_n$ partitions $\calH_n$ into $|S|$ fibers.
The average fiber size is $\bar{D}_n = |E|^n / |S|$.
Since $|S|$ is constant and $|E| > 1$, this grows exponentially in $n$.
\end{proof}

\begin{corollary}[Long Computations Are More Collapsed]
For fixed finite observable state space $S$, longer histories are increasingly
underdetermined by their final state. The fraction of historical information
retained by the final state decreases as $|S| / |E|^n \to 0$.
\end{corollary}

\begin{remark}
This is the formal content of the state illusion: as a computation proceeds,
each additional step multiplies the number of distinct histories consistent
with the current observable state. A debugger examining a running program
faces exponentially many possible execution paths that produced the current
state.
\end{remark}

\section{Fiber Dimension and Geometry}

\begin{definition}[Fiber Dimension]
The \emph{fiber dimension} of observation $o$ under rule $c$ is
$\dim(F_o) = \log_2 |c^{-1}(o)|$---the number of binary digits needed to
distinguish histories within the fiber.
\end{definition}

\begin{definition}[Projection Distortion]
The \emph{projection distortion} of collapse rule $c$ at history $H$ is
\[
\delta_c(H) = \kappa_c(H) = \log|c^{-1}(c(H))|.
\]
\end{definition}

\begin{theorem}[Projection Distortion Theorem]
Larger fibers imply larger projection distortion, and distortion is bounded
below by zero (Identity rule) and above by $\log|\calH_n|$ (trivial rule
collapsing all histories to one point).
\end{theorem}

\begin{proof}
$\delta_c(H) = \kappa_c(H) \geq 0$ with equality iff the fiber is a singleton
(injective collapse). The upper bound $\log|\calH_n|$ is achieved by the trivial
rule $c_\top$ that maps all histories to a single observable, producing one
fiber of size $|E|^n$.
\end{proof}

\chapter{The Prefix Topology and History Metrics}

\section{Prefix Topology on Histories}

Histories can be topologized by their common prefix structure, giving a natural
notion of ``closeness'' for computational trajectories.

\begin{definition}[Prefix Cylinder]
For a finite history $P$, define the \emph{prefix cylinder}
\[
U_P = \{H \in \calH \mid P \text{ is a prefix of } H\}.
\]
\end{definition}

\begin{theorem}[Prefix Cylinders Form a Basis]
The collection $\mathcal{B} = \{U_P \mid P \in \calH\}$ forms a basis for a
topology on $\calH$.
\end{theorem}

\begin{proof}
Every history $H$ belongs to $U_H$, so $\mathcal{B}$ covers $\calH$.

For the intersection property: suppose $H \in U_P \cap U_Q$. Then both $P$ and
$Q$ are prefixes of $H$. For finite sequences, if two sequences are both prefixes
of $H$, then one is a prefix of the other (they form a total order under the
prefix relation). Assume $P$ is a prefix of $Q$. Then $U_Q \subseteq U_P$, so
$H \in U_Q \subseteq U_P \cap U_Q$. For every point in the intersection, the
basis element $U_Q$ contains that point and is contained in the intersection.
Therefore $\mathcal{B}$ is a basis.
\end{proof}

\begin{remark}
In the prefix topology, two histories are topologically close when they share
a long initial segment. Computations with common early development occupy the
same local region of history space, even if they later diverge. This captures
the intuition that two program executions that run the same initialization code
are ``nearby'' even if their subsequent paths differ.
\end{remark}

\begin{proposition}[Continuation Neighborhoods]
For history $H$, the set $U_H$ is the basic open neighborhood of $H$ in the
prefix topology. Every continuation of $H$ lies in $U_H$. Diverging computations
that share prefix $P$ but not $H$ lie in $U_P \setminus U_H$.
\end{proposition}

\begin{proof}
A continuation of $H$ is any $H' = H \mathbin{++} H''$ for some non-empty $H''$.
Such $H'$ has $H$ as a prefix, so $H' \in U_H$. A computation in $U_P \setminus U_H$
shares prefix $P$ but diverges after $P$, before reaching the state corresponding
to $H$.
\end{proof}

\section{The Prefix Ultrametric}

\begin{definition}[Longest Common Prefix Length]
For histories $H_1, H_2$, let $\mathsf{lcp}(H_1, H_2)$ denote the length of
their longest common prefix.
\end{definition}

\begin{definition}[Prefix Distance]
Define
\[
d_p(H_1, H_2) = 2^{-\mathsf{lcp}(H_1, H_2)}.
\]
For identical finite histories, $d_p(H, H) = 0$.
\end{definition}

\begin{theorem}[Prefix Distance Is an Ultrametric]
The function $d_p$ satisfies the ultrametric inequality:
\[
d_p(H_1, H_3) \leq \max\{d_p(H_1, H_2),\, d_p(H_2, H_3)\}.
\]
\end{theorem}

\begin{proof}
Let $a = \mathsf{lcp}(H_1, H_2)$ and $b = \mathsf{lcp}(H_2, H_3)$.
Then $H_1$ and $H_3$ must share at least the first $\min(a, b)$ events,
since both agree with $H_2$ on that prefix. Hence
$\mathsf{lcp}(H_1, H_3) \geq \min(a, b)$. Therefore
$d_p(H_1, H_3) = 2^{-\mathsf{lcp}(H_1,H_3)} \leq 2^{-\min(a,b)} = \max\{2^{-a}, 2^{-b}\} = \max\{d_p(H_1,H_2), d_p(H_2,H_3)\}$.
\end{proof}

\begin{remark}
The ultrametric inequality is stronger than the ordinary triangle inequality:
$d(x,z) \leq \max\{d(x,y), d(y,z)\}$ implies $d(x,z) \leq d(x,y) + d(y,z)$
but not vice versa. Ultrametrics arise naturally in $p$-adic number theory
and in the study of hierarchical structures, both of which are relevant to
computation: program execution traces form a tree whose natural metric is
ultrametric.
\end{remark}

\section{Geodesics in History Space}

\begin{definition}[History Geodesic]
A \emph{geodesic} between histories $H_0$ and $H_n$ in the prefix metric
is a sequence $H_0, H_1, \ldots, H_n$ where each $H_{i+1}$ is obtained
from $H_i$ by appending exactly one event, minimizing the total path length
under the action functional $\mathcal{S}$.
\end{definition}

\begin{theorem}[Geodesic Repair Theorem]
The shortest repair path from a failed history $H$ to a target observable state
$o$ is the history geodesic from $H$ to the nearest element of $c^{-1}(o)$
in the prefix metric.
\end{theorem}

\begin{proof}
A repair morphism $\rho : H \to H'$ with $c(H') = o$ requires constructing
a new history $H'$ that maps to $o$ under $c$. The repair cost is
$d_{\mathrm{edit}}(H, H')$---the minimum number of event substitutions needed.
In the prefix metric, the nearest history in $c^{-1}(o)$ is the one with
the longest common prefix with $H$, minimizing the number of events that must
differ. This is precisely the minimum-edit-distance repair.
\end{proof}

\chapter{The Full Semantics of Bind}

\section{Motivation: Bind Has Syntax but No Ontology}

The critique identified what is perhaps the most important gap in the theory.
$\Bind(a, b)$ appears throughout the monograph as a dependency declaration,
but its operational semantics---what it actually \emph{does} to computation---has
not been fully specified.

Three interpretations are possible, and distinguishing them clarifies the design:

\begin{enumerate}
\item \emph{Bind as annotation}: $\Bind(a, b)$ is a pure record in the history,
with no immediate operational effect. Future agents reading the history can see
the dependency, but the runtime does not enforce it.
\item \emph{Bind as constraint}: $\Bind(a, b)$ establishes that if $a$ is refused,
$b$ becomes inadmissible. The runtime enforces dependency propagation.
\item \emph{Bind as synchronization}: $\Bind(a, b)$ requires that $a$ and $b$
be committed simultaneously in parallel contexts---a synchronization barrier.
\end{enumerate}

All three are coherent; they correspond to different use cases. This chapter
develops interpretation (2) as the primary operational semantics, with the
others as derived modes.

\section{Dependency Graphs}

\begin{definition}[Dependency Graph]
The \emph{dependency graph} of history $H$ is the directed graph
$G_H = (V, E)$ where $V = \Omega_0$ and
$(a, b) \in E$ iff $\Bind(a, b) \in H$.
\end{definition}

\begin{definition}[Dependency Closure]
The \emph{dependency closure} of a set $X \subseteq \Omega_0$ under $G_H$ is
\[
\overline{X}^{G_H} = \{y \in \Omega_0 \mid \text{there exists a directed path from some } x \in X \text{ to } y \text{ in } G_H\}.
\]
\end{definition}

\section{Refusal Propagation}

\begin{definition}[Propagation Rule]
Given $\Bind(a, b) \in H$ and a refusal $\Refuse(a, r) \in H$, the
\emph{propagated refusal} for $b$ is $\Refuse(b, \rho(r))$ where
$\rho : \mathsf{RefusalReason} \to \mathsf{RefusalReason}$ is a reason
transformation (e.g.\ $\rho(r) = \mathsf{DependencyFailed}(a, r)$).
\end{definition}

\begin{definition}[Refusal Closure]
The \emph{refusal closure} of history $H$ under dependency propagation is the
smallest extension $H^* \supseteq H$ such that for every $\Bind(a, b) \in H$
and $\Refuse(a, r) \in H^*$, we have $\Refuse(b, \rho(r)) \in H^*$.
\end{definition}

\begin{theorem}[Dependency Propagation Theorem]
\label{thm:dep-propagation}
Refusals propagate along dependency edges: if $\Bind(a, b) \in H$ and
$\Refuse(a, r) \in H$, then $b \notin \calA(H^*)$.
\end{theorem}

\begin{proof}
By definition of refusal closure, $\Refuse(b, \rho(r)) \in H^*$.
By definition of $\calA$, $b \notin \calA(H^*)$.
\end{proof}

\begin{theorem}[Dependency Closure Theorem]
In an acyclic dependency graph $G_H$, every refusal generates a unique
maximal refusal closure.
\end{theorem}

\begin{proof}
Acyclicity ensures that dependency propagation terminates: since every path in
$G_H$ has finite length (no cycles), the propagation reaches every dependent
node exactly once. Uniqueness follows from the determinism of $\rho$: each
refusal generates exactly one propagated reason at each dependent node.
Maximality holds because we take the closure under all dependency edges.
\end{proof}

\begin{corollary}[Deadlock Detection via Dependency Closure]
A dependency cycle $a_1 \to a_2 \to \cdots \to a_k \to a_1$ in $G_H$
with any $a_i$ refused implies that the refusal propagates around the cycle
indefinitely, making all $a_i$ inadmissible.
\end{corollary}

\begin{proof}
Starting from $\Refuse(a_i, r)$, propagation reaches $a_{i+1}, \ldots, a_k,
a_1, a_2, \ldots$ continuing around the cycle. All elements of the cycle
become inadmissible under the refusal closure.
\end{proof}

\section{Bind Symmetry and Asymmetry}

\begin{proposition}[Bind Is Not Symmetric]
In general, $\Bind(a, b)$ and $\Bind(b, a)$ have different operational effects:
the former propagates refusals from $a$ to $b$, the latter from $b$ to $a$.
\end{proposition}

\begin{proof}
Under the propagation rule, $\Bind(a, b)$ causes $\Refuse(a, r) \Rightarrow
\Refuse(b, \rho(r))$, while $\Bind(b, a)$ causes $\Refuse(b, r) \Rightarrow
\Refuse(a, \rho(r))$. These are distinct effects unless $\calA$ treats $a$ and
$b$ symmetrically.
\end{proof}

\begin{definition}[Symmetric Bind]
$\Bind(a, b)$ is \emph{symmetric} if both $\Bind(a, b)$ and $\Bind(b, a)$ are
in the history, creating bidirectional dependency.
\end{definition}

\begin{example}
A symmetric bind models mutual dependency: two modules that each require the
other to be non-refused. An asymmetric bind models unidirectional dependency:
a database connection that fails if the database server is refused, but the
server is not affected by failures in the client connection.
\end{example}

\section{Bind as Coordination: the Multi-Agent Case}

\begin{theorem}[Bind Creates Coupled Trajectory Spaces]
Let agents $A_1$ and $A_2$ have independent continuation spaces $\calC_1$ and
$\calC_2$. A bind $\Bind(a, b)$ with $a \in \Omega_1$ and $b \in \Omega_2$
creates a coupled trajectory space $\calC_{12}$ where
\[
|\calC_{12}| > |\calC_1| + |\calC_2|
\]
whenever new joint continuations are enabled by the coupling.
\end{theorem}

\begin{proof}
Without the bind, agents evolve independently. A joint continuation requiring
both $a$ to be committed and $b$ to be committed at related times is only
reachable when the agents are coupled. The coupling via Bind enables this
joint continuation, adding to the total reachability beyond the independent sum.
\end{proof}

\chapter{Reconstruction Geometry and Cost Models}

\section{Defining Reconstruction Cost Formally}

The reconstruction cost $C_r(c)$ has been used throughout the monograph as
a primitive. This chapter gives it a formal definition.

\begin{definition}[Reconstruction Problem]
Given an observable state $o \in O_c$ and collapse rule $c$, the
\emph{reconstruction problem} is: find the history $H \in c^{-1}(o)$ that
produced $o$.
\end{definition}

\begin{definition}[Reconstruction Cost]
The reconstruction cost $C_r(c, o)$ is the minimum number of oracle queries
to $c^{-1}(o)$ (each query reveals one event of the unknown history) needed
to identify $H$ with probability at least $1 - \varepsilon$, for fixed
error tolerance $\varepsilon > 0$.
\end{definition}

\begin{theorem}[Information-Theoretic Lower Bound]
\[
C_r(c, o) \geq \frac{S_c(o)}{\log|\Sigma|}
\]
where $\Sigma$ is the event alphabet and $S_c(o) = \log|c^{-1}(o)|$ is the
observational entropy.
\end{theorem}

\begin{proof}
The fiber $c^{-1}(o)$ contains $|c^{-1}(o)| = 2^{S_c(o)}$ candidate histories.
Each oracle query reveals one event (chosen from $\Sigma$), providing
$\log|\Sigma|$ bits of information. To distinguish $2^{S_c(o)}$ candidates
requires at least $S_c(o) / \log|\Sigma|$ queries.
\end{proof}

\begin{theorem}[Note Leverage as Entropy Reduction Rate]
The reachability leverage of a note $N$ satisfies
\[
\Lambda(N) = \frac{k}{C_{\mathrm{create}}(N)} \cdot (S_c(o) - S_c(o \mid N)),
\]
where $k = 1/\log|\Sigma|$ is the information density constant and
$S_c(o \mid N)$ is the entropy of observation $o$ when $N$ is available.
\end{theorem}

\begin{proof}
Note $N$ reduces the effective fiber size from $|c^{-1}(o)|$ to $|c^{-1}(o) \cap
\mathsf{Compatible}(N)|$. The reconstruction cost reduction is
$\Delta C_r = C_r(c,o) - C_r(c,o \mid N) \geq k(S_c(o) - S_c(o \mid N))$.
Dividing by $C_{\mathrm{create}}(N)$ gives $\Lambda(N) \geq k(S_c(o) - S_c(o \mid N)) / C_{\mathrm{create}}(N)$.
\end{proof}

\begin{corollary}[Optimal Note Theorem]
The optimal note $N^*$ for minimizing reconstruction cost of observations under
$c$ is the note that maximizes $(S_c(o) - S_c(o \mid N)) / C_{\mathrm{create}}(N)$:
maximum entropy reduction per unit creation cost.
\end{corollary}

\section{Reconstruction Distance}

\begin{definition}[Reconstruction Distance]
The \emph{reconstruction distance} between two histories $H_1, H_2 \in c^{-1}(o)$
is the minimum number of events that must be queried to distinguish $H_1$ from
$H_2$, given that both are in the same fiber.
\end{definition}

\begin{proposition}[Reconstruction Distance Equals Edit Distance in Fiber]
The reconstruction distance between $H_1$ and $H_2$ in the same fiber $c^{-1}(o)$
equals the Hamming distance $d_H(H_1, H_2)$ when the fiber consists of
equal-length histories.
\end{proposition}

\begin{proof}
Two histories in the same fiber are observationally equivalent under $c$:
they produce the same observable state. To distinguish them, an oracle must
identify the positions where they differ. The minimum number of queries needed
is the Hamming distance---the number of positions where they actually differ.
\end{proof}

\section{The Reachability Leverage Spectrum}

Different note classes achieve different positions on the leverage spectrum.

\begin{center}
\begin{tabular}{llll}
\toprule
Note Class & Entropy Reduction & Creation Cost & Typical $\Lambda$ \\
\midrule
Grocery list & Low & Very low & $\approx 100$ \\
Commit message & Medium & Low & $\approx 10{,}000$ \\
One-line equation & Very high & Low & $\approx 10^6$ \\
Mathematical proof & Extremely high & Medium & $\approx 10^9$ \\
Programming language & Vast & High & $\approx 10^{12}$ \\
Scientific theory & Unbounded & Very high & $\to \infty$ \\
\bottomrule
\end{tabular}
\end{center}

The divergence of $\Lambda$ for scientific theories reflects the fact that a
theory like Newtonian mechanics reduces the reconstruction cost of infinitely
many observations (planetary positions, projectile trajectories, tidal cycles)
to near zero for any agent fluent in the theory.

\chapter{Refusal Logic and the Algebra of Inadmissibility}

\section{Refusal Sets as a Monotone Structure}

\begin{definition}[Refusal Set]
For a history $H$, define the refusal set
$\calR(H) = \{(x, r) \mid \Refuse(x, r) \in H\}$.
\end{definition}

\begin{definition}[Refusal Order]
Define $H_1 \preceq_R H_2$ iff $\calR(H_1) \subseteq \calR(H_2)$.
\end{definition}

\begin{theorem}[Refusal Monotonicity]
If $H' = H \mathbin{++} [e]$ for any event $e$, then $H \preceq_R H'$.
\end{theorem}

\begin{proof}
If $e = \Refuse(x, r)$, then $\calR(H') = \calR(H) \cup \{(x,r)\} \supseteq \calR(H)$.
If $e$ is any other event, $\calR(H') = \calR(H)$. In both cases
$\calR(H) \subseteq \calR(H')$, so $H \preceq_R H'$.
\end{proof}

\begin{corollary}[No Silent Unrefusal]
No legal Spherepop event removes a refusal from the history. Once $(x, r) \in
\calR(H)$, we have $(x, r) \in \calR(H')$ for all legal extensions $H'$ of $H$.
\end{corollary}

\section{Refusal Algebra}

\begin{definition}[Refusal Conjunction]
Histories $H$ and $H'$ with refusal sets $\calR(H)$ and $\calR(H')$ have
\emph{refusal conjunction}
$\calR(H) \wedge \calR(H') = \calR(H) \cup \calR(H')$---the set of symbols
refused in either history.
\end{definition}

\begin{definition}[Refusal Disjunction]
The \emph{refusal disjunction} is
$\calR(H) \vee \calR(H') = \calR(H) \cap \calR(H')$---the set of symbols
refused in both histories.
\end{definition}

\begin{proposition}[Refusal Sets Form a Lattice]
Under conjunction ($\cup$) and disjunction ($\cap$), the refusal sets of all
histories over $\Omega_0$ form a distributive lattice with bottom $\varnothing$
(no refusals) and top $\Omega_0 \times \mathsf{Reasons}$ (everything refused).
\end{proposition}

\begin{proof}
$(\wp(\Omega_0 \times \mathsf{Reasons}), \cup, \cap)$ is the powerset lattice,
which is always distributive.
\end{proof}

\begin{definition}[Refusal Implication]
$r_1 \Rightarrow_R r_2$ (refusal $r_1$ implies refusal $r_2$) if every history
containing $r_1$ necessarily contains $r_2$ due to dependency propagation:
$\Bind(x_1, x_2) \in H$ and $r_1 = (x_1, r)$ imply $r_2 = (x_2, \rho(r))$.
\end{definition}

\begin{theorem}[Refusal Logic is Monotone]
The refusal implication relation $\Rightarrow_R$ is monotone: if $H \preceq_R H'$
and $\calR(H) \vdash_R (x, r)$ (i.e.\ the refusal closure of $\calR(H)$
contains $(x,r)$), then $\calR(H') \vdash_R (x, r)$.
\end{theorem}

\begin{proof}
Since $\calR(H) \subseteq \calR(H')$ and the refusal closure is monotone in
the refusal set (adding more refusals only adds more propagated refusals),
the conclusion $(x, r)$ follows from $\calR(H')$ whenever it follows from
$\calR(H)$.
\end{proof}

\section{Repair as Certified Transition}

Repair is not the erasure of refusal. It is the construction of a new admissible
path that acknowledges and resolves the refusal.

\begin{definition}[Repair Operator]
A \emph{repair operator} is a partial map
\[
\mathsf{Repair}_\rho : \Rfsd(T, r) \to \Adm(T)
\]
defined only when $\rho$ is a proof that the condition producing refusal
reason $r$ has been resolved.
\end{definition}

\begin{theorem}[Repair Soundness]
If $\Gamma \vdash t : \Rfsd(T, r)$ and $\rho$ is a valid repair certificate for
$r$, then $\Gamma \vdash \mathsf{Repair}_\rho(t) : \Adm(T)$.
\end{theorem}

\begin{proof}
The typing rule for repair is:
\[
\frac{\Gamma \vdash t : \Rfsd(T, r) \quad \Gamma \vdash \rho : \mathsf{Resolved}(r)}
     {\Gamma \vdash \mathsf{Repair}_\rho(t) : \Adm(T)}
\quad [\mathrm{T\text{-}REPAIR}]
\]
The proof $\rho : \mathsf{Resolved}(r)$ certifies that the condition causing
refusal $r$ has been remedied. Given this certificate, the repaired term
satisfies the admissibility requirement for its type.
\end{proof}

\begin{theorem}[Repair Does Not Violate Refusal Monotonicity]
If a refused term is repaired, the original refusal remains in the history.
\end{theorem}

\begin{proof}
Repair appends a new event $\mathsf{Repair}(x, \rho)$ to the history.
It does not remove the earlier $\Refuse(x, r)$ event. Thus
$\Refuse(x, r) \in H$ implies $\Refuse(x, r) \in H \mathbin{++}
[\mathsf{Repair}(x, \rho)]$.
\end{proof}

\begin{remark}
This is the difference between repair and denial. Repair preserves the evidence
that something failed while adding the evidence that it was fixed. Denial removes
or suppresses the failure record. Spherepop permits repair but forbids denial as
a history operation---the history is immutable, and past refusals cannot be
erased.
\end{remark}

\chapter{Partial-Order Histories and True Concurrency}

\section{The Limitation of Total-Order Histories}

The histories developed throughout this monograph are sequences: total orders
on events. This is adequate for sequential computation but misrepresents true
concurrency. In a distributed system where events at different nodes occur
simultaneously, there is no canonical total order on them---only a partial
order of causal precedence.

\begin{definition}[Partial-Order History]
A \emph{partial-order history} (or \emph{event structure}) is a pair
$\mathcal{E} = (E, \prec)$ where $E$ is a finite set of events and $\prec$
is a strict partial order (irreflexive, asymmetric, transitive) representing
causal precedence: $e_1 \prec e_2$ means $e_1$ causally precedes $e_2$.
\end{definition}

\begin{definition}[Concurrent Events]
Events $e_1$ and $e_2$ are \emph{concurrent} (written $e_1 \parallel e_2$) if
neither $e_1 \prec e_2$ nor $e_2 \prec e_1$.
\end{definition}

\section{Linearizations and the Linearization Theorem}

\begin{definition}[Linearization]
A \emph{linearization} of $\mathcal{E} = (E, \prec)$ is a total order $H$ on
$E$ that extends $\prec$: if $e_1 \prec e_2$ then $e_1$ appears before $e_2$
in $H$.
\end{definition}

\begin{theorem}[Linearization Existence]
Every partial-order history $\mathcal{E} = (E, \prec)$ has at least one
linearization.
\end{theorem}

\begin{proof}
This is Szpilrajn's extension theorem: every partial order on a set can be
extended to a total order on that set.
\end{proof}

\begin{remark}[The number of linearizations is exponential in $|E|$, not simply $2^k$]
The earlier claim that $k$ pairs of concurrent events force at least $2^k$
linearizations is false as stated: the orientations of distinct concurrent
pairs are not in general independent choices, since ordering one pair can
constrain another through shared elements and transitivity (fixing
$e_1 \prec' e_2$ and $e_2 \prec' e_3$ already fixes the relative order of
$e_1$ and $e_3$ even if that pair was itself ``concurrent'' in $\mathcal{E}$).
The correct statement is only the existence guarantee above, together with the
standard fact that the number of linear extensions of a poset can be
exponential in the number of events: an antichain on $n$ elements (no order
constraints at all) has $n!$ linear extensions, and more refined poset
families used in linear-extension counting (e.g. the ``standard examples''
studied in the combinatorics literature) exhibit extension counts exponential
in $n$ even with a comparatively small number of covering relations. This is
still the formal foundation of the interleaving explosion noted in Part~XV---the
number of linearizations can grow explosively---but the precise scaling law
depends on the shape of the partial order, not just the count of concurrent
pairs.
\end{remark}

\section{Partial-Order Collapse Rules}

\begin{definition}[Partial-Order Collapse Rule]
A \emph{partial-order collapse rule} $c : \mathsf{POHist} \to O_c$ maps
partial-order histories to observable states. The rule is
\emph{linearization-invariant} if $c(\mathcal{E}) = c(H)$ for every
linearization $H$ of $\mathcal{E}$.
\end{definition}

\begin{theorem}[Collapse Invariance Theorem]
The Accumulate rule $c_{\mathrm{acc}}$ is linearization-invariant: the count
of Pop events for each symbol is the same in every linearization of any
partial-order history.
\end{theorem}

\begin{proof}
The Accumulate rule counts events by type, ignoring their order. Since every
linearization of $\mathcal{E}$ contains exactly the same events (just in
different orders), the count for each symbol is identical across all
linearizations.
\end{proof}

\begin{theorem}[LastWrite Is Not Linearization-Invariant]
There exist partial-order histories $\mathcal{E}$ and linearizations $H_1, H_2$
such that $c_{LW}(H_1) \neq c_{LW}(H_2)$.
\end{theorem}

\begin{proof}
Let $e_1 = \Pop(x{=}a)$ from process $P_1$ and $e_2 = \Pop(x{=}b)$ from process
$P_2$, two \emph{distinct, value-carrying} writes to $x$ with $a \neq b$ and
$e_1 \parallel e_2$ (concurrent). Writing both as bare $\Pop(x)$, as an earlier
draft did, would make $H_1 = [e_1, e_2]$ and $H_2 = [e_2, e_1]$ identical as
event sequences up to relabeling, since two occurrences of the same event carry
no information distinguishing which one ``won''; the argument requires
distinguishable values. With $a \neq b$: in linearization $H_1 = [e_1, e_2]$,
$c_{LW}$ observes the last write to $x$ as $b$. In linearization
$H_2 = [e_2, e_1]$, the last write is $a$. Since $a \neq b$,
$c_{LW}(H_1) \neq c_{LW}(H_2)$. This is the formal content of the race
condition: the LastWrite collapse is sensitive to linearization order.
\end{proof}

\begin{corollary}
Only linearization-invariant collapse rules can be used in distributed Spherepop
systems without requiring total ordering of events. Systems requiring LastWrite
semantics (such as most database systems) must impose a total order (via
timestamps, Lamport clocks, or consensus protocols).
\end{corollary}

\section{Lamport Clocks as Spherepop Events}

\begin{definition}[Lamport Clock Event]
A \emph{Lamport clock event} is a Pop event of the form $\Pop(\mathsf{ts}(e, t))$
where $e$ is a computational event and $t$ is a logical timestamp satisfying:
$t > \max\{t' \mid \Pop(\mathsf{ts}(e', t')) \in H \text{ and } e' \prec e\}$.
\end{definition}

\begin{theorem}[Lamport Clocks Produce Total-Order History]
Adding Lamport clock events to a partial-order history $\mathcal{E}$ produces
a total-order history $H_{LC}$ consistent with the causal ordering:
$e_1 \prec e_2 \Rightarrow \mathsf{ts}(e_1) < \mathsf{ts}(e_2)$ in $H_{LC}$.
\end{theorem}

\begin{proof}
By the Lamport clock construction, a Lamport timestamp is always greater than
the timestamps of causally preceding events. The resulting total order on
timestamps induces a total order on events that extends the causal partial order.
\end{proof}


%% ============================================================
%% PART XXIII: OPERATOR THEORY AND ALGEBRAIC STRUCTURE
%% ============================================================
\part{Operator Theory and Algebraic Structure}

\chapter{The Operator Representation Theorem}

\section{Primitive Transformations of Possibility Space}

The four Spherepop operators were derived in Part~II from engineering requirements.
This chapter provides the deeper derivation: the operators emerge as structural
necessities from the requirement that possibility space be managed in a way that
preserves historical accountability.

\begin{definition}[Possibility Space Transformation]
A \emph{possibility space transformation} is a function
$f : (\calH \times \wp(\Omega_0)) \to (\calH \times \wp(\Omega_0))$
that maps world states to world states.
\end{definition}

\begin{definition}[Primitive Transformation Types]
The four primitive transformation types are:
\begin{enumerate}
\item \emph{Selection} ($\mathsf{Sel}$): removes one element from $\Omega$ and records its removal.
\item \emph{Exclusion} ($\mathsf{Exc}$): records that an element is inadmissible without removing it from $\Omega$.
\item \emph{Coupling} ($\mathsf{Cpl}$): records a dependency between two elements without consuming either.
\item \emph{Projection} ($\mathsf{Proj}$): applies a function to the history, producing an observable.
\end{enumerate}
\end{definition}

\begin{theorem}[Operator Representation Theorem]
\label{thm:operator-representation}
Every legal transformation of a Spherepop world state $(H, \Omega)$ that
(a) appends to history monotonically, (b) does not increase $|\Omega|$,
(c) preserves the admissibility record, and (d) can produce a new observable
decomposes as a composition of the four primitive transformations:
\[
f = \mathsf{Proj}^{n_4} \circ \mathsf{Cpl}^{n_3} \circ \mathsf{Exc}^{n_2} \circ \mathsf{Sel}^{n_1}
\]
for some $n_1, n_2, n_3, n_4 \geq 0$, which correspond to sequences of
$\Pop$, $\Refuse$, $\Bind$, and $\Collapse$ operations.
\end{theorem}

\begin{proof}
We show each primitive transformation is realized by exactly one operator.

\textit{Selection ($\Pop$)}: The only transformation that decreases $|\Omega|$
is $P_x$ for some $x \in \Omega$. By requirement (b), no transformation increases
$|\Omega|$, and by requirement (a), every transformation appends to history.
A transformation that decreases $|\Omega|$ by exactly one and appends one event
is precisely $\Pop$.

\textit{Exclusion ($\Refuse$)}: A transformation that leaves $|\Omega|$ unchanged
but modifies $\calA(H)$ must append a $\Refuse$ event (the only event that
changes admissibility without changing $\Omega$). By requirement (c), the
reason must be recorded. This is precisely $\Refuse(x, r)$.

\textit{Coupling ($\Bind$)}: A transformation that leaves $|\Omega|$ and
$\calA(H)$ unchanged but appends a dependency record must be a $\Bind$ event.
No other event type records binary relations between symbols.

\textit{Projection ($\Collapse$)}: A transformation that produces an observable
output by applying a function to the history is precisely $\Collapse(x, c)$.
Requirements (a)--(c) do not force this transformation; requirement (d) does.

Any compound transformation decomposes into a sequence of these four types by
structural induction on the sequence of events appended to the history.
\end{proof}

\begin{remark}
The Operator Representation Theorem elevates the four operators from language
design choices to structural necessities. Any system that (a) maintains a
monotonically growing history, (b) manages a shrinking option space, (c)
preserves admissibility records, and (d) can produce observations must contain
operations isomorphic to Pop, Refuse, Bind, and Collapse.
\end{remark}

\section{Algebraic Properties of the Operators}

\begin{proposition}[Refuse Idempotence]
Applying $\Refuse(x, r)$ twice has the same effect as applying it once:
\[
\calA(H \mathbin{++} [\Refuse(x,r)]) = \calA(H \mathbin{++} [\Refuse(x,r)] \mathbin{++} [\Refuse(x,r)]).
\]
\end{proposition}

\begin{proof}
Both histories have $x$ removed from the admissibility set after the first
Refuse. The second Refuse has no additional effect since $x$ is already
inadmissible.
\end{proof}

\begin{proposition}[Pop-Refuse Non-Commutativity]
In general, $\Pop(x)$ followed by $\Refuse(x, r)$ is not the same as
$\Refuse(x, r)$ followed by $\Pop(x)$.
\end{proposition}

\begin{proof}
After $\Pop(x)$: $x \notin \Omega$ and $x \in \calA$ (Pop removes from $\Omega$).
Then $\Refuse(x, r)$: $x \notin \calA$ (inadmissible and not in $\Omega$).

After $\Refuse(x, r)$: $x \in \Omega$ but $x \notin \calA$ (Refuse removes
from $\calA$ only). Then $\Pop(x)$ would fail the admissibility check: $x \in \Omega$
but $x \notin \calA$, so the Pop is rejected.

The two sequences have different outcomes.
\end{proof}

\begin{theorem}[Collapse Non-Commutativity]
For collapse rules $c_1$ and $c_2$, in general
$c_1(c_2(H)) \neq c_2(c_1(H))$ when treating $c_1$ and $c_2$ as functions
on the history.
\end{theorem}

\begin{proof}
Let $c_1 = c_{LW}$ and $c_2 = c_{\mathrm{acc}}$. For history $H = [\Pop(x), \Pop(x)]$:
$c_{LW}(H) = \{x \mapsto 1\}$ (last write wins, $x$ is present once);
$c_{\mathrm{acc}}(H) = \{x \mapsto 2\}$ (two pops counted).
Applying the other rule to these outputs gives different results depending on
the order, since the outputs have different structures.
\end{proof}

\begin{theorem}[Bind Commutativity Under Symmetric Dependencies]
If $\Bind(a, b)$ and $\Bind(b, a)$ are both in $H$, then the dependency between
$a$ and $b$ is symmetric: $a \notin \calA(H^*) \Leftrightarrow b \notin \calA(H^*)$
under the refusal closure.
\end{theorem}

\begin{proof}
By the Dependency Propagation Theorem, $\Refuse(a, r) \in H$ propagates to
$\Refuse(b, \rho(r)) \in H^*$ via $\Bind(a, b)$, and $\Refuse(b, r') \in H$
propagates to $\Refuse(a, \rho(r')) \in H^*$ via $\Bind(b, a)$.
Hence inadmissibility propagates in both directions, creating symmetric mutual
inadmissibility.
\end{proof}

\chapter{Counterfactuals and Causal Reasoning}

\section{Continuations and Counterfactuals}

Causal reasoning asks: what would have happened if we had done otherwise? In
Spherepop, this is a question about alternative histories---histories that agree
with the actual up to some point and then diverge.

\begin{definition}[Counterfactual History]
A \emph{counterfactual history} relative to $H$ at step $k$ is any history
$H'$ such that $H'$ agrees with $H$ on the first $k-1$ events and differs on
event $k$: $H[1:k-1] = H'[1:k-1]$ but $H[k] \neq H'[k]$.
\end{definition}

\begin{definition}[Counterfactual Continuation Difference]
The \emph{counterfactual difference} between histories $H_1$ and $H_2$ is
\[
\calC(H_1) \triangle \calC(H_2) = (\calC(H_1) \setminus \calC(H_2)) \cup (\calC(H_2) \setminus \calC(H_1)),
\]
the symmetric difference of their future continuation sets.
\end{definition}

\begin{theorem}[Counterfactual Distinguishability Theorem]
If $\calC(H_1) \neq \calC(H_2)$, then $H_1$ and $H_2$ are computationally
non-equivalent: there exists a future computation that is possible from one
but not the other.
\end{theorem}

\begin{proof}
$\calC(H_1) \neq \calC(H_2)$ implies $\calC(H_1) \triangle \calC(H_2) \neq \varnothing$.
There exists $c$ in the symmetric difference: $c \in \calC(H_1) \setminus \calC(H_2)$
(or vice versa). This means $c$ is reachable from $H_1$ but not from $H_2$.
The two histories therefore differ in their computational potential.
\end{proof}

\begin{definition}[Causal Influence]
Event $e_i$ in history $H$ \emph{causally influences} continuation $c$ if
removing $e_i$ from $H$ (and adjusting subsequent events consistently) changes
whether $c$ is reachable:
\[
c \in \calC(H) \text{ but } c \notin \calC(H \setminus \{e_i\}).
\]
\end{definition}

\begin{theorem}[Pop Events Are Causally Primary]
Every continuation $c \in \calC(H)$ is causally influenced by at least one
Pop event in $H$.
\end{theorem}

\begin{proof}
A continuation $c$ consists of further events extending $H$. The specific
form of $c$ depends on what has been committed (the Pop events in $H$),
since commitments remove alternatives and shape which future events are
admissible. Removing all Pop events from $H$ leaves only Refuse, Bind, and
Collapse events, which preserve $\Omega$ and do not commit to specific paths.
The set $\calC(H_{\mathrm{no-pop}})$ is a superset of $\calC(H)$, so specific
continuations reachable from $H$ are reachable only because specific Pop events
have narrowed the possibility space in the right direction.
\end{proof}

\section{Causal Models and Spherepop}

\begin{definition}[Spherepop Causal Model]
A \emph{Spherepop causal model} is a dependency graph $G_H$ together with
a labeling of each node $x$ with its Pop or Refuse status in history $H$.
An intervention on $x$ is a counterfactual history $H'$ that changes the
event at position $k$ (where $x$ first appears) from $\Pop(x)$ to
$\Refuse(x, \mathsf{Intervened})$ or vice versa.
\end{definition}

\begin{theorem}[Do-Calculus Correspondence]
The Spherepop intervention operation corresponds to Pearl's do-operator:
$\mathsf{do}(\Pop(x))$ sets $x$ to committed status regardless of its upstream
causes, which in Spherepop corresponds to adding $\Pop(x)$ to the history while
severing the incoming dependency edges to $x$ in $G_H$.
\end{theorem}

\begin{proof}
Pearl's $\mathsf{do}(X = x)$ removes the incoming arrows to $X$ in the causal
graph and sets $X$ to $x$. In Spherepop: severing incoming dependency edges
corresponds to removing Bind events $(y, x)$ from $G_H$; setting $X$ to
committed corresponds to adding $\Pop(x)$. The resulting modified history
$H'$ is the Spherepop representation of the interventional distribution
$P(\cdot \mid \mathsf{do}(X = x))$.
\end{proof}

\chapter{Notes as a Category}

\section{Note Morphisms}

\begin{definition}[Note Morphism]
A \emph{note morphism} from note $N_1$ to note $N_2$ is a function
$f : \calC(N_1) \to \calC(N_2)$ that maps continuations preserved by $N_1$
to continuations preserved by $N_2$, such that
$P_A(f(c), t \mid N_2) \geq P_A(c, t \mid N_1)$: the target note preserves
the continuation at least as well as the source.
\end{definition}

\begin{theorem}[Notes Form a Category]
Notes with note morphisms form a category $\mathsf{Notes}$:
\begin{itemize}
\item Objects: notes $N$ (artifacts increasing reachability of some continuation).
\item Morphisms: note morphisms $f : N_1 \to N_2$.
\item Composition: $(g \circ f)(c) = g(f(c))$, which is a valid note morphism
since reachability is preserved by composition.
\item Identity: $\mathsf{id}_N(c) = c$ is the identity note morphism.
\end{itemize}
\end{theorem}

\begin{proof}
Composition: if $f : N_1 \to N_2$ and $g : N_2 \to N_3$ are note morphisms,
then for any $c \in \calC(N_1)$:
$P_A(g(f(c)), t \mid N_3) \geq P_A(f(c), t \mid N_2) \geq P_A(c, t \mid N_1)$.
So $g \circ f$ is a valid note morphism. Identities and associativity are immediate.
\end{proof}

\section{Note Functors}

\begin{theorem}[The Continuation Functor]
There is a functor $\calC : \mathsf{Notes} \to \mathsf{Set}$ sending each note
to its continuation set and each morphism to the corresponding function on
continuation sets.
\end{theorem}

\begin{proof}
$\calC$ sends note $N$ to $\calC(N) \subseteq \calC_{\mathrm{total}}$ and
morphism $f : N_1 \to N_2$ to $f : \calC(N_1) \to \calC(N_2)$. This preserves
composition and identities by definition of note morphism.
\end{proof}

\begin{theorem}[Note Composition Increases Reachability]
The tensor product $N_1 \otimes N_2$ of two notes satisfies
\[
\calC(N_1 \otimes N_2) \supseteq \calC(N_1) \cup \calC(N_2)
\]
with strict inclusion when the notes enable joint continuations not reachable
from either alone.
\end{theorem}

\begin{proof}
A combined note system has access to the continuations of each component plus
any joint continuations enabled by having both. Since $\calC(N_1) \cup \calC(N_2)
\subseteq \calC(N_1 \otimes N_2)$ by definition and joint continuations strictly
add to the union, the inclusion is proper when $\calC_{\mathrm{bind}}(N_1, N_2)
\neq \varnothing$.
\end{proof}

\section{The Reachability Value of Knowledge}

\begin{definition}[Knowledge Value Functional]
The \emph{knowledge value} of a note $N$ for agent $A$ at time $t$ is
\[
V_K(N, A, t) = \sum_{c \in \calC} (P_A(c, t \mid N) - P_A(c, t)) = G_A(N, t),
\]
the total reachability gain.
\end{definition}

\begin{theorem}[Knowledge Leverage Theorem]
The knowledge value of a note equals the total increase in reachability it
provides:
\[
V_K(N, A, t) = |\calC(N)_{\mathrm{new}}| \cdot \bar{g},
\]
where $|\calC(N)_{\mathrm{new}}|$ is the number of new continuations enabled
and $\bar{g}$ is the average gain per continuation.
\end{theorem}

\begin{proof}
Decompose the sum: $V_K(N, A, t) = \sum_{c \in \calC(N)} G_A(N, c, t) +
\sum_{c \notin \calC(N)} G_A(N, c, t)$. For $c \notin \calC(N)$, $G_A(N, c, t) = 0$
by definition of the continuation set. For $c \in \calC(N)$, the average gain
is $\bar{g}$. The result follows.
\end{proof}


%% ============================================================
%% PART XXIV: CIVILIZATION AS CONTINUATION RESERVOIR
%% ============================================================
\part{Civilization as Continuation Reservoir}

\chapter{Libraries, Archives, and Scientific Infrastructure}

\section{Libraries as Continuation Reservoirs}

\begin{definition}[Continuation Reservoir]
A collection of notes $\calN$ is a \emph{continuation reservoir} for agent
population $A$ if it strictly increases total reachability:
\[
\sum_{c \in \calC} P_A(c, t \mid \calN) > \sum_{c \in \calC} P_A(c, t).
\]
\end{definition}

\begin{theorem}[Library Reachability Theorem]
If each note $N_i \in \calN$ has non-negative reachability leverage and at
least one note has positive leverage, then $\calN$ is a continuation reservoir.
\end{theorem}

\begin{proof}
Let $\Delta_i(c) = P_A(c, t \mid N_i) - P_A(c, t)$. Non-negative leverage
means $\sum_c \Delta_i(c) \geq 0$ for all $i$. Positive leverage for some $N_j$
means $\sum_c \Delta_j(c) > 0$. Summing over all notes:
$\sum_{c,i} \Delta_i(c) > 0$, which implies the total reachability increases.
\end{proof}

\begin{remark}
Civilization is not primarily the accumulation of facts. It is the accumulation
of preserved re-entry points into abandoned, deferred, repaired, and generative
continuations. A library's value is not the mass of paper it contains but the
volume of continuations it makes accessible.
\end{remark}

\section{Science as Civilizational Refusal}

\begin{theorem}[Scientific Progress as Directed Refusal]
Scientific progress is a sequence of Refuse events applied to the space of
possible theories: each failed experiment refuses a class of theories with
reason $\mathsf{ExperimentalRefutation}$.
\end{theorem}

\begin{proof}
A scientific experiment tests a prediction $P$ derived from theory $T$. If
the experiment fails ($\neg P$ is observed), the falsification principle
(Popper) requires adding $\Refuse(T, \mathsf{Refuted}(P))$ to the scientific
community's history. This monotonically shrinks the admissibility set of
viable theories. Scientific progress is therefore the directed shrinkage of
the theory-admissibility set, guided by experiment.
\end{proof}

\begin{corollary}[Science Increases Admissibility Entropy Compression]
As science progresses, the admissibility entropy of the theory space
$S_{\calA}(\text{theories})$ decreases: fewer theories are admissible given
the accumulated experimental record.
\end{corollary}

\begin{remark}
This formalizes the intuition that science constrains rather than merely
accumulates. Each confirmed refutation reduces the set of viable theories,
making the remaining theories more precisely constrained. The admissibility
entropy of scientific theories in 2025 is dramatically lower than in 1600,
not because we have proved more but because we have refused more.
\end{remark}

\section{Writing Systems as Compression Infrastructure}

\begin{theorem}[Writing Reduces Reconstruction Cost Across Time]
A writing system reduces the reconstruction cost $C_r(c, o)$ for observations
$o$ made before its invention by making the associated histories accessible
across the temporal gap between creation and reading.
\end{theorem}

\begin{proof}
Without writing, an observation at time $t_1$ produces a fiber $c^{-1}(o)$
at time $t_2 > t_1$ of size exponential in $(t_2 - t_1)$ (all histories
consistent with the biological memory trace are candidates). With writing,
the fiber is reduced to histories consistent with the written record, which
is exponentially smaller. Hence $C_r(c, o)$ decreases by an amount proportional
to the entropy reduction.
\end{proof}

\begin{example}[The Rosetta Stone as Repair Note]
The Rosetta Stone (196 BCE) is a repair note for the ancient Egyptian writing
system. When Egyptian hieroglyphics became unreadable (the writing system was
refused by history---practitioners died and the tradition collapsed), the
Rosetta Stone provided a Bind event: it coupled the Greek text (whose
continuation space was still accessible) to the Egyptian text (whose continuation
space had become unreachable). The stone did not restore the original writing
tradition, but it created a repair morphism from the collapsed tradition to
a reconstructed one.
\end{example}

\section{Programming Languages as Generative Notes}

\begin{theorem}[Programming Language as Maximum-Leverage Note]
A general-purpose programming language $L$ is a generative note with unbounded
reachability leverage: $\Lambda(L) \to \infty$ as the number of programs
expressible in $L$ grows.
\end{theorem}

\begin{proof}
The continuation set $\calC(L)$ includes all computations expressible in $L$.
For a Turing-complete language, $|\calC(L)|$ is infinite (countably many
programs). The creation cost $C_{\mathrm{create}}(L)$ is finite (the language
was designed and implemented once). Therefore $\Lambda(L) = \sum_c \Delta P(c) /
C_{\mathrm{create}}(L)$ involves dividing a diverging numerator by a finite
denominator: $\Lambda(L) \to \infty$.
\end{proof}

\section{Legal Systems as Constraint Ledgers}

\begin{theorem}[Legal Code as Coordinated Refusal Note]
A legal code is a Coordination Note that documents, for a community of agents,
which actions are admitted (legal), refused (illegal), and under what conditions
repair (amnesty, appeal, rehabilitation) is available.
\end{theorem}

\begin{proof}
A law $L$ that prohibits action $a$ is a Refuse event applied to all agents
in the jurisdiction: $\Refuse(a, \mathsf{IllegalUnder}(L))$. A law that
requires action $b$ is a Bind event: $\Bind(\text{citizenship}, b)$, coupling
the agent's status to the obligation. Enforcement is the execution of the
refusal: making refusal-violations actually inadmissible (with sanctions). The
legal system's history is the sequence of legislative acts (Pop events on the
space of possible legal rules) and judicial decisions (Refuse events on specific
behavioral patterns).
\end{proof}

\begin{example}[Precedent as Collapse Note]
A legal precedent (stare decisis) is a Collapse Note: it collapses the family
of future cases similar to the precedent case into a single observable state
(the holding). The precedent makes the collapse rule explicit: future courts
must observe similar cases through the lens of the established holding, and
may only refine or distinguish the rule (not ignore it).
\end{example}

\section{Markets as Distributed Prospective Notes}

\begin{theorem}[Prices as Distributed Prospective Notes]
A market price $p$ for good $x$ is a Prospective Note: it preserves access
to the continuation of exchanging $x$ at price $p$ for all agents who observe
the price signal.
\end{theorem}

\begin{proof}
A price $p$ is a public Bind event: $\Bind(\text{seller of } x, \text{buyer of } x \text{ at } p)$.
It creates a prospective continuation---the exchange---that was not reachable
before the price was established. Without the price signal, buyers and sellers
must individually search for counterparties (high reconstruction cost). With
the price signal, the continuation of exchange is directly accessible
(low reconstruction cost). Market prices therefore reduce the reconstruction
cost of economic coordination.
\end{proof}

\chapter{History--Reachability Correspondence}

\section{Formal Statement}

Throughout this monograph, histories and reachability have been treated as
two aspects of the same underlying structure. This chapter makes the
correspondence explicit.

\begin{definition}[Reachability Volume of a History]
Let $\mathsf{Fut}(H)$ be the set of admissible continuations extending $H$
(finite sequences of events that are admissible given $H$). The
\emph{reachability volume} of $H$ is $V_R(H) = |\mathsf{Fut}(H)|$.
\end{definition}

\begin{theorem}[Pop Narrows Reachability]
If $H' = H \mathbin{++} [\Pop(x)]$, then $V_R(H') \leq V_R(H)$.
\end{theorem}

\begin{proof}
Every continuation extending $H'$ is also a continuation extending $H$, since
$H'$ itself extends $H$. Thus $\mathsf{Fut}(H') \subseteq \mathsf{Fut}(H)$
and $V_R(H') \leq V_R(H)$.
\end{proof}

\begin{theorem}[Refuse Narrows Admissible Reachability]
If $H' = H \mathbin{++} [\Refuse(x,r)]$, then $V_R(H') \leq V_R(H)$.
\end{theorem}

\begin{proof}
A continuation extending $H'$ must respect the refusal of $x$: any continuation
requiring an unrepaired Pop of $x$ is excluded from $\mathsf{Fut}(H')$.
Since all continuations extending $H'$ are continuations extending $H$ with
an additional admissibility constraint, $\mathsf{Fut}(H') \subseteq \mathsf{Fut}(H)$.
\end{proof}

\begin{corollary}[Computation as Reachability Shaping]
A Spherepop computation is not merely a path from input to output. It is a
monotone transformation of a reachability volume: commitment (Pop) and refusal
(Refuse) both narrow the reachability volume, while Bind and Collapse reshape
its structure without necessarily shrinking it.
\end{corollary}

\section{The History-Reachability Duality}

\begin{theorem}[History-Reachability Duality]
There is a duality between the lattice of histories ordered by prefix extension
and the lattice of reachability volumes ordered by subset inclusion:
\[
H_1 \leq H_2 \text{ (prefix)} \implies V_R(H_2) \leq V_R(H_1) \text{ (subset)}.
\]
\end{theorem}

\begin{proof}
If $H_1$ is a prefix of $H_2$, then $H_2 = H_1 \mathbin{++} H'$ for some
non-empty $H'$. By repeated application of Pop Narrows Reachability and
Refuse Narrows Admissible Reachability (one application per event in $H'$),
$V_R(H_2) \leq V_R(H_1)$.
\end{proof}

\begin{remark}
The History-Reachability Duality is the formal content of the monograph's
central thesis: computation is the progressive restriction of possibility.
As the history grows (computation proceeds), the reachability volume shrinks.
The history and the reachability volume are dual representations of the
same computational trajectory.
\end{remark}

\section{The History-Reachability Correspondence Theorem}

\begin{theorem}[History-Reachability Correspondence]
\label{thm:hrc}
Let $\Phi : \calH \to \wp(\calC)$ be the map sending history $H$ to its
reachability volume $\mathsf{Fut}(H)$. Then:
\begin{enumerate}
\item $\Phi$ is an order-reversing map between the prefix lattice on $\calH$
and the inclusion lattice on $\wp(\calC)$.
\item $\Phi$ preserves the Conservation Law: $|\Omega(H)| + |\mathsf{consumed}(H)| = |\Omega_0|$.
\item $\Phi(\varepsilon) = \calC_{\mathrm{total}}$ (the empty history has maximal reachability).
\item $\Phi(H_{\mathrm{terminal}}) = \varnothing$ when $\Omega(H_{\mathrm{terminal}}) = \varnothing$.
\end{enumerate}
\end{theorem}

\begin{proof}
(1) is the History-Reachability Duality theorem. (2) follows from the Conservation
of Possibility: committed events ($|\mathsf{consumed}(H)|$) plus remaining
options ($|\Omega(H)|$) equal the initial option count. (3) holds because the
empty history has imposed no commitments, so all continuations are reachable.
(4) holds because if $\Omega = \varnothing$, no further Pop events are possible,
so no continuation can extend $H$ by selection.
\end{proof}

\begin{remark}
Theorem~\ref{thm:hrc} is the unification theorem promised from the beginning of
this monograph. It shows that Spherepop's two primitive concepts---the history
that computation produces and the reachability volume that computation consumes---are
formally dual. Each determines the other. A complete history determines a
terminal reachability volume. A specified reachability volume constrains the
possible histories that could have produced it.

This duality is the same duality that appears throughout the monograph in
other guises: state versus history, observation versus trajectory, collapse
versus reconstruction, notes as anti-collapse artifacts.

All of these are the same duality, seen from different angles.
\end{remark}

%% ============================================================
%% THEOREM STATUS TABLE
%% ============================================================
\chapter*{Theorem Status Table}
\addcontentsline{toc}{chapter}{Theorem Status Table}

The following table records the status of principal results in this monograph,
implementing the reviewer's recommendation for epistemic transparency.

\begin{center}
\renewcommand{\arraystretch}{1.15}
\begin{longtable}{lll}
\toprule
Result & Chapter & Status \\
\midrule
Conservation of Possibility & 3 & Proved \\
Representation Adequacy of Four Operators (formerly ``Completeness'') & 3 & Sketch \\
Boolean Completeness of Refuse and Bind & 3 & Proved \\
Operator Representation Theorem & 45 & Proved \\
History Monoid (free monoid) & 4 & Proved \\
Hist as free strict monoidal category & 4 & Proved \\
Unique History Factorisation & 4 & Proved \\
History Monotonicity & 5 & Proved \\
Option Space Monotonicity & 5 & Proved \\
Collapse Soundness & 7 & Proved \\
Type Preservation (full) & 40 & Proved \\
Progress (full) & 40 & Proved \\
Type Soundness Corollary & 40 & Proved \\
Canonical Forms Lemma & 40 & Proved \\
Weakening Lemma & 40 & Proved \\
Exchange Lemma & 40 & Proved \\
Substitution Lemma & 40 & Proved \\
Admissibility Preservation & 40 & Proved \\
Pop Entropy Reduction & 42 & Proved \\
Refuse Admissibility Entropy Reduction & 42 & Proved \\
State Illusion (Proposition 2.2) & 2 & Proved \\
Finite Observation Non-injectivity (formerly ``No Direct Observation'') & 17 & Proved \\
Ephemerality Inversion & 13 & Proved \\
Collapse Functor Functoriality & 6 & Proved (Conditional)$^\dagger$ \\
Collapse as Reflection (adjunction) & 20 & Proved (Conditional)$^\|$ \\
Universal Property of Collapse & 6 & Proved \\
History-Observation Adjunction & 20 & Proved (Conditional)$^\|$ \\
Notes as Partial Right Adjoints & 25 & Sketch \\
State Illusion as Sheaf Non-gluing & 20 & Proved \\
Refusal Obstruction Prevents Gluing & 47 & Proved \\
Civilizational Collapse Theorem & 21 & Proved \\
Percolation Phase Transition & 38 & Sketch \\
Replay Equivalence & 10 & Proved \\
Compiler Full Faithfulness & 22 & Proved (Extended)$^\S$ \\
GC Safety (Certificate-Only) & 15 & Proved \\
Error Correction via Certified Refusal & 16 & Proved \\
Eve's Law as Pythagorean Theorem & 31 & Proved \\
Tower Property as Nested Projection & 31 & Proved \\
Dependency Propagation Theorem & 47 & Proved \\
Dependency Closure Theorem & 47 & Proved \\
Deadlock as Graph Cycle & 34 & Proved \\
Linearization Existence (formerly ``Linearization Theorem'') & 48 & Proved$^\P$ \\
Collapse Invariance (Accumulate) & 48 & Proved \\
History-Reachability Duality & 51 & Proved \\
History-Reachability Correspondence & 51 & Proved \\
Library Reachability Theorem & 50 & Proved \\
Scientific Progress as Directed Refusal & 50 & Proved \\
Repair Soundness & 47 & Proved \\
Notes Form a Category & 46 & Proved \\
Knowledge Leverage Theorem & 46 & Proved \\
Entropy Lower Bound on Reconstruction & 23 & Proved \\
Optimal Note Theorem & 48 & Proved \\
Bisimulation Is Congruence & 34 & Proved \\
CAP as Collapse Tradeoff & 34 & Sketch \\
CRDTs Achieve Eventual Consistency & 34 & Proved \\
Repair Cohomology (obstruction) & 26 & Conjecture \\
Dependent Process Types (full) & 8 & Programmatic \\
Distributed Spherepop Runtimes & 34 & Programmatic \\
CoC Strong Normalization & 24 & Sketch \\
Proof-Carrying Refusal (full system) & 8 & Programmatic \\
Admissibility Lattice Completeness & 8 & Programmatic \\
\bottomrule
\end{longtable}
\end{center}

\noindent
\textbf{Status codes:}
\emph{Proved}: formal proof given in this manuscript.
\emph{Sketch}: main argument given; details omitted or rely on standard references.
\emph{Conjecture}: proposed result, no proof given.
\emph{Programmatic}: research direction; result not yet formulated precisely enough to prove.

\smallskip
\noindent
$\dagger$ Proved only for collapse rules $c$ that respect composition
(monoid homomorphisms), $c(H_2 \mathbin{++} H_1) = c(H_2) \circ c(H_1)$; not a
general statement about arbitrary observation functions. See the remark
following the functoriality proof.\\
$\ddagger$ Written before the VIEW/Collapse separation of
Section~\ref{sec:view-vs-collapse}; repaired below, see $\|$.\\
$\S$ The original theorem covers history equality only. A corollary is now
proved alongside it (Corollary~\ref{cor:op-state-faithfulness}) extending the
result to operational-state equality, discharging the replay test introduced
in Chapter~\ref{ch:compiler}
(\texttt{world.materialized == fold\_operational(\&world.history)}); history
equality remains the strictly stronger of the two criteria.\\
$\P$ Existence (Szpilrajn) is proved; the earlier $2^k$ lower-bound claim on the
number of linearizations was false as stated and has been removed. See the
remark following the existence theorem.\\
$\|$ Proved for \emph{prefix-monotone} collapse rules equipped with an
order-reflecting canonical-representative section (Theorem~\ref{thm:reflection-repaired}),
not for arbitrary $c$: the Identity and Accumulate rules satisfy the hypotheses
(Accumulate worked out explicitly as the nontrivial case), LastWrite does not,
for the same reason it fails Collapse Functor Functoriality. The original
statements conflated two incompatible readings of ``$\Hist$'' (histories as
morphisms vs.\ histories as objects); the repair introduces the explicit prefix
category $\Hist_\sqsubseteq$ to resolve this. See
Section~\ref{sec:view-vs-collapse}'s companion chapter.

\begin{thebibliography}{99}

\bibitem{barendregt84}
Barendregt, H.~P. \emph{The Lambda Calculus: Its Syntax and Semantics}.
North-Holland, 1984.

\bibitem{borgman07}
Borgman, C.~L. \emph{Scholarship in the Digital Age}. MIT Press, 2007.

\bibitem{clark98}
Clark, A. and Chalmers, D. The extended mind. \emph{Analysis} 58, no.~1 (1998):
7--19.

\bibitem{eisenstein80}
Eisenstein, E.~L. \emph{The Printing Press as an Agent of Change}. Cambridge
University Press, 1980.

\bibitem{feldman58}
Feldman, J. On the absolute continuity of Gaussian measures. \emph{Zeitschrift
f\"ur Wahrscheinlichkeitstheorie} 1 (1958).

\bibitem{felleisen87}
Felleisen, M. and Friedman, D.~P. A calculus for assignments in higher-order
languages. \emph{Proceedings of the 14th ACM SIGACT-SIGPLAN Symposium on
Principles of Programming Languages}, 1987.


\bibitem{chenoweth25}
Chenoweth, C., Chenoweth, A.
MEM\textbar8: A Wave-Based Cognitive Architecture for Multimodal Memory Integration
and Consciousness Simulation. Zenodo, 2025.
\texttt{https://doi.org/10.5281/zenodo.16436298}

\bibitem{coquand88}
Coquand, T. and Huet, G. The calculus of constructions.
\emph{Information and Computation} 76 (1988): 95--120.

\bibitem{fowler05}
Fowler, M. Event sourcing. \texttt{martinfowler.com}, 2005.

\bibitem{girard89}
Girard, J.-Y., Lafont, Y., and Taylor, P. \emph{Proofs and Types}. Cambridge
University Press, 1989.

\bibitem{goody77}
Goody, J. \emph{The Domestication of the Savage Mind}. Cambridge University
Press, 1977.

\bibitem{hajek58}
H\'ajek, J. On a property of normal distributions of an arbitrary stochastic
process. \emph{Czechoslovak Mathematical Journal} 8 (1958).

\bibitem{hoare85}
Hoare, C.~A.~R. \emph{Communicating Sequential Processes}. Prentice Hall, 1985.

\bibitem{hofmann98}
Hofmann, M. and Streicher, T. The groupoid interpretation of type theory.
In \emph{Twenty-five Years of Constructive Type Theory}. Oxford University Press,
1998.

\bibitem{hutchins95}
Hutchins, E. \emph{Cognition in the Wild}. MIT Press, 1995.

\bibitem{knuth84}
Knuth, D.~E. Literate programming. \emph{The Computer Journal} 27, no.~2
(1984): 97--111.

\bibitem{maclane71}
Mac Lane, S. \emph{Categories for the Working Mathematician}. Springer, 1971.

\bibitem{martinlof84}
Martin-L\"of, P. \emph{Intuitionistic Type Theory}. Bibliopolis, 1984.

\bibitem{milner80}
Milner, R. \emph{A Calculus of Communicating Systems}. Springer, 1980.

\bibitem{naur60}
Naur, P. (ed.). Report on the algorithmic language ALGOL 60.
\emph{Communications of the ACM} 3, no.~5 (1960): 299--314.
(Original introduction of Backus--Naur Form.)

\bibitem{nordstrom90}
Nordstr\"om, B., Petersson, K., and Smith, J. \emph{Programming in
Martin-L\"of's Type Theory}. Oxford University Press, 1990.

\bibitem{ong82}
Ong, W.~J. \emph{Orality and Literacy: The Technologizing of the Word}.
Methuen, 1982.

\bibitem{pierce02}
Pierce, B.~C. \emph{Types and Programming Languages}. MIT Press, 2002.

\bibitem{reiter78}
Reiter, R. On closed world data bases. In \emph{Logic and Data Bases}. Plenum
Press, 1978.

\bibitem{santoro24}
Santoro, M., Waghmare, S., and Panaretos, V.~M. Kernel embeddings of functional
data and mutual singularity. Preprint, 2024.

\bibitem{shapiro11}
Shapiro, M., Pregui\c{c}a, N., Baquero, C., and Zawirski, M. Conflict-free
replicated data types. INRIA Technical Report, 2011.

\bibitem{strachey67}
Strachey, C. Fundamental concepts in programming languages. \emph{Higher-Order
and Symbolic Computation} 13 (2000): 11--49. Originally delivered 1967.

\bibitem{tulving72}
Tulving, E. Episodic and semantic memory. In \emph{Organization of Memory}.
Academic Press, 1972.

\bibitem{univalent13}
The Univalent Foundations Program. \emph{Homotopy Type Theory: Univalent
Foundations of Mathematics}. Institute for Advanced Study, 2013.

\bibitem{wright94}
Wright, A.~K. and Felleisen, M. A syntactic approach to type soundness.
\emph{Information and Computation} 115, no.~1 (1994): 38--94.

\end{thebibliography}


\end{document}
